\documentclass[11pt]{article}
\usepackage[T1]{fontenc}
\usepackage[utf8]{inputenc}
\usepackage{lmodern}
\usepackage[margin=1in]{geometry}
\usepackage{amsmath,amssymb,amsthm,mathtools,bm}
\usepackage{longtable,booktabs,array,multirow}
\usepackage{textcomp}
\usepackage{graphicx}
\usepackage{enumitem}
\usepackage{xcolor}
\usepackage{framed}
\usepackage{fancyvrb}
\usepackage{ragged2e}
\usepackage{hyperref}
\usepackage{xurl}
\input{release_info.tex}

% Springer/pdflatex compatibility: provide fallbacks for Unicode math/text glyphs.
\newcounter{none}
\providecommand{\theHnone}{\arabic{none}}
\providecommand{\textsubscript}[1]{$_{\mathrm{#1}}$}
\newcommand{\ellbar}{\bar{\ell}}
\newcommand{\ophid}[1]{\path{#1}}
\ifdefined\DeclareUnicodeCharacter
\DeclareUnicodeCharacter{00B1}{\ensuremath{\pm}}
\DeclareUnicodeCharacter{00B2}{\textsuperscript{2}}
\DeclareUnicodeCharacter{00B3}{\textsuperscript{3}}
\DeclareUnicodeCharacter{00B7}{\ensuremath{\cdot}}
\DeclareUnicodeCharacter{00B9}{\textsuperscript{1}}
\DeclareUnicodeCharacter{00D7}{\ensuremath{\times}}
\DeclareUnicodeCharacter{02B8}{\textsuperscript{y}}
\DeclareUnicodeCharacter{0393}{\ensuremath{\Gamma}}
\DeclareUnicodeCharacter{0394}{\ensuremath{\Delta}}
\DeclareUnicodeCharacter{039B}{\ensuremath{\Lambda}}
\DeclareUnicodeCharacter{03A3}{\ensuremath{\Sigma}}
\DeclareUnicodeCharacter{03A6}{\ensuremath{\Phi}}
\DeclareUnicodeCharacter{03A9}{\ensuremath{\Omega}}
\DeclareUnicodeCharacter{03B1}{\ensuremath{\alpha}}
\DeclareUnicodeCharacter{03B2}{\ensuremath{\beta}}
\DeclareUnicodeCharacter{03B3}{\ensuremath{\gamma}}
\DeclareUnicodeCharacter{03B4}{\ensuremath{\delta}}
\DeclareUnicodeCharacter{03B5}{\ensuremath{\varepsilon}}
\DeclareUnicodeCharacter{03B7}{\ensuremath{\eta}}
\DeclareUnicodeCharacter{03B8}{\ensuremath{\theta}}
\DeclareUnicodeCharacter{03BA}{\ensuremath{\kappa}}
\DeclareUnicodeCharacter{03BB}{\ensuremath{\lambda}}
\DeclareUnicodeCharacter{03BC}{\ensuremath{\mu}}
\DeclareUnicodeCharacter{03BD}{\ensuremath{\nu}}
\DeclareUnicodeCharacter{03BE}{\ensuremath{\xi}}
\DeclareUnicodeCharacter{03C0}{\ensuremath{\pi}}
\DeclareUnicodeCharacter{03C1}{\ensuremath{\rho}}
\DeclareUnicodeCharacter{03C3}{\ensuremath{\sigma}}
\DeclareUnicodeCharacter{03C4}{\ensuremath{\tau}}
\DeclareUnicodeCharacter{03C6}{\ensuremath{\phi}}
\DeclareUnicodeCharacter{03C7}{\ensuremath{\chi}}
\DeclareUnicodeCharacter{03C8}{\ensuremath{\psi}}
\DeclareUnicodeCharacter{03C9}{\ensuremath{\omega}}
\DeclareUnicodeCharacter{1D9C}{\textsuperscript{c}}
\DeclareUnicodeCharacter{2020}{\ensuremath{\dagger}}
\DeclareUnicodeCharacter{2070}{\textsuperscript{0}}
\DeclareUnicodeCharacter{2074}{\textsuperscript{4}}
\DeclareUnicodeCharacter{2075}{\textsuperscript{5}}
\DeclareUnicodeCharacter{2076}{\textsuperscript{6}}
\DeclareUnicodeCharacter{2077}{\textsuperscript{7}}
\DeclareUnicodeCharacter{2078}{\textsuperscript{8}}
\DeclareUnicodeCharacter{2079}{\textsuperscript{9}}
\DeclareUnicodeCharacter{207A}{\textsuperscript{+}}
\DeclareUnicodeCharacter{207B}{\textsuperscript{-}}
\DeclareUnicodeCharacter{207D}{\textsuperscript{(}}
\DeclareUnicodeCharacter{207E}{\textsuperscript{)}}
\DeclareUnicodeCharacter{207F}{\textsuperscript{n}}
\DeclareUnicodeCharacter{2080}{\textsubscript{0}}
\DeclareUnicodeCharacter{2081}{\textsubscript{1}}
\DeclareUnicodeCharacter{2082}{\textsubscript{2}}
\DeclareUnicodeCharacter{2083}{\textsubscript{3}}
\DeclareUnicodeCharacter{2085}{\textsubscript{5}}
\DeclareUnicodeCharacter{2086}{\textsubscript{6}}
\DeclareUnicodeCharacter{2087}{\textsubscript{7}}
\DeclareUnicodeCharacter{2088}{\textsubscript{8}}
\DeclareUnicodeCharacter{208B}{\textsubscript{-}}
\DeclareUnicodeCharacter{2102}{\ensuremath{\mathbb{C}}}
\DeclareUnicodeCharacter{210B}{\ensuremath{\mathcal{H}}}
\DeclareUnicodeCharacter{2113}{\ensuremath{\ell}}
\DeclareUnicodeCharacter{2124}{\ensuremath{\mathbb{Z}}}
\DeclareUnicodeCharacter{2192}{\ensuremath{\to}}
\DeclareUnicodeCharacter{2193}{\ensuremath{\downarrow}}
\DeclareUnicodeCharacter{2194}{\ensuremath{\leftrightarrow}}
\DeclareUnicodeCharacter{21D2}{\ensuremath{\Rightarrow}}
\DeclareUnicodeCharacter{2202}{\ensuremath{\partial}}
\DeclareUnicodeCharacter{2208}{\ensuremath{\in}}
\DeclareUnicodeCharacter{220F}{\ensuremath{\prod}}
\DeclareUnicodeCharacter{2212}{\ensuremath{-}}
\DeclareUnicodeCharacter{221A}{\ensuremath{\surd}}
\DeclareUnicodeCharacter{221D}{\ensuremath{\propto}}
\DeclareUnicodeCharacter{221E}{\ensuremath{\infty}}
\DeclareUnicodeCharacter{2229}{\ensuremath{\cap}}
\DeclareUnicodeCharacter{222B}{\ensuremath{\int}}
\DeclareUnicodeCharacter{223C}{\ensuremath{\sim}}
\DeclareUnicodeCharacter{2245}{\ensuremath{\cong}}
\DeclareUnicodeCharacter{2248}{\ensuremath{\approx}}
\DeclareUnicodeCharacter{2261}{\ensuremath{\equiv}}
\DeclareUnicodeCharacter{2264}{\ensuremath{\leq}}
\DeclareUnicodeCharacter{2265}{\ensuremath{\geq}}
\DeclareUnicodeCharacter{226A}{\ensuremath{\ll}}
\DeclareUnicodeCharacter{2272}{\ensuremath{\lesssim}}
\DeclareUnicodeCharacter{2295}{\ensuremath{\oplus}}
\DeclareUnicodeCharacter{2297}{\ensuremath{\otimes}}
\DeclareUnicodeCharacter{2609}{\ensuremath{\odot}}
\DeclareUnicodeCharacter{2713}{\ensuremath{\checkmark}}
\DeclareUnicodeCharacter{27E8}{\ensuremath{\langle}}
\DeclareUnicodeCharacter{27E9}{\ensuremath{\rangle}}
\DeclareUnicodeCharacter{27F9}{\ensuremath{\Longrightarrow}}
\fi

\hypersetup{colorlinks=true,linkcolor=blue,citecolor=blue,urlcolor=blue}
\setcounter{tocdepth}{2}
\setcounter{secnumdepth}{3}
\setlength{\LTleft}{0pt}
\setlength{\LTright}{0pt}
\setlength{\tabcolsep}{1.2pt}
\setlength{\emergencystretch}{3em}
\renewcommand{\arraystretch}{1.05}

\providecommand{\tightlist}{%
 \setlength{\itemsep}{0pt}\setlength{\parskip}{0pt}}
\definecolor{shadecolor}{RGB}{248,248,248}
\newenvironment{Shaded}{\begin{snugshade}}{\end{snugshade}}
\DefineVerbatimEnvironment{Highlighting}{Verbatim}{commandchars=\\\{\}}
\newcommand{\NormalTok}[1]{#1}
\newcommand{\BuiltInTok}[1]{#1}
\newcommand{\CommentTok}[1]{#1}
\newcommand{\ExtensionTok}[1]{#1}
\newcommand{\AttributeTok}[1]{#1}
\newcommand{\SU}{\mathrm{SU}}
\newcommand{\U}{\mathrm{U}}

\title{Observers Are All You Need}
\author{B. M\"uller \and Kai Xue \and Peter Nguyen}
\date{\OPHPaperReleaseDate}

\begin{document}
\maketitle
\OPHPaperReleaseBanner

\begin{abstract}
Observer-Patch Holography (OPH) starts from one idea: no observer sees the whole world at once, so physics must be rebuilt from local descriptions that agree on overlaps. This synthesis paper gathers the OPH program into one view. It gives the fixed-cutoff theorem package for patch gluing, higher-gauge defects, consensus, records, and observer microphysics; a conditional route to Lorentzian spacetime and general relativity; and a conditional reconstruction of the realized Standard Model quotient \(\SU(3)\times\SU(2)\times\U(1)/\mathbb Z_6\) with the exact hypercharge lattice and the counting chain \(N_g=3\), \(N_c=3\). On the particle side it fixes the massless photon, gluons, and graviton, closes the electroweak \(W/Z\) lane, carries quantitative Higgs/top and neutrino branches, and states the charged, quark, and hadron boundaries explicitly. The same local input \(P\) also organizes the local unification surface for \(W\), \(Z\), \(H\), and, on the declared exact-release extension surface, the gravity coupling \(G\), while the invariant causal speed \(c\) is structural from the Lorentz branch.
\end{abstract}

\tableofcontents
\newpage

\section{Introduction}

This synthesis paper reads the OPH suite on one unified synthesis surface. It brings
together the recovered core from Ref.~\cite{ophcompact}, the finite patch-net consensus spine from
Ref.~\cite{ophconsensus}, the screen-microphysics engineering lane from Ref.~\cite{ophmicrophysics},
and the downstream particle-spectrum lane from Ref.~\cite{ophparticles}. The guiding
question is whether overlap consistency for local observer descriptions on a finite-capacity screen
is strong enough to recover the structural features of low-energy physics, and if so, how much of
that recovery is theorem-level and how much depends on explicit scaling-limit,
categorical, or continuation premises.

The answer is sharply tiered. Fixed-cutoff collar structure, higher-gauge defect transport,
the consensus/fixed-point formulation together with observable-level quotient confluence on the
declared physical algebras, the
record-algebra / Born-L\"uders package with its quantitative stability surface, and the
checkpoint/restoration package are theorem-bearing on their stated finite
surfaces. The Lorentz/null-modular/Einstein chain and compact gauge reconstruction remain
conditional continuum or categorical branches. The particle lane is the most tiered part of
the suite: the structural carrier skeleton and realized Standard Model branch are fixed, the
electroweak stage remains a forward-emitted non-core Phase-II calibration branch whose internal
transmutation data are fixed from \(P\), with target-free public \(W/Z\) rows and a compare-only
frozen validation pair, the Higgs/top stage is a secondary quantitative branch, the neutrino lane
closes on its weighted-cycle theorem branch with diagnostic sidecars beneath it, the charged-lepton lane is fixed
only up to exact centered readback plus a closed common-shift no-go, and the quark,
and hadron lanes remain explicitly unfinished for different reasons.

The paper is organized by lane. Section~2 gives the synthesis-level map. Section~3 carries the
recovered-core trunk: axioms, premise ledger, collar tools, gluing, the Lorentz/gravity branch,
and the compact-gauge/Standard-Model branch. Section~4 records the consensus lane. Section~5
summarizes the particle lane. Section~6 records the screen-microphysics and observer lane.
Section~7 records the conditional worldsheet lane. Section~8 gives the cross-lane theorem
boundary and open directions. The appendices collect extra architecture and interpretive
material that is useful in the synthesis paper but not part of the recovered core itself.

\section{Overview of Results}

The OPH suite supports the following lane-level split.

\begin{enumerate}[leftmargin=2em]
\item \textbf{At fixed cutoff, local structure and repair are sharp.} Technically, the
collar theorem package is explicit and finite-range, the genuinely noncentral higher-gauge branch
is closed at fixed cutoff, and the consensus lane gives a unique schedule-independent quotient
normal form together with schedule-independent physical observable values on the declared
fixed-cutoff physical algebras once the declared recovery-derived repair moves are committed by the
touched-overlap local-fit contract that yields strict decrease of the inconsistency potential
\(\Phi\), disjoint accepted repairs commute by support locality, the fixed-cutoff center-sector /
higher-gauge gluing package makes overlapping repairs parenthesization-independent on the physical
quotient, and repair completeness holds on the stated finite patch-net branch. Microscopic
representative lifts may still differ by gauge or sector relabelings inside one declared
quotient-local glued state without changing the resulting observable values on those algebras.
\item \textbf{Relativity appears as a controlled continuum branch.}
Technically, the Lorentz/null-modular/Einstein chain is a refinement/scaling-limit branch with
explicit canonical technical premises.
\item \textbf{The Standard Model gauge skeleton appears as a constrained consistency output.}
Technically, compact gauge reconstruction and the realized
\(\SU(3)\times\SU(2)\times\U(1)/\mathbb Z_6\) branch follow under the stated transportability and
categorical premises, with the realized counting chain \(N_g=3\) then \(N_c=3\).
\item \textbf{The particle lane has real outputs and a nonuniform closure profile.}
Technically, the structural carriers are exact, the electroweak stage remains on a non-core
Phase-II calibration branch with target-free public \(W/Z\) rows and a compare-only frozen
validation pair, its internal transmutation data are fixed forward from \(P\) rather than read
back from measured couplings, the Higgs/top stage is quantitatively emitted on a secondary branch,
the charged-lepton lane carries exact centered readback plus a closed common-shift no-go, the
neutrino lane closes on its weighted-cycle theorem branch with diagnostic sidecars below it, and the
quark and hadron lanes carry named remaining closure objects;
at the local numerical level the same declared input \(P\) already organizes the bosonic mass
trunk and the gravity-side bridge candidate.
\item \textbf{The screen-microphysics lane is explicit and concrete.} Technically, it
contains an explicit regulated reference architecture, sharpened patch-net and edge-sector
branches, a closed fixed-cutoff record-algebra / Born-L\"uders package with quantitative stability
bounds, and a closed checkpoint/restoration package.
\item \textbf{The UV picture is sharper at fixed cutoff than in the continuum.} Technically, OPH
fixes the physical UV branch only modulo implementation hiding and inert ancillary refinement; the
UV burden is the continuum BW/geometric lift and the refinement-limit transportable-sector
lift.
\item \textbf{The string/worldsheet lane remains conditional.} Technically, the fixed-cutoff
heat-kernel edge-sector result stays on the declared compact-gauge branch, any worldsheet-like
reorganization requires a controlled large-\(N_{\mathrm{edge}}\) regime distinct from that proved
chain, and any critical-superstring lift remains conjectural rather than part of the recovered core.
\end{enumerate}

\subsection*{Companion Papers}

The detailed depth surfaces used throughout this synthesis are:
\begin{enumerate}[leftmargin=2em]
\item \emph{Recovering Relativity and the Standard Model from the OPH Package Rooted in Observer Consistency}~\cite{ophcompact}, which carries the premise ledger, the recovered-core Lorentz/gravity route, and the compact-gauge/Standard-Model branch.
\item \emph{Reality as a Consensus Protocol: The Fixed-Point Computation That Implements Physics}~\cite{ophconsensus}, which carries the finite patch-net fixed-point, defect, quotient, and operator-record consensus package.
\item \emph{Screen Microphysics, Patches, and Observer Synchronization in OPH}~\cite{ophmicrophysics}, which carries the regulated reference architecture, the fixed-cutoff edge heat-kernel / Casimir theorem, and the measurement and observer checkpoint/restoration packages.
\item \emph{Deriving the Particle Zoo from Observer Consistency}~\cite{ophparticles}, which carries the structural-to-family particle route and the per-sector claim-tier ledger.
\end{enumerate}

\def\OPHSkipEmbeddedReferences{1}
\section{Recovered Core: Foundations and Structural Branches}\label{observer-patch-holography}

OPH is an observer-centric reconstruction program for fundamental physics. Physical data live on a horizon screen \(S^2\), and reality is encoded by the mutual consistency of overlapping patch descriptions. This paper works with five OPH axioms, two explicit external inputs used only in the quantitative branches, and theorem-local technical premises stated where needed for the scaling, gauge, and transport arguments below.

\subsection{Overview}\label{overview}

The recovered core of the paper is the D1--D5 and D7--D9 chain: overlap-consistency normal form, collar and edge-center structure, the conditional Lorentz and Einstein branches, compact gauge reconstruction, and the realized Standard Model branch with exact hypercharge structure, \(N_g=3\), and \(N_c=3\). Here and throughout the synthesis, labels such as D1--D12 are internal documentary node labels for the OPH derivation ledger; the local key appears in Section~\ref{16-summary-complete-axiom-set}, the SM/GR derivation paper \emph{Recovering Relativity and the Standard Model from the OPH Package Rooted in Observer Consistency} gives the same ledger in shorter form~\cite{ophcompact}, and the public paper ledger has no separate D11 row. D6 is a separate input-dependent corollary tied to total screen capacity. D12 collects phenomenological continuations kept distinct from the recovered core.

The Lorentz/null-modular/Einstein chain is a scaling-limit statement with explicit fixed-cutoff control. The gauge chain reconstructs a compact group from the transportable refinement-directed sector colimit supplied by \textbf{T6} and then specializes to the realized MAR-admissible branch. The worldsheet/string material and the flavor, neutrino, hadron, dark-sector, and spectroscopy branches are kept separate from the recovered core.

Two external continuous inputs appear only where declared. The pixel area \(P \equiv a_{\rm cell}/\ell_P^2\) enters deferred quantitative branches, and the screen capacity \(N_{\mathrm{scr}} \equiv \log \dim \mathcal H_{\mathrm{tot}}\) enters the cosmological-capacity branch. They are implementation inputs, not additional axioms.

\subsection{Model and Axioms}\label{1-model-and-axioms}

\subsubsection{Observers and access model}\label{11-observers-and-access-model}

An observer \(O\) is a tuple \((P_O, \mathcal{A}(P_O), \rho_O, R_O)\) where:

\begin{itemize}
\tightlist
\item
  \(P_O \subset S^2\) is a connected screen patch (the observer\textquotesingle s access region).
\item
  \(\mathcal{A}(P_O)\) is the von Neumann algebra associated to \(P_O\).
\item
  \(\rho_O\) is the local state, obtained by restricting the global state to \(\mathcal{A}(P_O)\).
\item
  \(R_O\) is a finite record algebra generated by stable internal record projectors within \(P_O\).
\end{itemize}

Observers are internal patterns in the global state. Different observers correspond to different patches and their compatible marginals.

\subsubsection{Screen, patches, and algebra net}\label{12-screen-patches-and-algebra-net}

We work in a single static patch with a horizon screen \(S^2\). Each connected subregion \(P \subset S^2\) is assigned a von Neumann algebra \(\mathcal{A}(P)\). The net satisfies isotony:

\[
P \subset Q \implies \mathcal{A}(P) \subset \mathcal{A}(Q).
\]

A global state \(\omega\) is a positive linear functional on the inductive-limit algebra. Overlap consistency is imposed algebraically: for overlaps \(P_1 \cap P_2\), \(\omega\) restricted to \(\mathcal{A}(P_1 \cap P_2)\) is the same from either side.

\subsubsection{Five OPH axioms}\label{13-core-axioms}

This paper uses five OPH axioms. The only labeled technical premises used below are the canonical premises \textbf{T1}--\textbf{T6} stated in Section 1.4. Other theorem-local hypotheses are written out in prose where needed rather than assigned a separate label family.

\textbf{1. Screen Net.} A horizon screen \(S^2\) carries a net of algebras \(P \mapsto \mathcal{A}(P)\).

\textbf{2. Overlap Consistency.} Local states agree on shared observables for any overlap.

\textbf{3. Local MaxEnt and Refinement Stability.} At the UV scale the realized state maximizes entropy subject to a finite family of gauge-invariant local constraints, and the resulting family of states is stable under coarse-graining so symmetry-allowed relevant operators are not held at zero by unexplained fine tuning.

\textbf{4. Recoverable Generalized Entropy.} A generalized entropy functional exists on caps, obeys quantum focusing on null generators, and comes with the recoverability structure used throughout the manuscript: collar tripartitions have small conditional mutual information with controlled recovery maps, with stronger collar/null-strip hypotheses stated explicitly where needed.

\textbf{5. Minimal Admissible Realization (MAR).} On the admissible low-energy branch, the realized sector package is the lexicographically minimal one under the complexity vector \(C(\mathfrak S)\). MAR is an explicit structural-economy axiom on admissible realized branches, not a theorem derived from the preceding four axioms.

\subsubsection{External inputs and theorem-local technical premises}\label{14-assumptions-and-external-inputs}

The quantitative branches discussed here use exactly two external continuous inputs:

\[
P \equiv a_{\mathrm{cell}}/\ell_P^2 = 1.63094,
\qquad
N_{\mathrm{scr}} \equiv \log \dim \mathcal H_{\mathrm{tot}} \sim 10^{122}.
\]

Here \(P\) and \(N_{\mathrm{scr}}\) are external implementation inputs, not theorem targets. \(P\) fixes the local UV pixel scale used by the deferred forward D10 branch and its descendants, while \(N_{\mathrm{scr}}\) fixes the global screen-capacity scale used by the cosmological-capacity branch. Any retained capacity-level neutrino side estimate on that input is legacy bookkeeping rather than the live weighted-cycle theorem lane. These are implementation inputs. They are not extra axioms.

Later theorem statements use only the canonical technical-premise labels \textbf{T1}--\textbf{T6}. Fixed-cutoff regulator bookkeeping, quasi-local propagation, endpoint control, collar decomposition, and related branch-internal statements are cited directly from axioms and earlier theorems rather than through a second premise-label system.

\textbf{Local MaxEnt at the regulator scale.} At the regulator scale \(\ell_{\mathrm{UV}}\), the global state \(\omega\) maximizes von Neumann entropy subject to:

\begin{enumerate}
\def\labelenumi{\arabic{enumi}.}
\tightlist
\item
  A finite set \(\{O_a\}\) of gauge-invariant local operators, each supported on a ball of radius \(\le r_0 = O(\ell_{\mathrm{UV}})\).
\item
  Constraint equations \(\langle O_a(x) \rangle = c_a\) for each cell \(x\) in the UV lattice.
\item
  Optionally, a finite number of global constraints (total energy, charge).
\end{enumerate}

This is the minimal specification needed to derive the local Gibbs form of Lemma 2.6 from Axiom 3.

\textbf{Clarification (MaxEnt \(\neq\) thermal equilibrium).} MaxEnt here is \textbf{local state selection} given constraints, not "the universe is in thermal equilibrium." The Lagrange multipliers (inverse temperatures) may vary slowly in space and time. Non-equilibrium physics appears as gradients in these multipliers and as controlled violations of exact Markov additivity under the explicit mixing hypotheses stated later in Section 2.3. Equilibrium is an approximation regime with explicit error terms.

\textbf{Rotationally invariant constraint branch.} Constraint sets are \(\mathrm{SO}(3)\)-invariant on \(S^2\).

\textbf{Gauge as overlap redundancy at fixed cutoff.} Overlap identifications are not unique; the freedom that leaves overlap observables invariant forms a local groupoid.

\textbf{Central-defect subbranch.} On triple overlaps, when the failure of strict coherence is central, one has

\[
\varphi_{ij} \varphi_{jk} \varphi_{ki} = \mathrm{Ad}(z_{ijk}),\qquad z_{ijk} \in Z(\mathcal{A}_{ijk}).
\]

\textbf{Collar double-scaling hypothesis.} There exists a UV length \(\ell_{\mathrm{UV}}\) such that for any cap \(C\) and collar width \(\delta\), in the refinement limit \(\delta \to 0\) and \(\ell_{\mathrm{UV}} \to 0\) with \(\delta/\ell_{\mathrm{UV}} \to \infty\), the Markov error satisfies

\[
I(A_\delta:D_\delta \mid B_\delta)_\omega \le \varepsilon(\delta/\ell_{\mathrm{UV}}),
\qquad \varepsilon(x) \to 0 \ \text{as} \ x \to \infty.
\]

See Section 2.3 for the collar tripartition definitions and the regulated EC proof that yields this limit from the fixed-cutoff realized presentation.

\textbf{Geometric-subnet BW lift.} The continuum Lorentz statement is asked only on the extracted prime geometric subnet for caps, together with the controlled collar/refinement limit that carries the Markov and recovery remainders explicitly. In the canonical numbered list below, the only global item here is the scaling-limit scope clause \textbf{T2}; the geometric cap-pair extraction and ordered cut-pair rigidity inputs are theorem-local UV objects written out where used.

\textbf{Derived quasi-local propagation and modular-locality control.} At scale \(\ell_{\mathrm{UV}}\), the dynamics generated by the effective Hamiltonian \(H_{\mathrm{eff}}\) (from Lemma 2.6) has a finite Lieb-Robinson velocity \(v_{\mathrm{LR}}\): for local operators \(A, B\) supported on regions separated by distance \(d\),

\[
\|[A(t), B]\| \le c \|A\| \|B\| \min(|R_A|, |R_B|) e^{-(d - v_{\mathrm{LR}}|t|)/\xi}
\]

for \(d > v_{\mathrm{LR}}|t|\), where \(\xi = O(\ell_{\mathrm{UV}})\).

This is the branch-internal control statement that turns the quasi-local structure of the
local-Gibbs branch into explicit support control for time-evolved operators. It is not an extra
premise beyond the third OPH axiom.

\begin{quote}
\textbf{Approximate modular covariance on the geometric branch.} Under the fixed-cutoff realized presentation, Lemma 2.6 (local Gibbs), the derived quasi-local propagation bound above, and the collar double-scaling plus mixing hypotheses of Section 2.3, the modular flow \(\sigma_t^{\omega,C}\) maps \(\mathcal{A}(R)\) into a slightly thickened region algebra:

\[
\sigma_t^{\omega,C}(\mathcal{A}(R)) \subseteq \mathcal{A}(R^{+v_{\mathrm{mod}}|t|})
\quad \text{up to error } \eta(d - v_{\mathrm{mod}}|t|),
\]

where \(R^{+s}\) denotes the \(s\)-neighborhood thickening, \(v_{\mathrm{mod}}\) is a "modular propagation velocity" controlled by \(v_{\mathrm{LR}}\) and local norm bounds, and \(d\) is the distance from \(R\) to \(\partial C\).

\textbf{Heuristic motivation.} The modular Hamiltonian \(K_C = -\log \rho_C\) is quasi-local by Lemma 4.1a-b (modular additivity localizes it to the collar). The derived Lieb-Robinson control then bounds support spreading under \(e^{iK_C t}\). In the double-scaling collar limit (\(\delta/\ell_{\mathrm{UV}} \to \infty\)), the thickening vanishes in macroscopic units. This is heuristic support for the tangent-limit package used later.

\textbf{Continuum-limit heuristic.} Define the induced region flow

\[
f_t^C(R) := \lim_{\ell_{\mathrm{UV}} \to 0} R^{+v_{\mathrm{mod}}|t|}.
\]

Then \(\sigma_t^{\omega,C}(\mathcal{A}(R)) = \mathcal{A}(f_t^C(R))\) becomes exact in the continuum limit, with error controlled by

\[
\eta(\delta) \lesssim 2\sqrt{\ln 2 \cdot c \cdot |\partial C|_{\mathrm{UV}}} \, e^{-\delta/(2\xi)}.
\]
\end{quote}

The discussion above is heuristic support for the tangent-limit package used in Theorem 4.2.

\textbf{Refinement-stable MaxEnt branch.} There exists a family of coarse-graining channels \(\Phi_{\ell\to L}\) between UV scale \(\ell\) and IR scale \(L\) such that the MaxEnt-selected states are self-similar under refinement,

\[
\Phi_{\ell\to L}(\omega_{\ell}) = \omega_{L},
\]

with the constraint set fixed and finite. Because the same finite constraint family is preserved, the regulator states lie in one common finite-dimensional multiplier family, and the realized branch is the persistent trajectory or invariant subset selected inside that multiplier family under the induced refinement maps. This is a state-side persistence statement; it does not by itself upgrade fixed-cutoff edge labels to a transportable refinement-persistent sector colimit.

\textbf{Fixed-cutoff realized presentation.} At a UV scale \(\ell_{\mathrm{UV}}\), local patch algebras are type-I with finite-dimensional Hilbert spaces, and gauge-as-gluing is realized as a boundary action on a lifted finite-dimensional presentation. Section 2.3 proves the resulting fixed-cutoff collar package directly from overlap consistency on this realized presentation, while distinguishing the lifted fixed-point algebra from the sector-preserving collar algebra used in the EC theorem.

External mathematical inputs: SSA and recovery theorems (Petz 1986, 1988; Fawzi and Renner 2015), Jacobson\textquotesingle s entanglement-equilibrium derivation (Jacobson 1995, 2016), and one of the following EFT bridges: (i) the null-surface modular route (Section 5.2), or (ii) a UV CFT regime on sufficiently small caps (Section 5.3). For SM contact we also use the Doplicher-Roberts reconstruction (Doplicher and Roberts 1989, 1990) once localized transportable sectors are assumed in the small-region limit. Full citations appear in the References.

\textbf{Canonical technical premises.} The only labeled technical premises used in this manuscript are:
\begin{itemize}
\tightlist
\item \textbf{T1}: the loop-coherence branch criterion kept explicit whenever the ordinary compact-group route is invoked, written \([z]=0\) on the central-defect branch; on the genuinely noncentral fixed-cutoff branch the parallel criterion is \(q_\Sigma=0\), but that branch is handled by crossed-module data rather than by the ordinary DR field-algebra route.
\item \textbf{T2}: all Lorentz/null-modular/Einstein statements are claims about a controlled refinement/scaling limit of the regulator net, not literal exactness at fixed finite cutoff.
\item \textbf{T3}: fixed-cap generalized-entropy stationarity for the admissible first-variation class used in the Jacobson branch.
\item \textbf{T4}: symmetric braiding in the \(3+1\)D EFT branch.
\item \textbf{T5}: a bosonic Tannakian fiber-functor premise, or else an explicit super-Tannakian fork.
\item \textbf{T6}: whenever compact gauge reconstruction is invoked, the sector category means the directed colimit of transportable edge sectors with objectwise finite-dimensional fibers.
\end{itemize}

\subsubsection{Notation}\label{15-notation}

\begin{itemize}
\tightlist
\item
  \(\rho_C\): reduced state on cap \(C\).
\item
  \(K_C := -\log \rho_C^{\omega}\): modular Hamiltonian of the reference state.
\item
  \(B_C\): geometric generator of the cap-preserving conformal dilation.
\item
  \(S_{\mathrm{gen}}(C)\): generalized entropy on a cap.
\item
  \(\ell_{\mathrm{UV}}\): UV length scale of the refined screen net.
\item
  \(\delta\): collar width around a cap boundary.
\end{itemize}

\subsubsection{Summary: premise ledger and dependency DAG}\label{16-summary-complete-axiom-set}

For reference, the paper uses the five OPH axioms above, the two external inputs \(P\) and \(N_{\mathrm{scr}}\) only in the branches that require them, and the theorem-local technical premises declared in Section 1.4.

{\def\LTcaptype{none}
\tiny
\setlength{\tabcolsep}{1pt}
\begin{longtable}{@{}>{\RaggedRight\arraybackslash}p{0.16\textwidth}>{\RaggedRight\arraybackslash}p{0.76\textwidth}@{}}
\toprule\noalign{}
Item & Meaning, typical use, and status \\
\midrule\noalign{}
\endhead
\bottomrule\noalign{}
\endlastfoot
Axiom 1 & Screen Net; use: global net structure; status: axiom. \\
Axiom 2 & Overlap Consistency; use: compatibility on overlaps; status: axiom. \\
Axiom 3 & Local MaxEnt and Refinement Stability; use: state selection and RG stability; status: axiom. \\
Axiom 4 & Recoverable Generalized Entropy; use: focusing plus recoverability structure; status: axiom. \\
Axiom 5 & Minimal Admissible Realization (MAR); use: selection on admissible gauge branches; status: axiom. \\
\(P\) & Pixel area \(a_{\mathrm{cell}}/\ell_P^2\); use: deferred quantitative forward input on the D10 branch; status: external input. \\
\(N_{\mathrm{scr}}\) & Total screen capacity; use: cosmological-capacity branch; status: external input. \\
\textbf{T1} & Loop-coherence branch criterion for the gauge route; use: central ordinary compact-group reconstruction, with the higher-gauge \(q_\Sigma\) analogue tracked separately at fixed cutoff; status: canonical technical premise when invoked. \\
\textbf{T2} & Controlled refinement/scaling-limit scope clause; use: Lorentz/null-modular/Einstein branch; status: canonical technical premise. \\
\textbf{T3} & Fixed-cap generalized-entropy stationarity; use: Jacobson branch; status: canonical technical premise. \\
\textbf{T4} & Symmetric braiding in the \(3+1\)D EFT branch; use: compact gauge reconstruction; status: canonical technical premise. \\
\textbf{T5} & Bosonic Tannakian fiber functor or explicit super-Tannakian fork; use: compact gauge reconstruction; status: canonical technical premise. \\
\textbf{T6} & Directed colimit of transportable edge sectors with objectwise finite-dimensional fibers; use: compact gauge reconstruction; status: canonical technical premise. \\
\end{longtable}
}

\textbf{Master dependency map.} The completion program is organized into the following nodes. The point of the table is to make explicit which results use the OPH axioms alone, which import standard mathematics, and which require extra physical premises or external inputs.

\begingroup
\scriptsize
{\def\LTcaptype{none}
\setlength{\tabcolsep}{1pt}
\begin{longtable}[]{@{}>{\RaggedRight\arraybackslash}p{0.07\textwidth}>{\RaggedRight\arraybackslash}p{0.18\textwidth}>{\RaggedRight\arraybackslash}p{0.13\textwidth}>{\RaggedRight\arraybackslash}p{0.13\textwidth}>{\RaggedRight\arraybackslash}p{0.17\textwidth}>{\RaggedRight\arraybackslash}p{0.12\textwidth}@{}}
\toprule\noalign{}
Node & Main-manuscript output & Immediate OPH ingredients & Standard mathematics used & Extra physical premises / external inputs & Claim tier \\
\midrule\noalign{}
\endhead
\bottomrule\noalign{}
\endlastfoot
\textbf{D1} & Theorem 3.1: schedule-independent normal-form statement for the declared recovery-derived repair relation & overlap-consistency setup on finite patch nets together with the declared fixed-cutoff collar recovery step, its touched-overlap local-fit acceptance contract, support-local disjoint commutation, and restriction-compatible union-collar gluing on the physical quotient & local-to-global Lyapunov descent on the inconsistency potential; local-diamond completion from overlap-associative center-sector / higher-gauge union-collar gluing; Newman's lemma & repair completeness; stated Petz support/CPTP control where that branch is used & structural theorem \\
\textbf{D2} & Section 2.3 and \S5.4: EC/Markov collar and generalized-entropy split & Axioms 1--4 plus the derived fixed-cutoff regulator/collar package when invoked & HJPW; Petz; Fawzi--Renner & exact-Markov or idealized recoverability hypotheses when exact identities are used; coarse-grained \(A/(4G)\) dictionary & structural lemma/\newline proposition layer \\
\textbf{D3} & Theorems 4.2--4.3: geometric modular flow and Lorentz branch & Axioms 1--4 together with SO(3)-invariant constraint data & Bisognano--Wichmann modular geometry;\newline conformal classification on \(S^2\) & T2 together with the theorem-local geometric cap-pair extraction and ordered cut-pair rigidity hypotheses of Theorem 4.2 & conditional scaling-limit\newline theorem \\
\textbf{D4} & Section 5.2: null modular bridge & D2+D3 & Borchers--Wiesbrock and endpoint differentiation & the theorem-local inherited-strip decomposition and exact-or-controlled Markov hypotheses of Section 5.2, plus the same-half-line local modular-Hamiltonian input; bounded-interval formulas still require the separate interval-preserving projective branch & conditional bridge theorem \\
\textbf{D5} & Theorem 5.1: Einstein branch & D3+D4 & Jacobson small-ball mathematics; null-to-tensor reconstruction modulo a metric term & T3 together with the locally Lorentzian \(d=4\) scaling regime and the separate interval-preserving projective branch used to transport the D4 half-line identification to bounded intervals & conditional scaling-limit theorem \\
\textbf{D6} & capacity/\(\Lambda\) reduction statements & D5 leaves \(+\Lambda g_{ab}\) undetermined & de Sitter entropy relation and dimensional analysis & external input \(N_{\mathrm{scr}}\) and the screen-capacity identification & input-dependent corollary / estimate \\
\textbf{D7} & Theorem 6.1: compact gauge reconstruction & Axioms 1--4 & Doplicher--Roberts / Tannaka reconstruction & T1 and T4--T6 on the ordinary or central-defect bosonic branch; the genuinely noncentral fixed-cutoff branch is handled separately by crossed-module data & conditional structural theorem \\
\textbf{D8} & Section 6.2: product gauge structure up to finite quotient & D7 + Axiom 5 (MAR) & compact Lie representation classification; Schur's lemma & the same ordinary or central-defect bosonic branch as D7, together with a connected admissible class, one connected abelian factor, and faithful action on the minimal coupled carrier & realized-branch theorem \\
\textbf{D9} & Theorem 6.13, Proposition 6.9, Theorem 6.14, Proposition 6.6: hypercharge lattice, the counting chain \(N_g=3\) then \(N_c=3\), and the \(\mathbb Z_6\) quotient & D8 & anomaly algebra; Witten anomaly; CKM CP counting & realized one-generation/one-Higgs package and MAR admissibility premises & realized-branch theorem/corollary chain \\
\textbf{D12} & downstream flavor,\newline dark-sector,\newline baryogenesis,\newline spectroscopy, and\newline the string/worldsheet lane beyond the fixed-cutoff heat-kernel / Yang--Mills form on the declared compact-gauge branch: conditional large-\(N_{\mathrm{edge}}\) worldsheet-like reorganization and conjectural critical-superstring extension & various subsets of D6 and D9 & branch-specific EFT\newline and phenomeno\-logical\newline manipulations & explicit additional ans\"atze, including any controlled\newline large-\(N_{\mathrm{edge}}\) regime and any extra worldsheet/CFT premises & phenomeno\-logical\newline continuation \\
\end{longtable}
}
\endgroup

In particular, D3--D5 are scaling-limit branches rather than fixed-cutoff matrix-algebra theorems; D7--D9 are realized-branch structural outputs in the bosonic internal-gauge sector; D6 is input-dependent; and D12 remains phenomenological.

\begin{center}\rule{0.5\linewidth}{0.5pt}\end{center}
\subsection{Information-Theoretic Tools}\label{2-information-theoretic-tools}

\subsubsection{Strong subadditivity and Markov states}\label{21-strong-subadditivity-and-markov-states}

For any tripartite state \(\rho_{ABC}\),

\[
I(A:C \mid B) := S(AB) + S(BC) - S(B) - S(ABC) \ge 0.
\]

Exact Markov states satisfy \(I(A:C \mid B) = 0\) and admit a recovery map:

\[
\rho_{ABC} = (\mathrm{id}_A \otimes \mathcal R_{B\to BC})(\rho_{AB}).
\]

\subsubsection{Approximate recovery}\label{22-approximate-recovery}

If \(I(A:C \mid B) \le \varepsilon\) (bits), there exists a CPTP recovery map \(\mathcal{R}\) with

\[
\| \rho_{ABC} - (\mathrm{id}_A \otimes \mathcal R)(\rho_{AB}) \|_1
\le 2 \sqrt{\ln 2\, \varepsilon}.
\]

\subsubsection{Collar refinement and sufficient mechanisms}\label{23-collar-refinement-and-sufficient-mechanisms}

Fix a cap \(C \subset S^2\) with boundary circle. For collar width \(\delta\) define

\[
B_\delta := \{x \in S^2 : \mathrm{dist}(x,\partial C) \le \delta\},
\qquad
A_\delta := C \setminus B_\delta,
\qquad
D_\delta := (S^2 \setminus C) \setminus B_\delta.
\]

Then \(S^2 = A_\delta \cup B_\delta \cup D_\delta\) with \(A_\delta\) and \(D_\delta\) interacting only through \(B_\delta\). The collar double-scaling hypothesis is the requirement that \(I(A_\delta:D_\delta \mid B_\delta) \to 0\) in the refinement limit. At fixed regulator scale, the same finite-dimensional realized presentation supplies the patch-net and overlap-gluing data used below. We therefore isolate that realized presentation first, then separate the exact-Markov and quantitative routes.

\textbf{Derived regulator realization.} Choose a finite UV cellulation at scale \(\ell_{\mathrm{UV}}\) and let \(R\) be a finite union of cells. Hilbertize each cell\textquotesingle s finite local data to a finite-dimensional space \(\tilde{\mathcal H}_i\cong \mathbb C^{n_i}\). The extended algebra before quotienting by overlap redundancy is

\[
\widetilde{\mathcal A}(R)=\mathcal B(\tilde{\mathcal H}_R),
\qquad
\tilde{\mathcal H}_R=\bigotimes_{i\subset R}\tilde{\mathcal H}_i.
\]

Overlap-preserving changes of local trivialization act only on cut data. On each finite-dimensional chart, their unitary image has compact closure, so one may represent the boundary gluing redundancy by a compact group \(G_{\partial R}\subset U(\tilde{\mathcal H}_R)\). The lifted boundary-invariant endomorphism algebra is

\[
A_{\mathrm{inv}}(R)=\widetilde{\mathcal A}(R)^{G_{\partial R}}
=\mathcal B(\tilde{\mathcal H}_R)^{G_{\partial R}}.
\]

This is the fixed-cutoff realized presentation used throughout the collar analysis. The later EC theorem works on the invariant-state realization and uses the sector-preserving induced collar algebra there rather than the entire fixed-point algebra on the unreduced tensor product.
When the consensus paper later speaks of quotient repair, the physical repair law is defined on
this fixed-point / overlap-invariant data and any representative-level map is only a lift of that
quotient update. Descent to the gauge quotient is therefore built into the physical-algebra
formulation, not added later as a separate covariance axiom on hidden representatives. On the
declared fixed-cutoff collar branch, the local repair step itself is read from exact Markov splice
or a declared Petz/Fawzi--Renner recovery channel. What remains above that step on the declared
fixed-cutoff branch is repair completeness, together with control on the declared Petz domain
where that branch is used. The declared touched-overlap acceptance contract makes \(\Phi\) the
finite-patch Lyapunov functional for accepted moves, and the support-local disjoint-commutation
plus restriction-compatible union-collar glue that turn overlap associativity into the local
diamond are part of the declared repair law itself. On the same fixed-cutoff quantum lift, the
terminal expectation functional on each declared physical observable algebra is therefore unique
even when microscopic representative lifts differ by gauge labels globally or by sector labels
inside one quotient-local glued state.
Stability of that compatibility package under
refinement or branch change is a separate question.

\textbf{Theorem 2.3 (EC from regulated overlap gluing).} For a collar \(B_\delta\) around a cap boundary \(\Sigma\), there is a canonical decomposition

\[
H_{B_\delta} = \bigl(\tilde{\mathcal H}_{B_L}\otimes \tilde{\mathcal H}_{B_R}\bigr)^{G_\Sigma}
= \bigoplus_{\alpha}
\bigl(H_{b_L^{\alpha}} \otimes H_{b_R^{\alpha}}\bigr),
\]

and the sector-preserving collar algebra
\[
A_{\mathrm{EC}}(B_\delta):=\bigoplus_\alpha
\bigl(B(H_{b_L^\alpha}) \otimes B(H_{b_R^\alpha})\bigr)
\]

with

\[
Z(A_{\mathrm{EC}}(B_\delta)) = \bigoplus_\alpha \mathbb C\,\mathbf 1_\alpha,
\]

such that \(\mathcal{A}(A_\delta B_\delta)\) acts only on \(H_{b_L^\alpha}\) and \(\mathcal{A}(B_\delta D_\delta)\) acts only on \(H_{b_R^\alpha}\) within each block.

\textbf{Proof.} Split the collar into half-collars \(B_L\) and \(B_R\) meeting on \(\Sigma = \partial C\). By the realized regulator presentation above, the physical collar Hilbert space is the diagonal invariant subspace \((\tilde{\mathcal H}_{B_L} \otimes \tilde{\mathcal H}_{B_R})^{G_\Sigma}\). Decompose each side into irreps:

\[
\tilde{\mathcal H}_{B_L} = \bigoplus_\alpha (V_\alpha \otimes H_{b_L^\alpha}),\qquad
\tilde{\mathcal H}_{B_R} = \bigoplus_\beta (V_\beta^* \otimes H_{b_R^\beta}).
\]

Then

\[
\tilde{\mathcal H}_{B_L} \otimes \tilde{\mathcal H}_{B_R}
= \bigoplus_{\alpha,\beta}
(V_\alpha \otimes V_\beta^*) \otimes
(H_{b_L^\alpha} \otimes H_{b_R^\beta}).
\]

By Schur\textquotesingle s lemma,

\[
(V_\alpha \otimes V_\beta^*)^{G_\Sigma}
\cong
\begin{cases}
\mathbb C, & \alpha=\beta,\\
0, & \alpha\ne\beta.
\end{cases}
\]

Therefore the invariant subspace is

\[
H_{B_\delta} = \bigoplus_\alpha (H_{b_L^\alpha} \otimes H_{b_R^\alpha}),
\]

as claimed. The induced sector-preserving collar algebra is

\[
A_{\mathrm{EC}}(B_\delta) = \bigoplus_\alpha
\bigl(B(H_{b_L^\alpha}) \otimes B(H_{b_R^\alpha})\bigr),
\]

so the center is generated by the block projectors. Adjacent region algebras act on the left or right factor only because the gluing action is supported on \(\Sigma\). QED.

\textbf{Remark.} On the central-defect subbranch, replace \(G_\Sigma\) by its central extension. The sector label \(\alpha\) then ranges over irreps of the extension; the decomposition is unchanged.

We refer to the decomposition in Theorem 2.3 as \textbf{edge-center completion (EC)}.

\textbf{Exact Markov route.} If the reference state on \(A_\delta B_\delta D_\delta\) is exact Markov, or in an explicitly stated idealized limit that reduces to exact Markovity, then

\[
\rho_{A_\delta B_\delta D_\delta}
= \bigoplus_\alpha p_\alpha
\bigl(\rho_{A_\delta b_L^\alpha} \otimes \rho_{b_R^\alpha D_\delta}\bigr),
\qquad
I_\omega(A_\delta:D_\delta \mid B_\delta)=0.
\]

This is the HJPW normal form applied to the EC blocks, and it is the exact identity used for
literal Markov-modular equalities.

Interpreting collar refinement as the inductive limit of these regulators with \(\delta/\ell_{\mathrm{UV}} \to \infty\), Theorem 2.3 supplies the kinematic edge-center decomposition. Exact Markovity is one idealized route. The following lemma and axiom provide the quantitative decay route used when the manuscript keeps the approximation explicit.

\begin{quote}
\textbf{Lemma 2.6 (MaxEnt with local constraints implies local Gibbs form).} Under the local MaxEnt branch of Axiom 3 and the finite-dimensional regulator realization above, the MaxEnt state has the Gibbs form

\[
\omega = \frac{e^{-H_{\mathrm{eff}}}}{\mathrm{Tr}\,e^{-H_{\mathrm{eff}}}},
\qquad
H_{\mathrm{eff}} = \sum_x \sum_a \lambda_a O_a(x) + \text{(global terms)},
\]

where the sum runs over UV cells \(x\) and constraint operators \(O_a\). The effective Hamiltonian \(H_{\mathrm{eff}}\) is quasi-local with range \(O(\ell_{\mathrm{UV}})\).

\textbf{Proof.} On a finite-dimensional algebra, the unique state maximizing \(S(\rho) = -\mathrm{Tr}(\rho \log \rho)\) subject to linear constraints \(\mathrm{Tr}(\rho O_i) = c_i\) is given by Lagrange multipliers:

\[
\rho = \frac{e^{-\sum_i \lambda_i O_i}}{\mathrm{Tr}\,e^{-\sum_i \lambda_i O_i}}.
\]

Strict concavity of von Neumann entropy ensures uniqueness. When the constraints are "translated local" (the same \(O_a\) at each cell \(x\)), the exponent is a sum of local terms. QED.
\end{quote}

\textbf{Exponential mixing hypothesis.} There exist constants \(c\) and correlation length \(\xi = O(\ell_{\mathrm{UV}})\) such that

\[
I_\omega(A_\delta:D_\delta \mid B_\delta) \le
c\,\lvert\partial C\rvert_{\mathrm{UV}}\,e^{-\delta/\xi},
\qquad
\lvert\partial C\rvert_{\mathrm{UV}} \sim \frac{\mathrm{length}(\partial C)}{\ell_{\mathrm{UV}}}.
\]

This is the standard fixed-cutoff clustering/mixing condition used for local Gibbs states. One
sufficient realization on a chosen fixed-cutoff collar family is a Dobrushin-type uniqueness estimate
for the corresponding local Gibbs specifications, or a spectral-gap estimate uniform across that
family. Used here, it is
a collar-local recoverability condition on the selected branch; it does not imply infinite-volume
uniqueness and it does not decide whether a separately supplied refinement-limit gauge-sector
colimit is trivial or nontrivial.

\begin{quote}
\textbf{Theorem 2.5 (Local Gibbs + mixing implies collar refinement).} Under Lemma 2.6 and the exponential mixing hypothesis above, the collar double-scaling hypothesis holds in the limit \(\delta \to 0\), \(\ell_{\mathrm{UV}} \to 0\) with \(\delta/\ell_{\mathrm{UV}} \to \infty\).

\textbf{Proof.} The mixing bound above has polynomial growth in \(\lvert\partial C\rvert_{\mathrm{UV}}\) and exponential decay in \(\delta/\ell_{\mathrm{UV}}\). In the double-scaling limit the exponential dominates, so \(I_\omega(A_\delta:D_\delta \mid B_\delta) \to 0\). QED.
\end{quote}

This bound is the quantitative hinge for constructive gluing.

\subsubsection{Concrete UV realization: quantum link models}\label{26-concrete-uv-realization-quantum-link-models}

Section 2.3 internalizes the regulator package: the fixed-cutoff type-I algebra and boundary fixed-point structure are the realized form of the screen net plus overlap gluing. Quantum link models are included here only as an explicit microscopic example of that structure, not as a separate axiom layer.

\textbf{UV regulator.} Triangulate \(S^2\) at scale \(\ell_{\mathrm{UV}}\), giving vertices \(v\), oriented links \(\ell\), and plaquettes \(p\). Refinement corresponds to \(\ell_{\mathrm{UV}} \to 0\) with increasing lattice size.

\textbf{Degrees of freedom.} Attach to every oriented link \(\ell\) a \textbf{finite-dimensional} Hilbert space \(\mathcal{H}_\ell\). In ordinary Wilson lattice gauge theory, the continuum/refinement-limit edge description is modeled by \(L^2(G)\) (infinite-dimensional for continuous \(G\)), but the microscopic OPH regulator is instead a \textbf{quantum link model} with finite-dimensional link Hilbert spaces that preserve gauge symmetry in operator form~\cite{cw97}. Optionally attach matter Hilbert spaces \(\mathcal{H}_v\) at vertices. Then:

\[
\tilde{\mathcal{H}}_{\mathrm{total}} = \bigotimes_\ell \mathcal{H}_\ell \otimes \bigotimes_v \mathcal{H}_v,
\]

finite-dimensional on any finite lattice. This is a concrete realization of the extended type-I presentation used in Section 2.3.

\textbf{Boundary gluing as Gauss-law invariants.} Define a local gauge transformation group \(G_v\) at each vertex \(v\) acting on incident links (and matter at \(v\)). Physical states satisfy:

\[
|\psi\rangle \in \mathcal{H}_{\mathrm{phys}} \quad\Longleftrightarrow\quad U(g_v)|\psi\rangle = |\psi\rangle \;\;\forall\, v,\, g_v \in G_v.
\]

Equivalently: \(\mathcal{H}_{\mathrm{phys}} = \tilde{\mathcal{H}}_{\mathrm{total}}^{\prod_v G_v}\).

\textbf{Region algebras.} For any region \(R \subset S^2\), define an extended Hilbert space \(\tilde{\mathcal{H}}_R\) from the links/vertices in \(R\). The \textbf{boundary gauge group} \(G_{\partial R}\) acts on the cut degrees of freedom (the "half-links" ending on \(\partial R\)). Define:

\[
\mathcal{A}_{\mathrm{inv}}(R) = \mathcal{B}(\tilde{\mathcal{H}}_R)^{G_{\partial R}}.
\]

\textbf{This is the concrete quantum-link instance of the lifted fixed-point presentation used in Section 2.3.} The same gauge constraints then induce the sector-preserving EC algebra on collars.

\textbf{Why EC follows immediately.} Take a cap \(C\) and a collar \(B_\delta\) around \(\partial C\). Because the \emph{only} coupling between inside and outside is through the boundary gauge constraint, the collar Hilbert space decomposes into superselection blocks labeled by boundary irreps:

\[
\mathcal{H}_{B_\delta} \cong \bigoplus_\alpha (H_{b_L^\alpha} \otimes H_{b_R^\alpha}),
\]

with center generated by the projectors \(P_\alpha\). This is precisely the Schur-lemma mechanism of Theorem 2.3. The labels \(\alpha\) are the familiar edge-mode / electric-flux labels appearing whenever one factorizes gauge theories across an entangling cut~\cite{donnellyWall15}. Exact Markovity depends on the state and is not forced by the decomposition alone.

\textbf{Dynamics and MaxEnt.} The natural Hamiltonian is a 2+1D lattice gauge Hamiltonian on the screen worldvolume: plaquette ("magnetic") terms, electric terms on links, vertex Gauss terms as constraints, plus local matter couplings. In quantum link form this remains finite-dimensional per link while behaving like gauge theory in the continuum limit. Then the MaxEnt assumption becomes concrete: the MaxEnt state is a Gibbs state \(\rho \propto e^{-\sum_i \lambda_i O_i}\) with quasi-local \(O_i\), precisely the local-Gibbs regime.

\textbf{Geometry and \(G\).} This microphysics naturally supplies the emergent geometric objects:

\begin{itemize}
\tightlist
\item
  \textbf{Edge entropy / area operator:} \(L_C = \sum_\alpha (\log d_\alpha) P_\alpha\) becomes "log of boundary irrep dimension" in the gauge link model.
\item
  \textbf{Newton constant \(G\):} the conversion factor between edge entropy density per boundary UV cell and macroscopic geometric area.
\end{itemize}

Thus area is an operator living in the center of the boundary algebra, because in gauge systems the center is where the cut labels live.

\textbf{What this example does and does not close.} The quantum-link realization makes the fixed-cutoff matrix/fixed-point bookkeeping concrete. What it does \textbf{not} automatically guarantee is that modular flow on caps becomes geometric conformal dilation with the \(2\pi\) KMS normalization. The remaining continuum burden is the \textbf{T2} scaling-limit package together with the theorem-local geometric-subnet extraction and ordered cut-pair rigidity objects stated later. Viable architectures for this include holographic quantum error-correcting codes~\cite{pastawski15} and quantum double / string-net Hamiltonians~\cite{levinWen05}.

\subsubsection{Conformal-modular fixed-point microphysics}\label{27-conformal-modular-fixed-point-microphysics-cmfp}

After Section 2.3, the remaining gap is narrower. On the local finite-constraint MaxEnt branch of the third OPH axiom, the logarithm of the selected state is a quasi-local UV generator, and the refinement-stable branch lies inside one common finite-dimensional multiplier family. The unresolved question is which algebraic phase the refinement limit realizes. At each regulator stage the patch and cap algebras are type-I matrix algebras, but the finite regulator class is not refinement-closed, so the scaling-limit observer algebra may leave that class. The fixed-cutoff analysis therefore proves a local-Gibbs/refinement-stable branch with propagation and endpoint control, but it does \emph{not} prove that every such branch emits the prime geometric cap pair with geometric modular flow. The conditional boundary is the \textbf{T2} scaling-limit clause together with the theorem-local geometric-subnet extraction and ordered cut-pair rigidity inputs: the target algebra is the extracted geometric subnet, while the internal extraction of its realized scaling-limit cap pair depends on the derived carried-collar schedule above the transported fixed-local-collar Markov/faithfulness datum and the later weak-\(*\) extraction plus GNS gluing step.

\textbf{Local finite-constraint MaxEnt branch.} The constraint family \(\mathcal C\) is generated by finitely many gauge-invariant local densities \(\{O_a(x)\}\) of UV range \(O(\ell_{\mathrm{UV}})\), with the same finite label set retained under refinement. This is not an extra postulate beyond the third OPH axiom; it is that axiom unpacked at regulator level.

\textbf{Theorem 2.6 (Local constraints imply a local-Gibbs form).} If the MaxEnt constraints are expectations of finitely many quasi-local operators \(\{O_a\}\) with bounded support size at scale \(\ell_{\mathrm{UV}}\), then the entropy maximizer is
\[
\omega \propto \exp\!\left(-\sum_a \lambda_a O_a\right),
\]
so the MaxEnt generator \(H_{\mathrm{MaxEnt}}=-\log\omega\) is a UV-range quasi-local sum. This is exactly the local-Gibbs form used later.

\textbf{Proof.} Standard exponential-family result: maximum entropy subject to linear constraints \(\langle O_a\rangle=c_a\) yields the Gibbs state with Lagrange multipliers \(\lambda_a\). QED.

\textbf{Derived propagation control on the same branch.} Because \(H_{\mathrm{MaxEnt}}\) is a finite-range or quasi-local sum on a finite type-I regulator net, standard Lieb--Robinson estimates~\cite{liebRobinson72} apply to the automorphism group it generates, or to any branch generator lying in the same bounded-support algebraic closure. Thus there is a finite propagation velocity \(v_{\mathrm{LR}}\) and constants \(C,\xi\) such that for local observables \(A_X,B_Y\),
\[
\bigl\|[\tau_t(A_X),B_Y]\bigr\|
\le
C\,\|A_X\|\,\|B_Y\|\,\min(|X|,|Y|)\,
e^{-(d(X,Y)-v_{\mathrm{LR}}|t|)/\xi}.
\]
Because changing a bounded interval endpoint only adds or removes an \(O(|\Delta v|)\) collar of local terms once the central endpoint piece is separated off, the same local branch also gives bounded-interval endpoint-Lipschitz matrix elements,
\[
\left|\langle\psi,(K[I']-K[I])\phi\rangle\right|
\le
C_{\psi,\phi,I_{\max}}\,|I'\Delta I|,
\]
used later in the null-modular bridge. So quasi-local propagation and endpoint control are not external regularity selectors; they are the local-constraint MaxEnt branch written in dynamical form.

\textbf{Refinement-stable multiplier branch.} Because the same finite constraint family is preserved under coarse-graining, the regulator states lie in one common finite-dimensional multiplier family rather than in unrelated state spaces at different cutoffs. The ``refinement-stable'' language used later therefore means persistence along the stable or fixed branch of this multiplier family. This state-side notion is enough to compare realized states across cutoffs. It does \emph{not} by itself upgrade fixed-cutoff edge labels to transportable refinement-persistent sectors, and it does not show that the directed colimit of such sectors is nontrivial. Accordingly, whenever the gauge derivation speaks of a refinement-stable directed colimit of sectors, the entire transportable directed-system/fiber clause is the explicit content of \textbf{T6}; the third OPH axiom supplies only the realized state branch along which one asks whether such sector transport persists.

\textbf{Scaling limit and algebraic type.} The regulator presentation gives a family of finite type-I algebras
\[
\mathcal A_{\ell}(C)\cong \mathcal B(\mathcal H_{C,\ell}).
\]
A scaling limit of this family need not remain type I. In the continuum-QFT case of interest one expects the local limit algebra \(\mathcal A_\infty(C)\) to be non-type-I, typically type III. Accordingly the fundamental modular datum in the limit is the automorphism group of the pair \((\mathcal A_\infty(C),\omega_\infty^C)\), not a density matrix inside \(\mathcal A_\infty(C)\).

\textbf{Proposition 2.6 (Geometric modular action on caps on the extracted prime geometric subnet).} Under the \textbf{T2} scaling-limit clause and the theorem-local geometric cap-pair extraction and ordered cut-pair rigidity hypotheses used later in Theorem 4.2, for any extracted scaling-limit geometric cap pair \((\mathcal A_\infty^{\mathrm{geo}}(C),\omega_\infty^{\mathrm{geo},C})\), let \(\alpha_{\lambda_C(s)}\) be the automorphism induced by the standard cap-preserving conformal subgroup \(\lambda_C(s)\). Then
\[
\sigma_t^{\omega_\infty^{\mathrm{geo},C}}
=
\alpha_{\lambda_C(2\pi t)}.
\]
If the limit cap algebra happens to remain type I, this may be written as \(K_C=2\pi B_C\). In the generic continuum case, the same statement is an outer modular action on a non-type-I algebra.

\textbf{Proof.} Ordered cut-pair rigidity on the extracted prime geometric cap pair identifies the scaling-limit cap modular group with the standard geometric cap subgroup up to normalization, and the modular KMS condition fixes the normalization to \(2\pi\). If the limit algebra is type I, the automorphism statement may be represented by a modular Hamiltonian \(K_C\); otherwise the automorphism statement is the full content. QED.

\textbf{Alternative derivation via net regularity.} The same modular-covariance property can also be read off from a scaling-limit support map when the net satisfies the outer-regularity / minimal-support condition used later. This does \emph{not} complete the BW lift by itself; it clarifies how the geometric labeling of the extracted limit net is read once the geometric cap pair has been emitted.

\textbf{(NR) Outer regularity / minimal support.} For any operator \(O\), the intersection of all connected regions \(P\) with \(O\in\mathcal A(P)\) is again a connected region, denoted \(\mathrm{supp}(O)\).

\textbf{Proposition 2.7 (Modular covariance from net regularity).} Under (NR), define for any region \(R\subset C\)
\[
f_t^C(R)
:=
\bigcup_{O\in \mathcal A(R)}
\mathrm{supp}\!\left(\sigma_t^{\omega,C}(O)\right).
\]
Then \(\sigma_t^{\omega,C}(\mathcal A(R))=\mathcal A(f_t^C(R))\), which is exactly the desired modular-covariance property.

\textbf{Proof.} Since \(\sigma_t^{\omega,C}\) is an automorphism of \(\mathcal A(C)\), and (NR) allows one to read support from the net labeling, the map \(R\mapsto f_t^C(R)\) is well-defined and consistent. QED.

\textbf{Null-surface modular structure.} On the same \textbf{T2} scaling branch and extracted geometric-subnet branch, the null-surface modular machinery narrows as follows:
\begin{itemize}
\item
\textbf{Fixed-cutoff null-strip bridge.} The null-strip package proves transferred cut-center data, theorem-local inherited left/right strip structure, exact-or-controlled four-term strip additivity on one inherited strip model, renormalized endpoint control up to the weak tail generator, and the derived half-sided modular pair whose Borchers--Wiesbrock consequence is an explicit positive null-translation generator on its Stone domain.
\item
\textbf{Derived half-sided modular inclusion.} Nested null half-line algebras satisfy half-sided modular inclusion on the declared geometric scaling branch by Corollary 5.2e; Borchers--Wiesbrock then identifies the positive null-translation generator on its Stone domain together with the affine half-line modular relation~\cite{wiesbrock93}.
\item
\textbf{Weak continuity and finite variation.} The bounded-interval and half-line endpoint control follow from the local MaxEnt branch and define the weak tail generator at fixed cutoff. The \textbf{T2} scaling branch together with the extracted geometric cap pair supplies the continuum null-generator setting in which that weak-tail data can be matched to the geometric null modular action of the relativistic phase.
\end{itemize}

\textbf{Constraint set specification.} On the local finite-constraint branch, the ``correct fixed-cap constraint set'' becomes explicit: constraints are the local conserved charges of the symmetries used in the derivation:
\begin{enumerate}
\def\labelenumi{\arabic{enumi}.}
\item
\textbf{Edge/cap label constraints:} fix the distribution of collar-sector labels, equivalently \(\langle L_C\rangle\) for each cap size, giving the area term.
\item
\textbf{Gauge charges:} fix boundary flux or charge operators.
\item
\textbf{Geometric (conformal) charges:} fix the expectation of the conformal Killing charges that preserve the cap, i.e.\ the generator \(B_C\) or its microscopic lattice approximation.
\end{enumerate}

MaxEnt therefore selects the unique finite-stage invariant state compatible with those conserved charges. What it does \emph{not} prove by itself is that the scaling-limit state is the vacuum or canonical cap state. The stronger statement is the \textbf{T2} scaling-limit clause together with the theorem-local extraction of a realized geometric cap pair and the later ordered cut-pair rigidity step. When the manuscript uses the phrase ``BW/canonical cap phase,'' it means only the case in which such a geometric cap pair has been emitted; if the limit algebra is non-type-I, the geometric modular action on that branch is outer.

\textbf{QNEC internalization.} QNEC has rigorous QFT proofs in broad settings~\cite{bousso16}. On the \textbf{T2} scaling branch with the extracted geometric cap pair in place, the Recoverable Generalized Entropy axiom can be supported internally by:
\begin{itemize}
\item
EC + MaxEnt derive \(S_{\mathrm{gen}}=S_{\mathrm{bulk}}+\langle L_C\rangle\) (Section 5.4).
\item
Once the downstream null-stress bridge and Einstein branch are in place, focusing becomes a semiclassical consequence in the same scaling regime.
\end{itemize}

\textbf{Summary.} The CMFP package therefore separates the dependency structure into two layers:
\begin{itemize}
\item
\textbf{Internal to OPH:} the local-Gibbs form, quasi-local propagation, bounded-interval endpoint control, and the realized state-side refinement branch along which later transport questions are asked, all from the local finite-constraint MaxEnt branch.
\item
\textbf{Theorem-external branch statement:} that the scaling limit emits the prime geometric cap pair and that ordered cut-pair rigidity holds on it. On that branch the limit algebra may be non-type-I, the modular action may be outer, and the \(2\pi\) normalization and later null half-sided-inclusion bridge follow.
\end{itemize}

The internal extension route is a two-step chain, and the pressure point is the first object. First, one needs a canonical scaling-limit observer cap-pair realization on the geometric subnet from transported marginals. The theorem contract visible in the manuscript is that the realized transported cap-local system packages the geometric cap-local test family, the projectively compatible transported marginal family, and the asymptotic transport-equivalence certificate. The remaining emitted witness is a derived vanishing carried-collar schedule on fixed local collars, explicitly decomposed into the constructive-recovery remainder and the faithful modular-defect term, namely convergence of the carried quantity
\[
r_{\mathrm{FR}}(\varepsilon_{n,m,\delta})
+
4\lambda_\ast^{-1}\,
\delta^{\mathrm M}_{A_{n,m,\delta}:B_{n,m,\delta}:D_{n,m,\delta}}
\bigl(\varepsilon_{n,m,\delta}\bigr)
\]
to zero after transport to each fixed local collar model; only then can local weak-\(*\) extraction and GNS gluing emit the realized scaling-limit geometric cap pair. Beneath that derived schedule, the sharper raw datum is the transported fixed-local-collar Markov/faithfulness package, and the single missing clause inside that lower datum is the eventual common floor on the finite modular-transport family feeding the faithful modular-defect term. No separate exact-Markov-reference faithfulness input is missing there: on one fixed collar model, the exact-Markov comparison marginals inherit the same eventual floor once the exact-Markov modulus tends to zero. The manuscript does not derive a refinement-uniform common floor on the finite modular-transport marginals; this single clause remains external to the emitted theorem chain. Second, one needs ordered cut-pair rigidity on that realized limit. That stage remains symbolic until the realized scaling-limit geometric cap pair exists. The weak-\(*\) extraction plus rigidity route is therefore an explicit internal extension route, not part of the proved theorem surface.
\subsection{Overlap Consistency and Gluing}\label{3-overlap-consistency-and-gluing}

The constructive part of overlap consistency is the tree-gluing theorem below. The structural part is the origin of the gluing redundancy itself. On the ordinary or central-defect branch, gauge-as-gluing is the finite-regulator overlap redundancy of local chart presentations: once the overlap algebras are realized in finite-dimensional charts, overlap-consistent rechartings of a cut form a compact unitary transition system. The collar theorem of Section 2.3 should therefore be read with its boundary group \(G_\Sigma\) understood as shorthand for this derived compact boundary action, not as a model-specific add-on. On the genuinely noncentral branch, the same weak gluing data are encoded by a compact crossed-module change system, so the fixed-cutoff collar theorem upgrades to a higher-gauge statement rather than failing. The fixed-cutoff topological package is therefore closed on all three branches.
A second structural point is UV underdetermination. If each patch Hilbert space is tensored with an inert finite ancillary factor and the observable patch algebras are embedded as \(\mathcal A(P)\otimes \mathbf 1\), then the physical observables, collar conditional mutual information, carried Markov errors, and quotient normal forms are unchanged. So the physical UV branch is fixed only modulo gauge or implementation hiding together with such ancillary stabilization, not a unique microscopic presentation; the unique theorem-grade object is the induced family of terminal expectation functionals on the declared physical observable algebras, not a preferred microscopic representative.

\subsubsection{Constructive gluing on tree covers}\label{31-constructive-gluing-on-tree-covers}

\textbf{Theorem 3.1 (tree gluing).} Let a rooted tree of patches satisfy a tree-ordered overlap structure and a tripartite factorization \((A_k, B_k, C_k)\) at step \(k\). If a target state \(\rho^{\ast}\) obeys \(I(A_k:C_k \mid B_k) \le \varepsilon_k\), then there exist recovery maps \(\mathcal{R}_k\) such that

\[
\| \rho^{\ast}_{A_k B_k C_k} - (\mathrm{id}_{A_k} \otimes \mathcal{R}_k)(\rho^{\ast}_{A_k B_k}) \|_1
\le \delta_k,
\]

with

\[
\delta_k = 2 \sqrt{\ln 2\, \varepsilon_k}.
\]

The iteratively glued state \(\hat{\rho}\) satisfies

\[
\| \hat{\rho} - \rho^{\ast} \|_1 \le \min\left(2, \sum_{k=2}^n \delta_k\right).
\]

\textbf{Proof.} Induct on \(k\). The recovery error contracts under CPTP maps, so the errors add. QED.

\subsubsection{Gauge-as-gluing and loops}\label{32-gauge-as-gluing-and-loops}

At finite regulator scale, the fixed-cutoff gauge-as-gluing package identifies the overlap-consistency redundancy of local chart presentations. Choose finite-dimensional local presentations of the overlap algebras on a connected cut \(\Sigma\). Because the overlap algebras are matrix algebras, any overlap-consistent change of chart is inner and is implemented by a unitary on the cut Hilbert space. Fixing a reference chart, the compact closure of the subgroup generated by all such recharting unitaries is the boundary gluing group \(K_\Sigma\); before fixing the reference chart, the same data form a compact unitary groupoid.

\textbf{Proposition 3.2a (Derived gauge-as-gluing at finite regulator).} Let a finite regulator chart be chosen for the patches meeting along a connected interface \(\Sigma\). Then overlap consistency determines a compact unitary transition system on the cut data. On the ordinary or central-defect branch, this transition system reduces to a compact boundary group \(K_\Sigma\), and when the triple-overlap defect is central its projective composition law lifts to a genuine action of a compact central extension \(\widehat K_\Sigma\). Gauge is therefore the overlap redundancy itself, not an extra primitive.

\textbf{Proof.} On a finite-dimensional matrix algebra every \(^*\)-automorphism is inner, so each overlap-consistent recharting is conjugation by a unitary on the cut Hilbert space. The subgroup generated by those unitaries has compact closure inside a finite-dimensional unitary group. If triple-overlap defects are central, the resulting projective composition law lifts to a central extension. QED.

\textbf{Lemma 3.2b (trees vs loops).} If the patch adjacency graph is a tree, one can choose local charts \(h_i\) so that the overlap labels satisfy \(g_{ij} = h_i^{-1} h_j\) on all edges. If loops exist, the loop holonomy

\[
H(\gamma) = g_{i_1 i_2} g_{i_2 i_3} \cdots g_{i_n i_1}
\]

is invariant under local frame changes. Nontrivial holonomy is the obstruction to global trivialization. QED.

\textbf{Corollary 3.2c (Collar consequence on the ordinary or central-defect branch).} Let \(B_\delta = B_L \cup B_R\) be a collar around a cap boundary \(\Sigma\), and set \(\widehat K_\Sigma = K_\Sigma\) on the ordinary branch. Then the EC theorem of Section 2.3 has the derived interpretation

\[
H_{B_\delta}
=
(\tilde{\mathcal H}_{B_L}\otimes \tilde{\mathcal H}_{B_R})^{\widehat K_\Sigma}
\cong
\bigoplus_\alpha
(H_{b_L^\alpha}\otimes H_{b_R^\alpha}),
\]

with

\[
Z(A_{\mathrm{EC}}(B_\delta)) = \bigoplus_\alpha \mathbb C\,\mathbf 1_\alpha.
\]

The right half-collar carries the contragredient action because it sees the inverse transport across the same cut. Exact Markov normal forms used later require the additional state hypothesis \(I_\omega(A_\delta:D_\delta \mid B_\delta)=0\) or the explicitly stated idealized recoverability limit; the block decomposition itself is kinematic and follows from the derived boundary action.

\textbf{Proof sketch.} Decompose the left boundary data into irreps \((V_\alpha \otimes H_{b_L^\alpha})\) of \(\widehat K_\Sigma\) and the right boundary data into the dual modules \((V_\beta^* \otimes H_{b_R^\beta})\). Then Schur's lemma leaves a singlet only when \(\alpha=\beta\), producing the displayed direct sum. QED.

\subsubsection{Loop obstruction class (central defect)}\label{33-loop-obstruction-class-central-defect}

On the central-defect subbranch, define central defects \(z_{ijk}\) by

\[
\varphi_{ij} \varphi_{jk} \varphi_{ki} = \mathrm{Ad}(z_{ijk}) \quad \text{on } \mathcal{A}_{ijk}.
\]

Then \(\{z_{ijk}\}\) is a \v{C}ech 2-cocycle, and its cohomology class \([z]\) is gauge invariant. Central defects do not obstruct the collar block decomposition: they only replace \(K_\Sigma\) by its central extension \(\widehat K_\Sigma\). Ordinary loop-coherent gluing exists iff \([z] = 0\). (A full proof appears in Section 6.4 below, in the algebra-net language.)

\subsubsection{Non-central obstruction (2-group cocycle)}\label{34-non-central-obstruction-2-group-cocycle}

When defects are not central, the natural coefficient data is a crossed module \((H \to G)\) with an action of \(G\) on \(H\) by conjugation. Here \(G\) is the reconstructed gauge group, and \(H\) is the unitary group acting on edge multiplicity spaces, with boundary map \(\partial: H \to G\).

A crossed module is a homomorphism \(\partial: H \to G\) together with an action of \(G\) on \(H\) such that

\[
\partial(g \triangleright h) = g\,\partial(h)\,g^{-1},
\qquad
\partial(h) \triangleright h' = h h' h^{-1}.
\]

On a good cover \(\{P_i\}\), a weakly coherent gluing is encoded by:

\[
g_{ij}: P_{ij} \to G,\qquad h_{ijk}: P_{ijk} \to H,
\]

obeying the 2-cocycle conditions

\[
g_{ij} g_{jk} = \partial(h_{ijk}) g_{ik},
\]

and on quadruple overlaps,

\[
h_{jkl} h_{ijl} = (g_{ij} \triangleright h_{ikl}) h_{ijk}.
\]

Gauge changes act by 1- and 2-cochains in the standard way for crossed-module cohomology.

\textbf{Theorem 3.4 (non-central obstruction).} Loop-coherent gluing exists iff the 2-cocycle \((g_{ij}, h_{ijk})\) is equivalent to the trivial cocycle in nonabelian \v{C}ech \(H^2\) with values in the crossed module \((H \to G)\).

\textbf{Proof sketch.} Strict gluing corresponds to \(h_{ijk}=1\) and \(g_{ij}g_{jk}=g_{ik}\). Gauge changes are exactly the crossed-module coboundaries, so strictification exists iff the 2-class is trivial. QED.

The central-defect case is the abelian truncation with \(H\) central and trivial action, which reduces to Section 3.3. A genuinely noncentral class is the point where the fixed-cutoff collar theorem upgrades from the ordinary-group package to the higher-gauge one below.

\textbf{Corollary 3.4a (Fixed-cutoff higher-gauge EC and transportability).} For a finite regulator chart on a connected cut \(\Sigma\), the genuinely noncentral branch admits a compact crossed-module change system
\[
\mathcal T_\Sigma
=
C^1(N_\Sigma,H_\Sigma)\rtimes C^0(N_\Sigma,G_\Sigma).
\]
The physical higher-gauge collar is
\[
\mathcal H_B^{2g}
=
(\widetilde{\mathcal H}_{B_L}\otimes \widetilde{\mathcal H}_{B_R})^{\mathcal T_\Sigma}
\cong
\bigoplus_\lambda
(\mathcal H_{b_L^\lambda}\otimes \mathcal H_{b_R^\lambda}),
\]
with
\[
Z(\mathcal A_{2g}(B))
=
\bigoplus_\lambda \mathbb C\,\mathbf 1_\lambda,
\]
and the defect class
\[
q_\Sigma=[(g,h)]\in \check H^2(N_\Sigma,H_\Sigma\to G_\Sigma)
\]
is invariant under local rechartings, classifies fixed-cutoff genuinely noncentral sectors, and vanishes iff the defect is removable.

\textbf{Proof sketch.} Pair-overlap rechartings are inner, while triple-overlap associators strictify to compact crossed-module data. Finite-dimensional unitary \(\mathcal T_\Sigma\)-modules decompose semisimply, and Schur matching leaves only diagonal left/right blocks. The transport statement is the crossed-module \v{C}ech analogue of the central-defect case. QED.

\begin{center}\rule{0.5\linewidth}{0.5pt}\end{center}
\subsection{Modular Flow and Lorentz Kinematics}\label{4-modular-flow-and-lorentz-kinematics}

\subsubsection{Modular additivity in the Markov collar limit}\label{41-modular-additivity-in-the-markov-collar-limit}

Consider a collar tripartition \(A:B:D\) around a cap boundary, with the EC decomposition of Section 2.3 understood. Define, for a faithful reference state \(\omega\),
\[
\Delta K(\omega)
:=
K_{ABD}(\omega)-K_{AB}(\omega)-K_{BD}(\omega)+K_B(\omega).
\]

Exact modular additivity is not a consequence of EC alone. It is available only on the exact Markov set, or along a controlled fixed-cutoff family that approaches that set on one fixed collar model. Three quantities must therefore be kept separate:
\begin{enumerate}
\item the raw collar conditional mutual information \(I(A:D\mid B)\);
\item the constructive Fawzi--Renner comparison error
\[
r_{\mathrm{FR}}(\varepsilon):=
2\sqrt{1-e^{-\varepsilon}}
\le
2\sqrt{\varepsilon};
\]
\item the fixed-collar exact-Markov replacement modulus
\[
\delta^{\mathrm M}_{A:B:D}(\varepsilon)
:=
\sup\left\{
\inf_{\sigma\in\mathfrak M_{A:B:D}}\|\rho-\sigma\|_1:
I(A:D\mid B)_\rho\le \varepsilon
\right\},
\]
where
\[
\mathfrak M_{A:B:D}
:=
\left\{
\sigma_{ABD}: I(A:D\mid B)_\sigma=0
\right\}.
\]
\end{enumerate}

\textbf{Lemma 4.1a (Exact Markov implies exact additivity up to a central term).} If \(I(A:D \mid B)_\omega = 0\), then the EC decomposition of Section 2.3 puts the state in HJPW block form
\[
\omega_{ABD}
=
\bigoplus_{\alpha}
p_{\alpha}\,
\omega^{(\alpha)}_{A b_L^{\alpha}}
\otimes
\omega^{(\alpha)}_{b_R^{\alpha} D},
\]
and \(\Delta K(\omega)\) is central. On the canonical HJPW block model one may take
\[
\Delta K(\omega)=0.
\]
Equivalently, there exists a central operator \(K_{\partial,ABD}(\omega)\in Z(\mathcal A(B))\) such that
\[
K_{ABD}(\omega)
=
K_{AB}(\omega)+K_{BD}(\omega)-K_B(\omega)+K_{\partial,ABD}(\omega).
\]

\textbf{Proof.} On each HJPW block the modular Hamiltonians of \(ABD\), \(AB\), \(BD\), and \(B\) are the logarithms of tensor-product states with the same classical block weight \(p_\alpha\). The blockwise logarithms therefore cancel exactly in the combination \(K_{ABD}-K_{AB}-K_{BD}+K_B\). If one keeps the blockwise endpoint-label bookkeeping explicit rather than absorbing it into the canonical block identification, the remainder is a direct sum of block constants, hence central. QED.

\textbf{Proposition 4.1b (Controlled exact-Markov replacement on a fixed collar).}
Because the state space is compact and conditional mutual information is continuous in finite dimension,
\[
\delta^{\mathrm M}_{A:B:D}(\varepsilon)\longrightarrow 0
\qquad
(\varepsilon\downarrow0).
\]
Thus small collar CMI converges to the exact HJPW normal form only in this controlled fixed-cutoff sense; the manuscript does not claim a dimension-free one-shot trace-norm bound from small \(I(A:D\mid B)\) directly to an exact Markov state.

\textbf{Corollary 4.1c (Carried collar defect operator).} Let \(\omega_\varepsilon\) satisfy \(I(A:D\mid B)_{\omega_\varepsilon}\le \varepsilon\), and choose \(\sigma_\varepsilon\in\mathfrak M_{A:B:D}\) with
\[
\|\omega_\varepsilon-\sigma_\varepsilon\|_1
\le
\delta^{\mathrm M}_{A:B:D}(\varepsilon).
\]
Define the carried defect operator
\[
\mathfrak D_{A:B:D}(\omega_\varepsilon,\sigma_\varepsilon)
:=
\Delta K(\omega_\varepsilon)-\Delta K(\sigma_\varepsilon).
\]
Then
\[
K_{ABD}(\omega_\varepsilon)
=
K_{AB}(\omega_\varepsilon)+K_{BD}(\omega_\varepsilon)-K_B(\omega_\varepsilon)
+K_{\partial,ABD}(\sigma_\varepsilon)
+\mathfrak D_{A:B:D}(\omega_\varepsilon,\sigma_\varepsilon),
\]
where \(K_{\partial,ABD}(\sigma_\varepsilon)=\Delta K(\sigma_\varepsilon)\) is central. For every bounded observable \(X\) supported on \(A\cup B\),
\[
\left|
\operatorname{Tr}\!\left[X(\omega_\varepsilon-\sigma_\varepsilon)\right]
\right|
\le
\|X\|_\infty\,\delta^{\mathrm M}_{A:B:D}(\varepsilon).
\]
If the relevant marginals of \(\omega_\varepsilon\) and \(\sigma_\varepsilon\) are uniformly faithful with lower spectral bound \(\lambda_\ast>0\), then
\[
\|\mathfrak D_{A:B:D}(\omega_\varepsilon,\sigma_\varepsilon)\|_\infty
\le
4\lambda_\ast^{-1}\,\delta^{\mathrm M}_{A:B:D}(\varepsilon).
\]

Independently, Fawzi--Renner recovery supplies a constructive comparison state \(\omega^{\mathrm{rec}}_\varepsilon\) with
\[
\|\omega_\varepsilon-\omega^{\mathrm{rec}}_\varepsilon\|_1
\le
r_{\mathrm{FR}}(\varepsilon),
\qquad
r_{\mathrm{FR}}(\varepsilon):=2\sqrt{1-e^{-\varepsilon}}\le 2\sqrt{\varepsilon}.
\]
This \(O(\varepsilon^{1/2})\) term is the finite-stage observable error. The modulus \(\delta^{\mathrm M}_{A:B:D}(\varepsilon)\) is instead the fixed-collar quantity that justifies replacing the physical state by an exact Markov normal form in the later geometric modular arguments.

Accordingly, the manuscript's later exact collar formulas take one of two forms:
\begin{enumerate}
\item literal exactness when the reference state is exact Markov on the relevant collar; or
\item a controlled collar family for which \(\delta^{\mathrm M}_{A:B:D}(\varepsilon_\delta)\to0\), with the constructive Fawzi--Renner remainder and the carried operator defect \(\mathfrak D_{A:B:D}\) kept explicit at finite stage.
\end{enumerate}
\subsubsection{\texorpdfstring{Theorem: \(\mathrm{BW}_{S^2}\) on the extracted prime geometric subnet}{4.2 Theorem: \textbackslash mathrm\{BW\}\_\{S\^{}2\} on the extracted prime geometric subnet}}\label{42-theorem-mathrmbw_s2-from-markov-locality-symmetry-regularity}

The collar analysis proves a fixed-cutoff statement on the finite type-I regulator net: the reduced cap state has a literal density matrix, its modular Hamiltonian exists, and its nonadditive part is confined to a shrinking collar up to carried errors. The Lorentz claim is therefore not a fixed-cutoff matrix-algebra identity. Its target is the refinement-limit geometric cap pair \((\mathcal A_\infty^{\mathrm{geo}}(C),\omega_\infty^{\mathrm{geo},C})\) singled out by the controlled scaling clause \textbf{T2} together with the theorem-local geometric cap-pair extraction. Nothing in the local MaxEnt branch by itself proves that the scaling-limit cap algebra remains type I, and nothing in the fixed-cutoff MaxEnt/collar package by itself emits that prime geometric cap pair. Ordered cut-pair rigidity is then what turns the extracted pair into the BW statement.

Fix a cap \(C\subset S^2\) and a shrinking collar family \((A_\delta,B_\delta,D_\delta)\) around \(\partial C\). Write
\[
\varepsilon_\delta:=I(A_\delta:D_\delta\mid B_\delta)_\omega,
\qquad
r_{\mathrm{FR}}(\varepsilon_\delta):=
2\sqrt{1-e^{-\varepsilon_\delta}}
\le 2\sqrt{\varepsilon_\delta},
\]
and, on each fixed faithful collar model,
\[
\eta_\delta^{\mathrm M}
:=
4\lambda_\ast^{-1}\,
\delta^{\mathrm M}_{A_\delta:B_\delta:D_\delta}(\varepsilon_\delta),
\]
where \(\lambda_\ast>0\) is the lower spectral bound used to compare modular Hamiltonians. All exact collar formulas in this subsection are therefore to be read either literally at exact Markovity or along a controlled collar family satisfying
\[
\delta^{\mathrm M}_{A_\delta:B_\delta:D_\delta}(\varepsilon_\delta)\to 0
\qquad(\delta\downarrow 0),
\]
with \(r_{\mathrm{FR}}(\varepsilon_\delta)\) and \(\eta_\delta^{\mathrm M}\) carried explicitly.

For each cap \(C\), let \(\lambda_C(s)\subset \mathrm{Conf}^+(S^2)\) denote the standard cap-preserving conformal one-parameter subgroup, normalized so that the null blow-up near a smooth cut acts by \(v\mapsto e^{-s}v\), and let \(\alpha_{\lambda_C(s)}\) denote the induced automorphism of the scaling-limit cap net.

\textbf{Theorem 4.2 (\(\mathrm{BW}_{S^2}\) on the extracted prime geometric subnet).} Let \(C\subset S^2\) be a round cap. Assume the OPH axioms, the derived fixed-cutoff regulator/collar package established earlier, the controlled scaling-limit scope clause \textbf{T2}, and the theorem-local UV inputs that:
\begin{enumerate}
\item the transported cap marginals on the geometric subnet admit a realized scaling-limit geometric cap pair \((\mathcal A_\infty^{\mathrm{geo}}(C),\omega_\infty^{\mathrm{geo},C})\) through the carried-collar local weak-\(*\) / GNS extraction package; and
\item ordered cut-pair rigidity holds on that extracted prime geometric pair.
\end{enumerate}
Let \(\lambda_C(s)\) denote the standard cap-preserving conformal one-parameter subgroup, normalized so that the null blow-up near a smooth cut acts by \(v\mapsto e^{-s}v\). Then the scaling-limit modular automorphism group is
\[
\sigma_t^{\omega_\infty^{\mathrm{geo},C}}
=
\alpha_{\lambda_C(2\pi t)}.
\]
At finite regulator stage the nonadditive part of
\[
K_C^{(\delta)}:=-\log \rho_C^{(\delta)}
\]
is confined to the shrinking collar up to the carried errors \(r_{\mathrm{FR}}(\varepsilon_\delta)\) and \(\eta_\delta^{\mathrm M}\). No separate cap-isotropy/\(\mathrm{SO}(2)\)-equivariance selector or Euclidean-regularity premise is used. If the scaling-limit cap algebra remains type I, the same automorphism identity may be represented by
\[
K_C=2\pi B_C;
\]
otherwise the automorphism identity is the full statement and the action is outer. Consequently, for every controlled collar family with \(\delta^{\mathrm M}_{A_\delta:B_\delta:D_\delta}(\varepsilon_\delta)\to0\) and surviving faithfulness bound, replacing the finite-stage modular action by the geometric cap-dilation action incurs only the carried errors \(r_{\mathrm{FR}}(\varepsilon_\delta)\) and \(\eta_\delta^{\mathrm M}\), and these vanish in the refinement limit.

\textbf{Proof.} Markov locality localizes the modular defect to the shrinking collar, with carried finite-stage errors \(r_{\mathrm{FR}}(\varepsilon_\delta)\) and \(\eta_\delta^{\mathrm M}\). By \textbf{T2}, the theorem concerns only scaling-limit cap pairs. By the theorem-local extraction hypothesis, the relevant limit object is the prime geometric cap pair. Ordered cut-pair rigidity identifies the scaling-limit modular automorphism group with the standard cap-preserving conformal flow up to normalization:
\[
\sigma_t^{\omega_\infty^{\mathrm{geo},C}}
=
\alpha_{\lambda_C(\kappa_C t)}
\]
for some \(\kappa_C>0\). Because \(s\) is normalized by the null blow-up \(v\mapsto e^{-s}v\), the modular KMS condition fixes \(\kappa_C=2\pi\). Thus
\[
\sigma_t^{\omega_\infty^C}
=
\alpha_{\lambda_C(2\pi t)}.
\]
If the limit algebra remains type I, this automorphism identity is represented by \(K_C=2\pi B_C\); otherwise the automorphism identity is the full statement and the action is outer. QED.

\textbf{Scope note.} The theorem fixes the \(2\pi\) normalization and the carried refinement limit on the extracted geometric cap pair. It does \emph{not} prove that MaxEnt alone emits that pair. The remaining UV-side scaffold is exactly a realized transported geometric cap-local system whose carried-collar schedule
\[
r_{\mathrm{FR}}(\varepsilon_{n,m,\delta})
+
4\bar\lambda_{m,\delta}^{-1}
\delta^{\mathrm{M}}_{A_{n,m,\delta}:B_{n,m,\delta}:D_{n,m,\delta}}
\bigl(\varepsilon_{n,m,\delta}\bigr)
\]
is derived from transported CMI together with an eventual fixed-local-collar modular-transport common floor, followed by ordered cut-pair rigidity on the emitted scaling-limit geometric cap pair. No separate cap-isotropy or Euclidean-regularity premise is used, and no second comparison-state faithfulness clause is missing.
\subsubsection{\texorpdfstring{Theorem: \(\mathrm{BW}_{S^2}\) implies Lorentz kinematics}{4.3 Theorem: \textbackslash mathrm\{BW\}\_\{S\^{}2\} implies Lorentz kinematics}}\label{43-theorem-mathrmbw_s2-implies-lorentz-kinematics}

\textbf{Theorem 4.3 (Lorentz kinematics on the screen).} Under the hypotheses of Theorem~\ref{42-theorem-mathrmbw_s2-from-markov-locality-symmetry-regularity},
\[
\mathrm{Conf}^+(S^2)\cong \mathrm{PSL}(2,\mathbb C)\cong \mathrm{SO}^+(3,1).
\]
The cap modular flows are therefore the standard one-parameter Lorentz boost/dilation subgroups in the celestial-sphere realization, and the induced local kinematic group is the connected Lorentz group.

\textbf{Proof.} Orientation-preserving conformal maps of \(S^2\) are exactly the M\"obius transformations, so
\[
\mathrm{Conf}^+(S^2)\cong \mathrm{PSL}(2,\mathbb C)\cong \mathrm{SO}^+(3,1).
\]
By Theorem~\ref{42-theorem-mathrmbw_s2-from-markov-locality-symmetry-regularity}, on the extracted prime geometric subnet the realized scaling-limit cap modular automorphism groups are the standard cap-preserving conformal dilations with the \(2\pi\) normalization fixed internally. Hence the local kinematics induced by modular flow is the connected Lorentz group. QED.
\subsection{Gravity from Entanglement Equilibrium}\label{5-gravity-from-entanglement-equilibrium}

\subsubsection{Cap first law}\label{51-cap-first-law}

Section 4.2 fixes the cap modular statement first at the automorphism level. On the extracted prime geometric subnet, the scaling-limit cap modular flow is geometric with the standard \(2\pi\) normalization,
\[
\sigma_t^{\omega_\infty^C}
=
\alpha_{\lambda_C(2\pi t)}.
\]
If the realized scaling-limit cap algebra happens to remain type I, one may choose a density matrix \(\rho_C^\omega\) and modular Hamiltonian
\[
K_C:=-\log\rho_C^\omega,
\]
and then for a perturbation \(\rho(\varepsilon)\) with \(\rho(0)=\omega\) the first-law identity reads
\[
\delta S_C=\delta\langle K_C\rangle
=
2\pi\,\delta\langle B_C\rangle.
\]
In the generic continuum/non-type-I case, Section 4.2 does not supply an inner operator identity \(K_C=2\pi B_C\); the theorem surface there is only the geometric modular automorphism statement, so the cap first-law formula is conditional on the special type-I realization.

\subsubsection{Null-surface modular bridge}\label{52-null-surface-modular-bridge}

This subsection extends the fixed-cutoff weak-tail-generator boundary to the derived positive null-translation stage and the exact half-line generator/charge identification. At fixed cutoff one can transfer cut-center data to regulated null strips, obtain exact or controlled strip additivity on one inherited strip model, and derive a weak tail generator for the renormalized half-line family. On the OPH geometric scaling branch of Theorem 4.2, the null half-line blow-up net then inherits geometric dilation and therefore half-sided modular inclusion, so Borchers--Wiesbrock supplies an explicit positive null-translation generator on its Stone domain. The remaining downstream boundary is narrower: bounded-interval formulas still require the separate interval-preserving projective branch, and the later tensor upgrade still carries the null-invisible metric ambiguity. Accordingly the c.12-local bridge consists of:
\begin{itemize}
\tightlist
\item
  null-cut center transfer and the inherited split condition;
\item
  exact or controlled four-term strip additivity on one fixed inherited strip model;
\item
  endpoint-Lipschitz control for the renormalized half-line family and the resulting weak tail generator;
\item
  derived half-sided modular inclusion on the null half-line blow-up net, and the resulting Borchers positive translation generator.
\end{itemize}
The later density formula and \(T_{kk}\) bridge are kept below only as explicit downstream templates rather than as fixed-cutoff theorems proved in this subsection. Lemma 5.2f makes the positive null-translation generator itself explicit: it records the Borchers unitary group, positivity and self-adjointness on the Stone domain, the corrected affine commutator relation, and the half-line modular-Hamiltonian identity \(K_a=K_0-2\pi aP_\Omega\). Bounded-interval formulas remain downstream.

\textbf{Proposition 5.2a (Null-cut center transfer and inherited split).} Fix a regulated null tripartition
\[
I_-=(v_1,v_2),\qquad
J=(v_2,v_3),\qquad
I_+=(v_3,v_4),
\]
with cuts \(\Gamma_-:=\{v=v_2\}\) and \(\Gamma_+:=\{v=v_3\}\). Under the fixed-cutoff realized presentation, assume that the two null cuts inherit the ordinary or central-defect boundary-redundancy data used in the spatial collar branch, so that in a compatible type-I regulator presentation one has
\[
\tilde{\mathcal H}_{I_-}
\cong
\bigoplus_{\alpha_-}
W_{\alpha_-}\otimes \mathcal H_{i_-^{\alpha_-}},
\]
\[
\tilde{\mathcal H}_{J}
\cong
\bigoplus_{\alpha_-,\alpha_+}
W_{\alpha_-}^{*}\otimes
\mathcal H_{j^{\alpha_-,\alpha_+}}\otimes
W_{\alpha_+},
\]
\[
\tilde{\mathcal H}_{I_+}
\cong
\bigoplus_{\alpha_+}
W_{\alpha_+}^{*}\otimes \mathcal H_{i_+^{\alpha_+}},
\]
where opposite sides of each cut carry inverse transport. If the strip algebra is the commutant of the transported cut actions, then
\[
\mathcal A(J)
\cong
\bigoplus_{\alpha_-,\alpha_+}
\mathcal B(\mathcal H_{j^{\alpha_-,\alpha_+}}),
\qquad
Z(\mathcal A(J))
=
\bigoplus_{\alpha_-,\alpha_+}\mathbb C\,P_{\alpha_-,\alpha_+}.
\]
If, in addition, each multiplicity space factors as
\[
\mathcal H_{j^{\alpha_-,\alpha_+}}
\cong
\mathcal H_{j_L^{\alpha_-,\alpha_+}}
\otimes
\mathcal H_{j_R^{\alpha_-,\alpha_+}},
\]
with \(\mathcal A(I_-\cup J)\) acting blockwise only on
\(\mathcal H_{i_-^{\alpha_-}}\otimes \mathcal H_{j_L^{\alpha_-,\alpha_+}}\)
and \(\mathcal A(J\cup I_+)\) acting blockwise only on
\(\mathcal H_{j_R^{\alpha_-,\alpha_+}}\otimes \mathcal H_{i_+^{\alpha_+}}\),
then
\[
\mathcal A(J)
\cong
\bigoplus_{\alpha_-,\alpha_+}
\mathcal B(\mathcal H_{j_L^{\alpha_-,\alpha_+}})
\otimes
\mathcal B(\mathcal H_{j_R^{\alpha_-,\alpha_+}}),
\]
and
\[
\mathcal H_{I_-\cup J\cup I_+}
\cong
\bigoplus_{\alpha_-,\alpha_+}
\mathcal H_{i_-^{\alpha_-}}
\otimes
\mathcal H_{j_L^{\alpha_-,\alpha_+}}
\otimes
\mathcal H_{j_R^{\alpha_-,\alpha_+}}
\otimes
\mathcal H_{i_+^{\alpha_+}}.
\]

\textbf{Proof.} Complete reducibility gives the displayed decompositions. Commuting with the two transported cut actions forces strip operators to preserve the sector pair \((\alpha_-,\alpha_+)\) and act only on the multiplicity space; Schur's lemma kills the off-diagonal intertwiners and leaves the direct-sum algebra above. Taking invariants across both cuts in the glued tripartition again uses Schur's lemma and leaves exactly the matching sector pairs. The left/right split of the multiplicity spaces is additional structure; it is not forced by the center transfer alone. QED.

\textbf{Corollary 5.2b (Exact or controlled four-term null modular relation on an inherited strip model).} On one fixed finite-dimensional strip model satisfying Proposition 5.2a, define
\[
\mathfrak M_{I_-:J:I_+}
:=
\left\{
\sigma:\ I(I_-:I_+\mid J)_\sigma=0
\right\},
\]
and
\[
\delta^{\mathrm M}_{I_-:J:I_+}(\varepsilon)
:=
\sup\left\{
\inf_{\sigma\in\mathfrak M_{I_-:J:I_+}}\|\rho-\sigma\|_1:
I(I_-:I_+\mid J)_\rho\le\varepsilon
\right\}.
\]
For any strip state \(\eta\), let
\[
\Delta K_J(\eta)
:=
K_{I_-\cup J\cup I_+}(\eta)-K_{I_-\cup J}(\eta)-K_{J\cup I_+}(\eta)+K_J(\eta).
\]

If the strip reference state \(\omega\) is exact Markov, then \(\Delta K_J(\omega)\) is central, and on the canonical inherited HJPW model one may take
\[
\Delta K_J(\omega)=0.
\]
Equivalently, there exists a central operator \(K_{\partial,J}(\omega)\in Z(\mathcal A(J))\) such that
\[
K_{I_-\cup J\cup I_+}(\omega)
=
K_{I_-\cup J}(\omega)+K_{J\cup I_+}(\omega)-K_J(\omega)+K_{\partial,J}(\omega).
\]

If instead \(I(I_-:I_+\mid J)_\omega\le\varepsilon\), choose \(\widetilde\omega_J\in\mathfrak M_{I_-:J:I_+}\) with
\[
\|\omega-\widetilde\omega_J\|_1
\le
\delta^{\mathrm M}_{I_-:J:I_+}(\varepsilon).
\]
Define
\[
\mathfrak D_J(\omega,\widetilde\omega_J)
:=
\Delta K_J(\omega)-\Delta K_J(\widetilde\omega_J).
\]
Then
\[
K_{I_-\cup J\cup I_+}(\omega)
=
K_{I_-\cup J}(\omega)+K_{J\cup I_+}(\omega)-K_J(\omega)
+K_{\partial,J}(\widetilde\omega_J)
+\mathfrak D_J(\omega,\widetilde\omega_J),
\]
with \(K_{\partial,J}(\widetilde\omega_J)=\Delta K_J(\widetilde\omega_J)\in Z(\mathcal A(J))\). Every bounded observable on \(I_-\cup J\) or \(J\cup I_+\) then differs from the exact-Markov reference by at most
\[
\|O\|_\infty\,\delta^{\mathrm M}_{I_-:J:I_+}(\varepsilon).
\]
If the relevant strip marginals are uniformly faithful with lower spectral bound \(\lambda_\ast>0\), then
\[
\|\mathfrak D_J(\omega,\widetilde\omega_J)\|_\infty
\le
4\lambda_\ast^{-1}\,\delta^{\mathrm M}_{I_-:J:I_+}(\varepsilon).
\]
Independently, Fawzi--Renner gives a recovered comparison state with trace-norm error
\[
r_{\mathrm{FR}}(\varepsilon)=2\sqrt{1-e^{-\varepsilon}}\le 2\sqrt{\varepsilon}.
\]
Thus the exact four-term strip relation is available only at exact Markovity or in a controlled strip family on one fixed inherited strip model with \(\delta^{\mathrm M}_{I_-:J:I_+}(\varepsilon_J)\to0\), while \(\mathfrak D_J\) is carried explicitly at finite stage.

\textbf{Proof.} On the inherited split of Proposition 5.2a, the exact-Markov case is the same HJPW block calculation as in the spatial collar theorem: the four logarithms cancel blockwise on the canonical model, and any residual bookkeeping term is central. The controlled replacement is the same fixed-model compactness argument used for ordinary collars, after relabeling \(A,B,D\mapsto I_-,J,I_+\); the operator-norm bound on \(\mathfrak D_J\) is the corresponding relabeling of the modular-transport estimate. QED.

\textbf{Definition (Renormalized null modular functional).} For a null interval \(I\) on generator \(\Omega\), let \(K_\partial(I,\Omega)\) denote the central endpoint-label term singled out by Corollary 5.2b on the inherited strip model, or by its controlled exact-Markov replacement when that model is used as reference. Define
\[
\widetilde K[I,\Omega]
:=
K[I,\Omega]-K_\partial(I,\Omega),
\qquad
\widetilde K_a(\Omega):=\widetilde K[(a,\infty),\Omega].
\]

\textbf{Proposition 5.2c (Endpoint-Lipschitz null modular families and weak tail generator).} For renormalized modular Hamiltonians on one null generator,
\[
\widetilde K[I,\Omega]
:=
K[I,\Omega]-K_\partial(I,\Omega),
\qquad
\widetilde K_a(\Omega):=\widetilde K[(a,\infty),\Omega],
\]
the branch-internal endpoint-control coming from the local finite-constraint MaxEnt branch gives:
\[
\bigl|
\langle\psi,(\widetilde K[(a',b'),\Omega]-\widetilde K[(a,b),\Omega])\phi\rangle
\bigr|
\le
C_{\psi,\phi,\Omega}\bigl(|a'-a|+|b'-b|\bigr)
\]
for bounded intervals in a compact endpoint window, and
\[
\bigl|
\langle\psi,(\widetilde K_{a'}(\Omega)-\widetilde K_a(\Omega))\phi\rangle
\bigr|
\le
C_{\psi,\phi,\Omega}|a'-a|
\]
for half-lines. Hence
\[
f_{\psi,\phi}(a):=\langle\psi,\widetilde K_a(\Omega)\phi\rangle
\]
is locally Lipschitz and therefore absolutely continuous, and its distributional derivative defines the weak tail generator
\[
\langle\psi,q(a,\Omega)\phi\rangle
:=
-\frac{1}{2\pi}\,\partial_a f_{\psi,\phi}(a).
\]
For \(a<b\),
\[
\langle\psi,(\widetilde K_b(\Omega)-\widetilde K_a(\Omega))\phi\rangle
=
-2\pi\int_a^b \langle\psi,q(v,\Omega)\phi\rangle\,dv.
\]
At this stage \(q(a,\Omega)\) is a weak tail generator, not yet a local operator-valued density; its identification with the positive self-adjoint Borchers generator occurs only after Corollary 5.2e and Lemma 5.2f.

\textbf{Proof.} The derived endpoint-control estimate applies to the renormalized family after removal of the central endpoint term. For bounded intervals, the symmetric-difference length is bounded by \(|a'-a|+|b'-b|\); for half-lines it is exactly \(|a'-a|\). Therefore the matrix-element functions are Lipschitz in the endpoints. A Lipschitz function on a compact interval is absolutely continuous, so \(f_{\psi,\phi}(a)\) has an \(L^\infty_{\mathrm{loc}}\) derivative and obeys the fundamental theorem of calculus. Defining \(q\) as the rescaled negative derivative gives the displayed integral relation. QED.

\textbf{Downstream c.12--c.13 boundary.} The fixed-cutoff bridge established above is the end of this subsection's proved content.

\textbf{Lemma 5.2d (Downstream density-upgrade template).} If, in addition, the weak tail generator \(q(a,\Omega)\) is weakly differentiable in \(a\) and both \(f_{\psi,\phi}(a)\) and \(q_{\psi,\phi}(a)\) vanish at \(+\infty\), then one may define a density
\[
\langle\psi,p(a,\Omega)\phi\rangle
:=
-\partial_a\langle\psi,q(a,\Omega)\phi\rangle
=
\frac{1}{2\pi}\,\partial_a^2 \langle\psi,\widetilde K_a(\Omega)\phi\rangle
\]
and obtain the half-line formula
\[
\langle\psi,\widetilde K_a(\Omega)\phi\rangle
=
2\pi\int_a^\infty (v-a)\,\langle\psi,p(v,\Omega)\phi\rangle\,dv.
\]
This is a downstream density-upgrade step and is not proved from the fixed-cutoff strip arguments above.

\textbf{Null half-line blow-up net.} Fix a smooth cut point and choose affine coordinate \(v\) on generator \(\Omega\) so that the cut sits at \(v=0\). For \(a\ge 0\), write
\[
H_a:=(a,\infty),
\qquad
\mathcal M_a(\Omega)
:=
\overline{\bigvee_{a<c<d<\infty}\mathcal A((c,d),\Omega)}.
\]
By isotony of the null interval net,
\[
a\le b
\quad\Longrightarrow\quad
\mathcal M_b(\Omega)\subseteq \mathcal M_a(\Omega).
\]

\textbf{Corollary 5.2e (Derived half-sided modular inclusion on null half-lines).} On the OPH geometric scaling branch of Theorem 4.2, the blow-up modular action near a smooth entangling cut acts on the null coordinate by
\[
v\mapsto e^{-2\pi t}v.
\]
Therefore, for every \(a\ge 0\),
\[
\sigma_t^\omega\bigl(\mathcal M_a(\Omega)\bigr)
=
\mathcal M_{e^{-2\pi t}a}(\Omega).
\]
Hence, for every \(a>0\),
\[
\sigma_t^\omega\bigl(\mathcal M_a(\Omega)\bigr)\subseteq \mathcal M_a(\Omega)
\qquad (t\le 0),
\]
so the inclusion \(\mathcal M_a(\Omega)\subset \mathcal M_0(\Omega)\) is half-sided modular. After the harmless convention change \(t\mapsto -t\), this is the standard positive-time half-sided-inclusion form. The half-sided modular inclusion is therefore derived from OPH null-interval structure, isotony, and the scaling-limit geometric action rather than imported separately.

\textbf{Proof.} Blow up the cap modular flow of Theorem 4.2 near a smooth cut and restrict to the chosen null generator. In the tangent limit the cap-preserving flow becomes the null dilation \(v\mapsto e^{-2\pi t}v\). For every bounded interval \((c,d)\subset H_a\), this sends \(\mathcal A((c,d),\Omega)\) to \(\mathcal A((e^{-2\pi t}c,e^{-2\pi t}d),\Omega)\), and taking the von Neumann closure of the interval net yields the displayed identity for \(\mathcal M_a(\Omega)\). If \(t\le 0\), then \(e^{-2\pi t}a\ge a\), so \(H_{e^{-2\pi t}a}\subseteq H_a\), and isotony gives \(\mathcal M_{e^{-2\pi t}a}(\Omega)\subseteq \mathcal M_a(\Omega)\). QED.

\textbf{Lemma 5.2f (Positive null-translation generator).} For the derived half-sided modular inclusion
\[
\mathcal M_a(\Omega)\subset \mathcal M_0(\Omega)
\qquad (a>0),
\]
let \(\Delta_0(\Omega)\) be the modular operator of the standard pair \((\mathcal M_0(\Omega),\omega)\) and define \(K_0(\Omega):=-\log \Delta_0(\Omega)\). Borchers--Wiesbrock then yields a unique strongly continuous one-parameter unitary group
\[
U_\Omega(a)=e^{iaP_\Omega},
\qquad a\in\mathbb R,
\]
such that
\[
U_\Omega(a)\omega=\omega,
\qquad
U_\Omega(a)\mathcal M_b(\Omega)U_\Omega(a)^*=\mathcal M_{a+b}(\Omega)
\quad (a,b\ge0),
\]
and whose generator \(P_\Omega\) is positive and self-adjoint on the Stone domain
\[
D(P_\Omega)
:=
\left\{
\psi\in\mathcal H:
\lim_{a\to0}\frac{U_\Omega(a)\psi-\psi}{ia}\ \text{exists}
\right\}.
\]
Moreover,
\[
\Delta_0(\Omega)^{it}U_\Omega(a)\Delta_0(\Omega)^{-it}
=
U_\Omega(e^{-2\pi t}a),
\qquad
\Delta_0(\Omega)^{it}P_\Omega\Delta_0(\Omega)^{-it}
=
e^{-2\pi t}P_\Omega,
\]
and on the common invariant analytic core of \(K_0(\Omega)\) and \(P_\Omega\),
\[
[K_0(\Omega),P_\Omega]=-\,i\,2\pi P_\Omega.
\]
If \(\Delta_a(\Omega)\) denotes the modular operator of \((\mathcal M_a(\Omega),\omega)\) and \(K_a(\Omega):=-\log \Delta_a(\Omega)\), then
\[
\Delta_a(\Omega)=U_\Omega(a)\Delta_0(\Omega)U_\Omega(a)^*,
\qquad
K_a(\Omega)=U_\Omega(a)K_0(\Omega)U_\Omega(a)^*,
\]
so
\[
K_a(\Omega)=K_0(\Omega)-2\pi a\,P_\Omega
\]
as a quadratic-form identity on \(D(K_0(\Omega))\cap D(P_\Omega)\). Equivalently,
\[
\langle\psi,(K_b(\Omega)-K_a(\Omega))\phi\rangle
=
-2\pi(b-a)\langle\psi,P_\Omega\phi\rangle
\]
for all \(\psi,\phi\in D(K_0(\Omega))\cap D(P_\Omega)\). In the canonical normalization of Corollary 5.2b, the weak endpoint derivative from Proposition 5.2c is therefore exactly the Borchers generator. Bounded-interval modular-Hamiltonian formulas remain downstream and require the separate interval-preserving branch recorded in Theorem 5.2g.

\textbf{Proof.} Corollary 5.2e gives the required half-sided modular inclusion. Borchers--Wiesbrock then yields the unique strongly continuous unitary group \(U_\Omega(a)\), its positive self-adjoint generator \(P_\Omega\), and the affine covariance relation with the modular group; Stone's theorem gives the displayed domain formula. Differentiating
\[
\Delta_0(\Omega)^{it}U_\Omega(a)\Delta_0(\Omega)^{-it}
=
U_\Omega(e^{-2\pi t}a)
\]
first at \(a=0\) and then at \(t=0\) on the common analytic core gives
\[
[K_0(\Omega),P_\Omega]=-\,i\,2\pi P_\Omega.
\]
Because \(U_\Omega(a)\omega=\omega\) and \(U_\Omega(a)\mathcal M_0(\Omega)U_\Omega(a)^*=\mathcal M_a(\Omega)\), uniqueness of modular data for the translated standard pair gives the displayed formulas for \(\Delta_a(\Omega)\) and \(K_a(\Omega)\). Differentiating \(K_a(\Omega)=U_\Omega(a)K_0(\Omega)U_\Omega(a)^*\) in \(a\) yields \(dK_a/da=-2\pi P_\Omega\) as a quadratic-form identity, and integrating from \(0\) to \(a\) gives the stated half-line modular-Hamiltonian relation. QED.

\textbf{Theorem 5.2g (The OPH half-line generator is the local null-stress charge).} In the canonical normalization of the preceding half-line construction,
\[
\langle\psi,P_\Omega\phi\rangle
=
-\frac{1}{2\pi}\frac{d}{da}\Big|_{a=0^+}
\langle\psi,\widetilde K_a(\Omega)\phi\rangle
\]
for all \(\psi,\phi\in D(K_0(\Omega))\cap D(P_\Omega)\). In the effective spacetime description of that same renormalized half-line family, the local null-stress charge is defined by the identical endpoint derivative,
\[
Q_{kk}(\Omega)
:=
-\frac{1}{2\pi}\frac{d}{da}\Big|_{a=0^+}\widetilde K^{\mathrm{eff}}_a(\Omega),
\]
so
\[
P_\Omega = Q_{kk}(\Omega)
\]
as a quadratic-form identity on the common domain. The only imported continuum input is the standard local null modular-Hamiltonian formula for that same half-line family. Bounded-interval transport through the affine-covariant kernel \(g_I(v)\) remains a separate projective-branch step, and reconstruction of a full tensor from the directional charges remains subject to the null-invisible metric ambiguity.

\textbf{Boundary of the conditional null bridge.} The null bridge is conditional rather than automatic. Its explicit burden is:
\begin{itemize}
\tightlist
\item
  transferred cut-center data alone give the central sector decomposition of the null strip, not yet the left/right HJPW split;
\item
  the extra decomposition-inheritance condition of Proposition 5.2a is exactly what upgrades those sectors to the same collar-type block structure used in the spatial branch;
\item
  the fixed-cutoff null-strip bridge is exact for exact Markov strip states on that inherited decomposition, and otherwise is replaced by a controlled limit with the carried defect operator \(\mathfrak D_J\);
\item
  the local finite-constraint MaxEnt branch gives endpoint-Lipschitz control strong enough to define a weak tail generator for renormalized half-lines, and the OPH geometric scaling branch then derives half-sided modular inclusion on the half-line blow-up net;
\item
  Borchers--Wiesbrock therefore supplies a genuine positive null-translation generator \(P_\Omega\) inside the bridge itself, together with its Stone domain and the half-line modular-Hamiltonian relation \(K_a=K_0-2\pi aP_\Omega\);
\item
  the half-line generator/charge identification is fixed inside the bridge itself, while bounded-interval transport and the later tensor upgrade remain downstream.
\end{itemize}
This is the route from null-strip factorization to the later Jacobson branch, with the half-line generator/charge identification proved inside the null modular bridge itself.
\subsubsection{Modular energy as stress-tensor charge (UV CFT)}\label{53-modular-energy-as-stress-tensor-charge-uv-cft}

The effective-theory input here is no longer a second operator added on top of the null bridge. In the canonical normalization of the preceding null-half-line construction, the positive OPH null generator is the endpoint derivative of the renormalized half-line modular family:
\[
\langle\psi,P_\Omega\phi\rangle
=
-\frac{1}{2\pi}\frac{d}{da}\Big|_{a=0^+}
\langle\psi,\widetilde K_a(\Omega)\phi\rangle
\]
for all \(\psi,\phi\in D(K_0(\Omega))\cap D(P_\Omega)\).

Now pass to the effective spacetime description of that same half-line family. There the local null-stress charge is defined by the identical endpoint derivative,
\[
Q_{kk}(\Omega)
:=
-\frac{1}{2\pi}\frac{d}{da}\Big|_{a=0^+}\widetilde K^{\mathrm{eff}}_a(\Omega),
\]
so on the common quadratic-form domain
\[
P_\Omega = Q_{kk}(\Omega).
\]

Thus the OPH null generator is exactly the local null-stress charge in the effective spacetime description; this step is not a separate scaling-limit premise. When the effective description admits a local stress tensor, for example in a UV CFT regime or any local Lorentzian regime with local modular Hamiltonians, the same operator is the standard null-stress charge associated with the chosen null generator.

For comparison, in a CFT vacuum on a ball one has the familiar local formula
\[
H_{\zeta}=\int_{\Sigma}T_{ab}\zeta^b\,d\Sigma^a,
\]
so the present identification is the null-half-line version of that same modular-energy locality rather than an additional EFT postulate.

What still lies downstream is narrower than before. The separate interval-preserving projective branch is still needed to transport this half-line statement to bounded null intervals with affine-covariant kernel \(g_I(v)\), and the null-to-tensor upgrade still determines \(T_{ab}\) only up to the familiar ambiguity
\[
T_{ab}\mapsto T_{ab}+\phi g_{ab}.
\]
But the generator/charge identification itself is fixed inside the OPH null bridge.

\subsubsection{Localized generalized entropy from Markov + MaxEnt}\label{54-localized-generalized-entropy-from-markov--maxent}

Using the collar decomposition and the collar double-scaling hypothesis, supplied either as a stated collar-limit premise or derived via Theorem 2.5 from the local-Gibbs form together with exponential mixing, the state takes the Markov normal form. MaxEnt selection maximizes entropy within each edge sector, producing

\[
\rho_C = \bigoplus_{\alpha} p_{\alpha}
\left(\rho_{\mathrm{bulk},C}^{\alpha} \otimes \frac{\mathbf 1_{\mathrm{edge}}^{\alpha}}{d_{\alpha}}\right).
\]

The entropy splits as

\[
S(\rho_C) = H(p_{\alpha}) + \sum_{\alpha} p_{\alpha} S(\rho_{\mathrm{bulk},C}^{\alpha}) + \sum_{\alpha} p_{\alpha} \log d_{\alpha}.
\]

\emph{Convention:} Throughout this paper, "log" denotes the natural logarithm (ln), so entropies are measured in \textbf{nats} (1 nat = 1/ln 2 \ensuremath{\approx} 1.443 bits). This is standard in thermodynamics and QFT; the Bekenstein-Hawking formula S = A/4G uses nats. When clarity requires it, we write log\textsubscript{2} explicitly for bits.

Define

\[
S_{\mathrm{bulk}}(C) := H(p_{\alpha}) + \sum_{\alpha} p_{\alpha} S(\rho_{\mathrm{bulk},C}^{\alpha}),
\]

and the central area operator

\[
L_C := \sum_{\alpha} (\log d_{\alpha}) P_{\alpha}.
\]

Then

\[
S_{\mathrm{gen}}(C) := \mathrm{Tr}(\rho L_C) + S_{\mathrm{bulk}}(C).
\]

\textbf{Deriving Newton\textquotesingle s constant from edge entropy density.}

Rather than normalize \(L_C\) by fiat, we \emph{derive} the relation to \(G\) from the UV edge structure. In the collar double-scaling limit, the edge contribution becomes extensive along the entangling surface \(\Sigma = \partial C\):

\[
\mathrm{Tr}(\rho L_C) \approx N_\Sigma \cdot \bar{\ell}(t),
\qquad
\bar{\ell}(t) := \sum_\alpha p_\alpha \log d_\alpha,
\]

where \(N_\Sigma\) is the number of UV cut elements covering \(\Sigma\) and \(\bar{\ell}(t)\) is the \textbf{single-cell edge entropy} from the heat-kernel distribution (Theorem 6.20). Similarly, the geometric area is extensive:

\[
A(C) \approx N_\Sigma \cdot a_{\mathrm{cell}},
\]

where \(a_{\mathrm{cell}}\) is the area per UV cut element in the emergent metric.

Matching these expressions gives the \textbf{derived formula for Newton\textquotesingle s constant}:

\[
G = \frac{a_{\mathrm{cell}}}{4 \, \bar{\ell}(t)}
\]

where:

\begin{itemize}
\tightlist
\item
  \(a_{\mathrm{cell}}\) is fixed operationally from the UV correlation/mixing length \(\xi\) via \(a_{\mathrm{cell}} \sim \xi^2\) (from the exponential mixing hypothesis),
\item
  \(\bar{\ell}(t) = \sum_R p_R(t) \log d_R\) is computed from the heat-kernel edge distribution with \(p_R \propto d_R e^{-t\lambda_R}\).
\end{itemize}

Explicitly:

\[
\bar{\ell}(t) = \frac{\sum_R d_R e^{-t\lambda_R} \log d_R}{\sum_R d_R e^{-t\lambda_R}}.
\]

Thus \(G\) is the inverse edge-entropy density per geometric area, computable from the UV regulator and the reference-state Gibbs parameter \(t\), rather than a normalization convention.

\subsubsection{Entanglement equilibrium from MaxEnt}\label{55-entanglement-equilibrium-from-maxent}

MaxEnt selection implies that for variations preserving cap labels (fixed size and charges),

\[
\delta S_{\mathrm{gen}}(C) = 0.
\]

Using the split above and the first law for the bulk term,

\[
\delta S_{\mathrm{gen}}(C) = \delta \langle L_C \rangle + \delta \langle K_{\mathrm{bulk}} \rangle.
\]

\subsubsection{Einstein equation from cap equilibrium}\label{56-einstein-equation-from-cap-equilibrium}

In the d=4 scaling regime, the small-ball bridge is internal to the derived gravity chain once the separate bounded-interval projective branch is included. The only extra standard geometric input is the fixed-volume area-variation identity for a small geodesic ball. The modular kernel itself is not imported from a separate EFT law: on the extracted prime geometric subnet, a sufficiently small cap is the tangent causal diamond \(D_\ell\), whose preserving conformal Killing field is
\[
\xi_{D_\ell}
=
\frac{1}{2\ell}
\Bigl((\ell^2-r^2-t^2)\,\partial_t-2t\,x^i\partial_i\Bigr).
\]
On the t=0 slice its lapse is
\[
\xi_{D_\ell}\!\cdot n=\frac{\ell^2-r^2}{2\ell}.
\]
The bounded-interval kernel comes from combining the D3 geometric cap generator with the separate interval-preserving projective branch used downstream of the D4 half-line stress bridge, so the D3 geometric branch and the D4 stress bridge together yield
\[
\delta S_{\mathrm{bulk}}(C)
=
2\pi\int_{B_\ell}\frac{\ell^2-r^2}{2\ell}\,\delta\langle T_{00}\rangle\,d^3x
+
\delta\langle E^{(\eta)}_{C,\ell}\rangle.
\]
If \(\delta\langle T_{00}\rangle\) is approximately constant across the ball and the carried remainder is \(o(\ell^4)\), then
\[
\delta S_{\mathrm{bulk}}(C)
=
\frac{8\pi^2\ell^4}{15}\,\delta\langle T_{00}\rangle
+O(\ell^5\partial T)+o(\ell^4).
\]

Stationarity of generalized entropy therefore gives the observer-covariant scalar relation
\[
\delta\!\Bigl[(G_{ab}+\Lambda g_{ab})u^a u^b\Bigr]
=
8\pi G\,u^a u^b\,\delta\langle T_{ab}\rangle
\]
for the local diamond rest-frame four-velocity \(u^a\). In the adapted rest frame \(u^a=(1,0,0,0)\), this is
\[
\delta\!\left(G_{00}+\Lambda g_{00}\right)=8\pi G\,\delta\langle T_{00}\rangle.
\]

\subsubsection{Overlaps supply all timelike directions}\label{57-overlaps-supply-all-timelike-directions}

This scalar-to-tensor upgrade is internal once the Lorentz branch is in place. Overlapping observers through the same bulk point realize every local timelike four-velocity u. Define
\[
Y_{ab}:=G_{ab}+\Lambda g_{ab}-8\pi G\,\langle T_{ab}\rangle.
\]
The rest-frame relation of the previous subsection says
\[
u^a u^b\,\delta Y_{ab}=0
\]
for every such u. In local inertial coordinates,
\[
u^a=\gamma(1,v^i),
\qquad
|\vec v|<1,
\]
so
\[
0
=
u^a u^b\,\delta Y_{ab}
=
\gamma^2\Bigl(\delta Y_{00}+2v^i\delta Y_{0i}+v^iv^j\delta Y_{ij}\Bigr)
\]
for all \(|\vec v|<1\). The polynomial therefore vanishes coefficientwise, hence
\[
\delta Y_{00}=0,\qquad
\delta Y_{0i}=0,\qquad
\delta Y_{ij}=0.
\]
So the full tensor variation vanishes:
\[
\delta Y_{ab}=0.
\]

Along any connected scaling branch of reference states, \(Y_{ab}\) is therefore constant. The only purely local freedom left by the null reconstruction is the metric term, and fixing that term on one maximally symmetric reference state gives
\[
G_{ab}+\Lambda g_{ab}=8\pi G\,\langle T_{ab}\rangle.
\]
Thus the tensor upgrade does not require a separate EFT handoff once OPH supplies the geometric branch, the null bridge, and the overlap-complete timelike family.

\subsubsection{Non-tunable numerical constants}\label{58-non-tunable-numerical-constants}

The gravity chain yields specific numerical constants as rigid outputs of the axiom chain.

\textbf{The \(2\pi\) KMS normalization.} From Theorem 4.2, once the OPH refinement branch is geometric, cap modular flow is fixed by the KMS condition to the standard \(2\pi\) normalization. This is the same rigidity that fixes Unruh/Hawking temperature normalization. The period is determined by the stated scaling-limit package together with the geometric-branch hypothesis.

\textbf{The geometric coefficient \(\Omega_{d-2}/(d^2 - 1)\).} This coefficient appears in both (a) the CFT-ball modular Hamiltonian weight integral and (b) the geometric area-variation identity. It is an exact integral identity:

\[
\int_{B^{d-1}_\ell} \frac{\ell^2 - r^2}{2\ell} \, d^{d-1}x
= \frac{\Omega_{d-2} \, \ell^d}{d^2 - 1}.
\]

In \(d = 4\):

\[
\frac{\Omega_2}{4^2 - 1} = \frac{4\pi}{15} \approx 0.8377580409572781.
\]

This is the reason prefactors cancel cleanly when going from \(\delta S_{\mathrm{gen}} = 0\) to the Einstein equation (leaving \(8\pi G\) with the \(2\pi\) fixed by Theorem 4.2).

\textbf{What is predicted.} The framework cleanly separates:

\begin{itemize}
\tightlist
\item
  \textbf{Non-tunable constants}: \(2\pi\) (KMS period), \(\Omega_{d-2}/(d^2-1)\) (geometric coefficient), the existence of the Einstein form.
\item
  \textbf{Micro-dependent constants}: \(G\) (Newton\textquotesingle s constant) is the conversion between edge entropy and geometric area (a density of edge degrees of freedom per geometric area), which is model-dependent. \(\Lambda\) is fixed by the reference state/constraints. These require specifying the microscopic model.
\end{itemize}

\subsubsection{Quantitative Markov error and controlled corrections}\label{59-quantitative-markov-error-and-controlled-corrections}

The role of this subsection is to keep three distinct quantities separate:

\begin{enumerate}
\def\labelenumi{\arabic{enumi}.}
\tightlist
\item
the raw collar conditional mutual information
\[
\varepsilon_\delta:=I(A_\delta:D_\delta\mid B_\delta)_\omega;
\]
\item
the constructive recovery error
\[
r_{\mathrm{FR}}(\varepsilon_\delta)
:=
2\sqrt{1-e^{-\varepsilon_\delta}}
\le
2\sqrt{\varepsilon_\delta};
\]
\item
the exact-Markov replacement error on one fixed faithful collar model,
\[
\eta_\delta^{\mathrm M}
:=
4\lambda_\ast^{-1}\,
\delta^{\mathrm M}_{A_\delta:B_\delta:D_\delta}(\varepsilon_\delta),
\]
where \(\lambda_\ast>0\) is the lower spectral bound needed to compare modular Hamiltonians.
\end{enumerate}

The exact identity for the modular defect expectation is unchanged:
\[
\langle \Delta K_\delta\rangle_\omega
=
-I(A_\delta:D_\delta\mid B_\delta)_\omega
=
-\varepsilon_\delta,
\]
with
\[
\Delta K_\delta
:=
K_{A_\delta B_\delta D_\delta}
-K_{A_\delta B_\delta}
-K_{B_\delta D_\delta}
+K_{B_\delta}.
\]

What changes is the interpretation of later ``exact'' collar formulas. For each fixed finite-dimensional collar one may choose an exact Markov reference state \(\sigma_\delta\) with
\[
\|\omega_\delta-\sigma_\delta\|_1
\le
\delta^{\mathrm M}_{A_\delta:B_\delta:D_\delta}(\varepsilon_\delta),
\]
and a constructive recovered comparison state \(\omega_\delta^{\mathrm{rec}}\) with
\[
\|\omega_\delta-\omega_\delta^{\mathrm{rec}}\|_1
\le
r_{\mathrm{FR}}(\varepsilon_\delta).
\]
These two comparisons do different jobs:

\begin{quote}
\textbf{Bounded-observable control.} For every bounded observable \(X\) supported on \(A_\delta\cup B_\delta\),
\[
\left|
\operatorname{Tr}\!\left[X(\omega_\delta-\omega_\delta^{\mathrm{rec}})\right]
\right|
\le
\|X\|_\infty\,r_{\mathrm{FR}}(\varepsilon_\delta),
\]
and
\[
\left|
\operatorname{Tr}\!\left[X(\omega_\delta-\sigma_\delta)\right]
\right|
\le
\|X\|_\infty\,
\delta^{\mathrm M}_{A_\delta:B_\delta:D_\delta}(\varepsilon_\delta).
\]
The first bound is constructive and dimension-free; the second is the fixed-collar route to the exact Markov set.
\end{quote}

\begin{quote}
\textbf{Modular-additivity control.} If the relevant marginals are uniformly faithful with lower spectral bound \(\lambda_\ast>0\), then
\[
\|\Delta K_\delta(\omega)-\Delta K_\delta(\sigma_\delta)\|_\infty
\le
\eta_\delta^{\mathrm M}.
\]
Since \(\Delta K_\delta(\sigma_\delta)\) is central by exact Markovity, the finite-stage modular defect is within \(\eta_\delta^{\mathrm M}\) of the exact splice or additivity value.
\end{quote}

Therefore the manuscript's exact collar identities arise as limits with a carried remainder
\[
\eta_\delta
:=
r_{\mathrm{FR}}(\varepsilon_\delta)
+
\eta_\delta^{\mathrm M},
\]
together with the separate long-wavelength derivative remainders present in the small-ball expansion. Exact identities at finite collar width require exact Markovity; otherwise one works with \(\eta_\delta\)-controlled approximations and then lets \(\eta_\delta\to0\) in the controlled collar limit.

\textbf{Modified Einstein relation with an explicit carried remainder.} In the small-ball rest-frame calculation, write
\[
\delta S_C
=
\frac{8\pi^2\ell^4}{15}\,\delta\langle T_{00}\rangle
+
\delta\langle E_C^{(\delta)}\rangle
+
O(\ell^5\partial T),
\]
where \(E_C^{(\delta)}\) packages the carried collar remainder. Then the fixed-cap stationarity argument gives
\[
\delta\!\left(G_{00}+\Lambda g_{00}\right)
=
8\pi G\,\delta\langle T_{00}\rangle
+
\mathcal E_\delta,
\]
with
\[
\mathcal E_\delta
:=
\frac{15G}{\pi\ell^4}\,
\delta\langle E_C^{(\delta)}\rangle
+
O(\ell\,\partial T),
\]
and, by the bounds above,
\[
\bigl|
\delta\langle E_C^{(\delta)}\rangle
\bigr|
\le
C_C\,\eta_\delta
\]
for some collar-dependent constant \(C_C\) on the fixed faithful model.

If one wishes to package this carried remainder as an effective anomalous energy density, the natural definition is
\[
\delta\langle T_{00}^{\mathrm{anom}}\rangle
:=
\frac{15}{8\pi^2\ell^4}\,
\delta\langle E_C^{(\delta)}\rangle.
\]
Then
\[
\delta\!\left(G_{00}+\Lambda g_{00}\right)
=
8\pi G\,
\delta\!\left(
\langle T_{00}\rangle
+
\langle T_{00}^{\mathrm{anom}}\rangle
\right)
+
O(\ell\,\partial T).
\]
This is a bookkeeping definition of the finite-stage remainder, not an upgrade of the recovered-core theorem. Its significance is exactly that it is controlled:
\[
\bigl|
\delta\langle T_{00}^{\mathrm{anom}}\rangle
\bigr|
\le
\frac{15\,C_C}{8\pi^2\ell^4}\,
\eta_\delta.
\]

The framework carries the Markov error through three distinct controls:
\begin{enumerate}
\def\labelenumi{\arabic{enumi}.}
\tightlist
\item
\(r_{\mathrm{FR}}(\varepsilon_\delta)\) controls bounded observables at one stage;
\item
\(\delta^{\mathrm M}_{A_\delta:B_\delta:D_\delta}(\varepsilon_\delta)\) controls convergence to the exact Markov normal form on a fixed collar;
\item
\(\eta_\delta^{\mathrm M}\) controls modular-Hamiltonian replacements once faithfulness is assumed;
\item
the Einstein-branch remainder is a carried \(O(\eta_\delta)\) term, not a silent exact identity at finite cutoff.
\end{enumerate}

This is the concrete bridge from ``axioms about screens'' to ``precision GR predictions plus a disciplined carried remainder''.

\subsubsection{Focusing/QNEC internalization via relative entropy}\label{510-focusingqnec-internalization-via-relative-entropy}

Once the fixed-cutoff null modular structure of Section 5.2 is combined with the exact half-line generator/charge identification and the bounded-interval projective branch, focusing constraints follow from information-theoretic principles in the same scaling-limit branch rather than from a separate fixed-cutoff postulate.

\textbf{Derivation chain.} QNEC and focusing are supported conditionally within the same branch:

\begin{itemize}
\tightlist
\item
  Fixed-cutoff null modular bridge (\S{}5.2) \ensuremath{\Longrightarrow} exact-or-controlled strip additivity and a weak tail generator
\item
  Derived half-sided modular pair + exact half-line generator/charge identification \ensuremath{\Longrightarrow} local null-stress charge on half-lines and {[}K, P{]} = i 2\ensuremath{\pi} P
\item
  Relative entropy monotonicity \ensuremath{\Longrightarrow} QNEC
\item
  Einstein (Thm 5.1) + Raychaudhuri \ensuremath{\Longrightarrow} QFC for S\_gen
\end{itemize}

\textbf{Relative entropy monotonicity argument.} The key input is the monotonicity of relative entropy under partial trace, which is pure information theory:

\[
S(\rho_{AB} \| \sigma_{AB}) \ge S(\rho_A \| \sigma_A).
\]

For null deformations parameterized by \(\lambda\), consider nested null regions \(R(\lambda) \subset R(\lambda')\) obtained by varying the entangling cut along \(v\). The modular Hamiltonian \(K_\lambda\) generates the modular flow, and relative entropy satisfies convexity:

\[
\frac{d^2}{d\lambda^2} S(\rho_\lambda \| \sigma_\lambda) \ge 0.
\]

\begin{quote}
\textbf{Proposition 5.10a (Conditional internal QNEC).} Under the fixed-cutoff null bridge of \S{}5.2 together with the downstream density-upgrade and resulting null-stress-identification premises, the second null variation of von Neumann entropy satisfies

\[
\frac{d^2 S_{\mathrm{bulk}}}{d\lambda^2} \le 2\pi \langle T_{kk}(\lambda) \rangle,
\]

with the \(2\pi\) normalization fixed by Theorem 4.2.

\textbf{Proof.} Use the derived half-sided modular-inclusion statement of Corollary 5.2e and its Borchers-Wiesbrock consequence from Lemma 5.2f. This gives the translation structure with \([K, P] = - i 2\pi P\).

Consider the relative entropy \(S(\rho_\lambda \| \omega_\lambda)\) between the state \(\rho\) restricted to \(R(\lambda)\) and the reference state \(\omega\). Monotonicity under restriction to smaller regions (\(\lambda' > \lambda\)) gives:

\[
S(\rho_{R(\lambda)} \| \omega_{R(\lambda)}) \le S(\rho_{R(\lambda')} \| \omega_{R(\lambda')}).
\]

Using the first law \(\delta S = \delta \langle K \rangle\) and the geometric half-line modular form supplied by Section 5.2, expand to second order in the deformation. The convexity of relative entropy yields the QNEC inequality. The bound saturates for coherent states in the standard way. QED.
\end{quote}

\begin{quote}
\textbf{Corollary 5.10b (QFC for generalized entropy).} With the central area operator \(L_C\) from EC/MaxEnt (Section 5.4), define

\[
S_{\mathrm{gen}} = \mathrm{Tr}(\rho L_C) + S_{\mathrm{bulk}}.
\]

Given the conditional Einstein branch (Theorem 5.1) and the classical Raychaudhuri identity for null congruences, the Quantum Focusing Conjecture (QFC) follows within the same scaling regime: the generalized expansion \(\Theta_{\mathrm{gen}}\) is non-increasing along null generators.

\textbf{Proof sketch.} The Raychaudhuri equation relates expansion evolution to \(R_{kk}\). Einstein\textquotesingle s equation gives \(R_{kk} = 8\pi G (T_{kk} - \frac{1}{2}g_{kk}T)\). For null \(k\), this simplifies to \(R_{kk} = 8\pi G T_{kk}\). The QNEC (Prop 5.10a) then bounds the bulk entropy production, ensuring \(d\Theta_{\mathrm{gen}}/d\lambda \le 0\). QED.
\end{quote}

\textbf{Significance.} This section shows how the Recoverable Generalized Entropy focusing package is internally supported inside the same null-modular branch once the stress-tensor identification and Einstein relation are in place. That package remains part of the stated axiom set.

\textbf{Theorem 5.1 (Observer-consistency implies a conditional scaling-limit Einstein branch).} Under the five OPH axioms, the collar and null-strip hypotheses of Sections 2 and 5, the hypotheses of Theorem 4.2, the canonical technical premise \textbf{T3}, the half-line generator/charge identification of Section 5.2 together with the bounded-interval projective branch used there, and admissible first variations about a maximally symmetric reference state, the cap equilibrium condition implies the rest-frame first-variation relation

\[
\delta\!\left(G_{00}+\Lambda g_{00}\right)=8\pi G\,\delta\langle T_{00}\rangle.
\]

If this relation holds for all local directions and reference states in the locally Lorentzian scaling regime, overlap consistency upgrades it to the semiclassical Einstein equation modulo the expected \(\Lambda g_{ab}\) ambiguity.

\textbf{Proof sketch.} The cap first law identifies \(\delta S_C\) with \(\delta\langle K_C\rangle\), while the conditional null bridge together with the bounded-interval projective branch relates the bulk modular variation to the required null stress-tensor charge on the diamond. Since the reference state is maximally symmetric, the fixed-volume small-ball area identity enters at first order in the perturbation, so entanglement equilibrium yields the displayed rest-frame relation. Overlap consistency then supplies all local directions, and the null-to-tensor lemma upgrades the relation to the tensor equation modulo \(\Lambda g_{ab}\). QED.

\subsubsection{Black-hole structural statements and continuation-level spectroscopy}\label{511-discrete-horizon-area-spectrum-and-hawking-emission-established}

This subsection mixes two claim tiers and is organized accordingly. The structural black-hole statements already supported in the manuscript are: the edge-center/collar obstruction to a naive inside-outside tensor factorization, the small-CMI recoverability statement at the level proved earlier, the Hawking/KMS normalization fixed by the geometric modular branch, the discrete area spectrum from integer edge-sector dimensions, and the Schwarzschild transition identity \(\Delta E = k_B T_H \ln(d'/d)\) for an actual sector change \(d \to d'\). Any clean comb structure, QNM selector, linewidth estimate, PBH burst template, or Kerr/LIGO horizon spectroscopy template remains continuation-level and uses extra premises stated below.

\textbf{Derived structural statements.}

\textbf{Area eigenvalues from edge sectors.} The central area operator (Section 5.4) is

\[
L_C = \sum_\alpha (\log d_\alpha) P_\alpha,
\]

where \(d_\alpha \in \mathbb{N}\) is the dimension of the edge Hilbert space in sector \(\alpha\). With the normalization \(\mathrm{Tr}(\rho L_C) = \langle A
\rangle / 4G\), the area eigenvalues are

\[
A_\alpha = 4G \log d_\alpha = 4\ell_p^2 \ln d_\alpha,
\]

where \(\ell_p^2 = \hbar G/c^3\) is the Planck area. Since \(d_\alpha\) is a positive integer, \textbf{areas are discretely spaced} with logarithmic gaps.

\textbf{Hawking emission energy quantization.} For a Schwarzschild black hole with \(A(M) = 16\pi G^2 M^2/c^4\), a transition between sectors \(d \to d'\) changes the area by

\[
\Delta A = 4\ell_p^2 \ln(d'/d).
\]

The corresponding energy at infinity is \(\Delta E = c^2 \Delta M\), with \(\Delta M = \Delta A / (dA/dM)\). This gives

\[
\Delta E = \frac{\hbar c^3}{8\pi G M} \ln(d'/d).
\]

Using the Hawking temperature \(T_H = \hbar c^3 / (8\pi G k_B M)\), whose \(2\pi\) normalization is fixed by the BW\(_{S^2}\) tangent-limit normalization of Theorem 4.2:

\[
\Delta E = k_B T_H \ln(d'/d).
\]

\textbf{Continuation-level spectroscopy templates.} The remainder of this subsection is not part of the recovered core. It records what follows only if one adds extra discrete-horizon selection rules and, where stated, standard semiclassical inputs.

\textbf{Integer transitions.} If dominant transitions multiply the edge dimension by an integer \(k\) (i.e., \(d'/d = k\)), the spectrum becomes a discrete comb:

\[
\Delta E_k = k_B T_H \ln k, \qquad
\Delta f_k = \frac{c^3}{16\pi^2 G M} \ln k.
\]

\textbf{Structural condition: comb vs. generic discreteness.} The log-integer \emph{comb} structure requires the additional dynamical assumption that integer-multiplication transitions (d \ensuremath{\to} kd) dominate. If generic transitions between arbitrary integers dominate instead, the set of ln(d\textquotesingle/d) values becomes a dense log-rational set that may appear quasi-continuous after folding in linewidths and astrophysical effects. What follows directly from the axioms is the \emph{discrete area spectrum}; the clean comb pattern is conditional on the selection rule.

\textbf{Continuation-level template (Discrete Hawking spectrum).} If the additional integer-transition selection rule is realized, the Hawking emission spectrum consists of discrete lines with spacing \(\Delta E_k = k_B T_H \ln k\), where \(k\) is an integer characterizing the dominant sector transitions, instead of a continuous thermal profile.

\textbf{Mass-independent fractional linewidth (continuation-level estimate).} Using Page\textquotesingle s semiclassical calculation for emission power \(P(M) = p_0 \hbar c^6 / (G^2 M^2)\) with \(p_0 \approx 2 \times 10^{-4}\), the emission rate is \(\dot{N} \approx P /
\langle E \rangle\) where \(\langle E \rangle = a \, k_B T_H\) with \(a \sim
\mathcal{O}(1-10)\). The natural linewidth \(\Gamma \sim \hbar \dot{N}\) divided by the level spacing gives:

\[
\frac{\Gamma}{\Delta E_k} \approx \frac{64\pi^2 p_0}{a \ln k} \approx 3-5\%
\]

Within this continuation branch, the emission lines are narrow (few-percent fractional width) and the estimated fraction is mass-independent.

\textbf{Connection to quasinormal modes (interpretive continuation).} The highly-damped Schwarzschild quasinormal modes have asymptotic real part (Motl, 2002):

\[
\mathrm{Re}\,\omega \to \frac{c^3}{8\pi G M} \ln 3.
\]

This matches exactly the \(k = 3\) transition frequency \(\Delta E_3 / \hbar\). If one adopts a Bohr-type identification between quantum transition frequencies and asymptotic QNM frequencies, this selects

\[
\Delta A = 4\ell_p^2 \ln 3 \approx 4.39 \, \ell_p^2
\]

as the fundamental area quantum.

\textbf{Conditionality statement.} The area quantization follows from the edge-sector structure (derived). The \(k = 3\) selection requires the additional interpretive identification with QNM frequencies (not derived from axioms). The linewidth prediction uses standard semiclassical inputs.

\textbf{Numerical examples.} For \ensuremath{\Delta}f\_k = (c\textsuperscript{3}/16\ensuremath{\pi}\textsuperscript{2}GM) ln k:

\begin{itemize}
\tightlist
\item
  \textbf{M = 30 M\ensuremath{\odot}}: k=2 at 29.7 Hz, k=3 at 47.1 Hz
\item
  \textbf{M = 1 M\ensuremath{\odot}}: k=2 at 891 Hz, k=3 at 1412 Hz
\item
  \textbf{M = 10\textsuperscript{1}\textsuperscript{2} kg (primordial)}: k=2 at 7.3 MeV, k=3 at 11.6 MeV
\end{itemize}

Within the integer-transition continuation, these frequencies track \(k_B T_H \ln k\) exactly and are in principle distinguishable from a continuous thermal spectrum.

\textbf{Continuation-level PBH burst search template.}

Within the discrete-horizon continuation, the Hawking comb would provide a distinctive gamma-ray template. A template-level discriminant is \textbf{log-integer energy ratios}: if two emission lines are observed at energies \(E_2\) and \(E_3\), their ratio must satisfy

\[
\frac{E_3}{E_2} = \frac{\ln 3}{\ln 2} \approx 1.585
\]

exactly, independent of black hole mass. Within that continuation template, the ratio is fixed once the log-integer rule is assumed.

\textbf{Available instruments and energy coverage.} The \(k = 2\) line energy \(E_2 = k_B T_H \ln 2\) determines which instruments can see a given BH mass:

{\def\LTcaptype{none} % do not increment counter
\begin{longtable}[]{@{}>{\RaggedRight\arraybackslash}p{0.320\textwidth}>{\RaggedRight\arraybackslash}p{0.320\textwidth}>{\RaggedRight\arraybackslash}p{0.320\textwidth}@{}}
\toprule\noalign{}
Instrument & Energy band & BH mass range (k=2 in band) \\
\midrule\noalign{}
\endhead
\bottomrule\noalign{}
\endlastfoot
Fermi GBM (BGO) & 0.15--40 MeV & 2\ensuremath{\times}10\textsuperscript{1}\textsuperscript{1}--5\ensuremath{\times}10\textsuperscript{1}\textsuperscript{3} kg \\
Fermi LAT & 0.1--300 GeV & 2\ensuremath{\times}10\textsuperscript{7}--7\ensuremath{\times}10\textsuperscript{1}\textsuperscript{0} kg \\
H.E.S.S. & 0.1--100 TeV & 7\ensuremath{\times}10\textsuperscript{4}--7\ensuremath{\times}10\textsuperscript{7} kg \\
LHAASO-WCDA & 1--15 TeV & 5\ensuremath{\times}10\textsuperscript{5}--7\ensuremath{\times}10\textsuperscript{6} kg \\
\end{longtable}
}

\textbf{Detector resolution vs. intrinsic linewidth.} The continuation-level linewidth estimate is 3--5\% (mass-independent). Current detector energy resolutions:

\begin{itemize}
\tightlist
\item
  Fermi GBM: \(< 10\%\) (0.1--1 MeV), \(\sim 4\%\) at 10 MeV (BGO)
\item
  Fermi LAT: \(< 10\%\) (1--100 GeV)
\item
  H.E.S.S.: \(\sim 15\%\) (TeV)
\item
  LHAASO-WCDA: \(\sim 33\%\) (TeV)
\end{itemize}

Within this template, the comb could be resolvable with GBM/LAT; at TeV energies it would appear as moderately broad bumps rather than sharp lines.

\textbf{Search protocol.} A dedicated OPH-comb search would:

\begin{enumerate}
\def\labelenumi{\arabic{enumi}.}
\tightlist
\item
  Select burst-like candidates (10--120 s time windows, matching existing PBH burst search protocols).
\item
  Fit each candidate with null model (smooth continuum) vs. OPH comb model (peaks at \(E_k = E_0 \ln k\) convolved with detector response).
\item
  Scan over the single scale parameter \(E_0 = k_B T_H\) (equivalently, BH mass).
\item
  Require at least two lines satisfying log-integer ratio to claim detection.
\item
  Correct significance for trials (time windows \(\times\) sky positions \(\times\) \(E_0\) scan).
\end{enumerate}

\textbf{Observational context.} Dedicated PBH burst searches (H.E.S.S., LHAASO) report \textbf{no significant bursts}. An OPH-specific comb-template analysis of archival data would:

\begin{itemize}
\tightlist
\item
  Set upper limits on OPH-comb PBH burst rates
\item
  Demonstrate direct testability of the discrete spectrum prediction
\item
  Provide constraints comparable to or stronger than generic PBH burst limits
\end{itemize}

\textbf{Data availability.} Fermi GBM provides public Time-Tagged Event (TTE) burst data; Fermi LAT provides public photon event lists with documented analysis workflows. H.E.S.S. has a small public test data release.

\textbf{Continuation-level GW horizon spectroscopy template for Kerr remnants.}

The same continuation-level discrete-horizon branch extends to gravitational wave observables. For Kerr black holes, the thermodynamic first law is \(\delta M = T_H \delta S +
\Omega_H \delta J\), so the entropy change for absorbing a quantum with frequency \(\omega\) and azimuthal number \(m\) is:

\[
\delta S = \frac{\hbar(\omega - m\Omega_H)}{k_B T_H}.
\]

In the edge-sector framework, \(\delta S = \ln(d'/d)\), so the discreteness condition becomes:

\[
\hbar(\omega - m\Omega_H) = k_B T_H \ln k, \qquad k \in \{2, 3, 4, \ldots\}
\]

Under those additional discrete-horizon premises, this gives the \textbf{GW horizon spectroscopy comb}: a continuation-level set of discrete resonant frequencies where the horizon can efficiently absorb or emit energy.

\textbf{Kerr line frequencies.} For a remnant with mass \(M\) and dimensionless spin \(\chi = a_*/M\), define the spin correction factor:

\[
g(\chi) = \frac{2\sqrt{1-\chi^2}}{1+\sqrt{1-\chi^2}}, \qquad
\Omega_H(M,\chi) = \frac{c^3}{2GM} \cdot \frac{\chi}{1+\sqrt{1-\chi^2}}.
\]

The line frequencies are:

\[
f_{k,m}(M,\chi) = \frac{m \, \Omega_H(M,\chi)}{2\pi} + \frac{c^3}{16\pi^2 GM} \, g(\chi) \, \ln k
\]

Within this continuation template, once LIGO/Virgo infers \((M, \chi)\) for a remnant, the line pattern is fixed by the inferred remnant parameters and the assumed discrete-horizon rule. This is not a recovered-core theorem.

\textbf{Line weights from GR envelope + discretization.} In this continuation template, the line strengths are modeled by matching to the known GR greybody absorption spectrum in the semiclassical limit. The discretization rule gives bin width \(\Delta\omega_k \approx \omega_T \ln(1 + 1/k)\) where \(\omega_T = k_B T_H/\hbar\). The net line weight (absorption minus stimulated emission) is:

\[
W^{\mathrm{net}}_{k,\ell m} = \Gamma^{\mathrm{GR}}_{\ell m}(\omega_{k,m}) \cdot \Delta\omega_k \cdot \frac{k-1}{k}
\]

where \(\Gamma^{\mathrm{GR}}_{\ell m}\) is the standard GR greybody factor and the \((k-1)/k\) factor arises from KMS detailed balance with \(e^{(\omega-m\Omega_H)/T_H} = k\).

\textbf{Universal stacking coordinate.} Define the dimensionless rescaled frequency:

\[
x := \frac{GM}{c^3 g(\chi)}(\omega - m\Omega_H).
\]

Then the predicted line locations collapse to universal constants:

\[
x_k = \frac{\ln k}{8\pi} \qquad (k = 2, 3, 4, \ldots)
\]

Numerically: \(x_2 = 0.02758\), \(x_3 = 0.04371\), \(x_4 = 0.05516\), \(x_5 = 0.06404\).

\textbf{Stacking test.} Multiple BBH events can be mapped to this universal \(x\) coordinate and stacked. If the comb is real, peaks align across events with different \((M, \chi)\); detector noise does not stack coherently.

\textbf{Comparison to existing work.} Prior area-quantization searches~\cite{laghi21} used parameterized models with one free spacing constant. The OPH prediction is more constrained: multiple lines with exact \(\ln k\) ratios, plus the \((k-1)/k\) weight hierarchy from detailed balance.

\textbf{Numerical example (GW170608).} Remnant parameters: \(M_f \approx 18.0 M_\odot\), \(\chi_f \approx 0.69\). For \(m = 2\), the horizon rotation frequency is \(m\Omega_H/(2\pi) \approx 719\) Hz. The \textbf{thermal comb spacing} (the part that encodes the area quantization) is:

{\def\LTcaptype{none} % do not increment counter
\begin{longtable}[]{@{}>{\RaggedRight\arraybackslash}p{0.320\textwidth}>{\RaggedRight\arraybackslash}p{0.320\textwidth}>{\RaggedRight\arraybackslash}p{0.320\textwidth}@{}}
\toprule\noalign{}
k & \(\Delta f_k := \frac{c^3 g(\chi)}{16\pi^2 GM} \ln k\) (Hz) & Relative weight \((k-1)/k\) \\
\midrule\noalign{}
\endhead
\bottomrule\noalign{}
\endlastfoot
2 & 41.6 & 0.500 \\
3 & 65.9 & 0.667 \\
4 & 83.2 & 0.750 \\
5 & 96.5 & 0.800 \\
6 & 107.5 & 0.833 \\
\end{longtable}
}

The full physical frequencies are \(f_{k,2} = 719 + \Delta f_k\) Hz (i.e., 760--827 Hz), outside LIGO\textquotesingle s most sensitive band for this remnant. However, the \textbf{stacking analysis} uses the rescaled coordinate \(x = GM(\omega -
m\Omega_H)/(c^3 g(\chi))\), which maps the thermal spacing to universal constants \(x_k = \ln k / 8\pi\) regardless of the rotation offset.

\textbf{Template-matching criterion.} After rescaling by \((M, \chi)\), spectral features in this continuation template must satisfy \(f_k/f_2 =
\ln k / \ln 2\) exactly, independent of remnant parameters. Absence of coherent stacking at the predicted \(x_k\) values would challenge the discrete-horizon continuation template, not by itself the derived area-spectrum statement.

\subsubsection{Classical mechanics from emergent GR}\label{512-classical-mechanics-from-emergent-gr}

Once the Einstein equation is established, the framework inherits standard GR consequences. This section makes explicit how classical mechanics emerges.

\textbf{Stress-energy conservation is automatic.} The contracted Bianchi identity is geometric:

\[
\nabla^a G_{ab} = 0.
\]

Combined with the Einstein equation, this implies:

\[
\nabla^a \langle T_{ab} \rangle = 0.
\]

\textbf{Geodesic motion from dust limit.} For pressureless classical matter ("dust"), \(T^{ab} = \rho \, u^a u^b\). Conservation yields:

\[
\nabla_a(\rho u^a u^b) = 0
\quad \Rightarrow \quad
u^b \nabla_a(\rho u^a) + \rho \, u^a \nabla_a u^b = 0.
\]

Projecting orthogonally to \(u^b\) using \(h^b{}_c = \delta^b{}_c + u^b u_c\) kills the first term, giving:

\[
\rho \, u^a \nabla_a u^b = 0
\quad \Rightarrow \quad
u^a \nabla_a u^b = 0.
\]

This is the geodesic equation: free classical bodies follow spacetime geodesics.

\textbf{Newtonian limit from weak-field GR.} Take the weak-field, slow-motion limit with metric:

\[
g_{00} \approx -(1 + 2\Phi/c^2), \qquad
g_{0i} \approx 0, \qquad
g_{ij} \approx \delta_{ij}(1 - 2\Phi/c^2),
\]

and velocities \(|\mathbf{v}| \ll c\). Then \(G_{00} \approx 2\nabla^2\Phi/c^2\) (leading order), and \(T_{00} \approx \rho c^2\). The Einstein equation reduces to:

\[
\nabla^2 \Phi = 4\pi G \rho.
\]

Geodesic motion reduces to:

\[
\ddot{\mathbf{x}} = -\nabla \Phi.
\]

These are Newton\textquotesingle s gravitational law and Newton\textquotesingle s second law. Classical mechanics is recovered as a controlled limit of the emergent GR dynamics.

\textbf{Precision classical predictions.} Once the field equation is fixed to Einstein form, the framework inherits the standard GR precision toolbox (post-Newtonian expansion, lensing, time delay, etc.), with no free "shape" parameters beyond \(G\) and \(\Lambda\).

Selected precision predictions (in the regime where the GR derivation applies):

\emph{Light bending by mass \(M\):} For impact parameter \(b\),

\[
\Delta\theta = \frac{4GM}{c^2 b}.
\]

For the Sun with \(b \approx R_\odot\): \(\Delta\theta \approx 1.751\) arcsec.

\emph{Mercury perihelion advance:} Per orbit,

\[
\Delta\varpi = \frac{6\pi GM}{a(1-e^2)c^2}.
\]

Using Mercury\textquotesingle s orbital parameters: \(\Delta\varpi \approx 42.98\) arcsec/century.

\emph{Gravitational redshift:} Between two radii in a static potential,

\[
\frac{\Delta\nu}{\nu} \approx \frac{\Delta\Phi}{c^2}.
\]

For the Sun (surface to infinity): \(z \approx 2.12 \times 10^{-6}\).

These predictions are fixed functions of \(G\) and known source parameters, and are confirmed observationally to high precision. The framework inherits them automatically once the Einstein equation is derived.

\subsubsection{Precision gravity predictions and experimental bounds}\label{513-precision-gravity-predictions-and-experimental-bounds}

The gravity sector makes symmetry-protected exact-zero predictions that can be confronted with the tightest available experimental bounds. This section translates the theoretical predictions into the specific observables that experiments actually constrain.

\textbf{Speed of gravitational waves.} The derived GR regime implies massless gravitons propagating on the same null cones as photons:

\[
\frac{c_{\mathrm{GW}} - c}{c} = 0 \text{ exactly.}
\]

\emph{Current bound (GW170817 + GRB 170817A multi-messenger):}

\[
-3 \times 10^{-15} < \frac{c_{\mathrm{GW}} - c}{c} < +7 \times 10^{-16}
\quad (90\% \text{ credibility}).
\]

For a source at \(\sim 40\) Mpc, this fractional difference corresponds to only a few seconds of propagation-time mismatch across \(\sim 10^8\) years of travel.

\textbf{Graviton mass.} The gauge redundancy (diffeomorphism invariance) forbids a hard mass term:

\[
m_g = 0 \text{ exactly.}
\]

\emph{Current bound (GW dispersion analysis, PDG 2025):}

\[
m_g \le 1.76 \times 10^{-23} \text{ eV}/c^2 \quad (90\% \text{ credibility}).
\]

This corresponds to a reduced Compton wavelength \(\bar{\lambda}_C \gtrsim 1.6 \times 10^{16}\) m, i.e., order \(\sim 1.6\) light-years.

\textbf{No dipole radiation.} Many modified gravity theories predict extra channels (scalar/vector) producing dipolar radiation at \((-1)\)PN order. The derived GR limit predicts no such channel.

\emph{Current bound (GW170817 inspiral phasing, PDG 2025):}

\[
-4 \times 10^{-6} < \delta\hat{p}_{-2} < 2 \times 10^{-5}
\quad (90\% \text{ credibility}).
\]

\textbf{Only tensor polarizations.} The GR outcome means only the two tensor (helicity-2) modes propagate. Pure non-tensor hypotheses are disfavored by observational constraints, and mixed tensor-scalar/vector models are tightly constrained.

\textbf{Equivalence principle tests.} Additional null checks from the derived GR structure:

\begin{itemize}
\tightlist
\item
  Universality of free fall (space tests): precision \textasciitilde10\textsuperscript{-}\textsuperscript{1}\textsuperscript{5}
\item
  Nordtvedt parameter (\ensuremath{\eta} = 4\ensuremath{\beta} - \ensuremath{\gamma}): (0.47 \ensuremath{\pm} 0.55) \ensuremath{\times} 10\textsuperscript{-}\textsuperscript{4}
\item
  Binary pulsar radiative damping (PSR J0737-3039): 0.999963 \ensuremath{\pm} 0.000063
\end{itemize}

\subsubsection{Theory-side error propagation from Markov bounds}\label{514-theory-side-error-propagation-from-markov-bounds}

The framework provides exact-zero predictions and quantitative control over how well those predictions hold. The Markov/recovery machinery can be propagated through the entire GR emergence chain.

\textbf{The key quantitative hook.} From Theorem 3.1, if the target state satisfies

\[
I(A_k : C_k \mid B_k) \le \varepsilon_k,
\]

then recovery maps exist with trace-distance error

\[
\delta_k = 2\sqrt{\ln 2 \cdot \varepsilon_k}.
\]

Trace distance gives immediate bounds on observable errors. Using the standard dual norm inequality:

\[
|\langle O \rangle_\rho - \langle O \rangle_\sigma|
\le \|O\|_\infty \|\rho - \sigma\|_1 = 2 \|O\|_\infty D(\rho, \sigma),
\]

where \(D(\rho, \sigma) = \frac{1}{2}\|\rho - \sigma\|_1\) is the trace distance.

\textbf{Exponential decay from the mixing hypothesis.} The mixing assumption (Section 2.3) provides:

\[
I_\omega(A_\delta : D_\delta \mid B_\delta)
\le c \cdot |\partial C|_{\mathrm{UV}} \cdot e^{-\delta/\xi}.
\]

Combining these gives an explicit precision dial:

\[
\delta_{\mathrm{step}} \lesssim 2\sqrt{\ln 2 \cdot c \cdot |\partial C|_{\mathrm{UV}}}
\cdot e^{-\delta/(2\xi)}.
\]

\textbf{What precision requires.} To match the GW speed bound (\(\sim 10^{-15}\) fractional accuracy), the recovery-map error must satisfy:

\[
\delta \lesssim 10^{-15}
\quad \Rightarrow \quad
\varepsilon \lesssim \frac{(\delta/2)^2}{\ln 2} \approx 3.6 \times 10^{-31}.
\]

This is extremely small, but achievable: with a macroscopic boundary (\(|\partial C|_{\mathrm{UV}} \sim 10^{35}\) for a meter-scale boundary at Planck UV scale), the exponential decay \(e^{-\delta/\xi}\) with \(\delta/\xi \sim\) a few hundred easily pushes below \(10^{-31}\) once the prefactor is included.

\textbf{Precision summary.} The framework provides:

\begin{enumerate}
\def\labelenumi{\arabic{enumi}.}
\tightlist
\item
  Exact-zero predictions (\(m_g = 0\), \(c_{\mathrm{GW}} = c\)) from symmetry protection.
\item
  Translation of those zeros into the specific observables experiments constrain.
\item
  Explicit bounds on how far derived geometric statements can drift, using the conditional mutual information \ensuremath{\to} trace distance \ensuremath{\to} observable error chain.
\end{enumerate}

This is the concrete path from "axioms about screens" to "precision GR predictions with quantitative error control."

\subsubsection{Dark-sector response from the modular anomaly}\label{515-dark-matter-as-modular-anomaly}

This subsection studies the benchmark scale obtained by adding a specific deep-IR response ansatz on top of the modular-anomaly term. Neither the observational dark-matter identification nor MOND/RAR-like dynamics are derived here.

The modular anomaly term \(T_{ab}^{\mathrm{anom}}\) derived in Section 5.9 supplies one structural ingredient for a possible dark-sector continuation, without introducing new particle species.

\textbf{The identification.} The anomalous stress-energy contribution

\[
\langle T_{00}^{\mathrm{anom}} \rangle
= \frac{15}{8\pi^2} \cdot \frac{\delta \langle K_C^{\mathrm{(anom)}} \rangle}{\ell^4}
\]

is "dark" by construction: it arises from information-theoretic/gravitational structure (modular Markov imperfections), not from Standard Model fields. It gravitates but does not couple electromagnetically. These properties explain why the term is explored as a dark-sector ingredient, but they do not by themselves close the observational identification.

\textbf{Connection to the cosmological constant.} The framework makes \(\Lambda\) a global capacity parameter:

\[
\Lambda = \frac{3\pi}{G N_{\mathrm{scr}}},
\quad
r_{dS} = \sqrt{\frac{3}{\Lambda}}.
\]

This introduces an unavoidable IR length scale \(r_{dS}\). Any galaxy-scale continuation built from the anomaly term would therefore be an IR phenomenon, appearing only when accelerations are small and distances are large.

\textbf{Acceleration benchmark.} If one assumes the relevant deep-IR response is controlled only by \(r_{dS}\), \(c\), and the fixed anomaly prefactor, then the benchmark scale must:

\begin{enumerate}
\def\labelenumi{\arabic{enumi}.}
\tightlist
\item
  Vanish if \(r_{dS} \to \infty\) (infinite capacity, no de Sitter scale)
\item
  Be controlled by \(r_{dS}\) as the only new IR scale
\item
  Carry non-tunable coefficients from the derivation
\end{enumerate}

The anomaly enters with prefactor \(\frac{15}{8\pi^2}\). The corresponding benchmark acceleration scale constructible from \((\Lambda, c)\) is:

\[
\boxed{
a_0^{\mathrm{(OPH)}} := \frac{15}{8\pi^2} \cdot c^2 \sqrt{\frac{\Lambda}{3}}
= \frac{15}{8\pi^2} \cdot \frac{c^2}{r_{dS}}
}
\]

\textbf{Normalization estimate.} Using Planck 2018 \(\Lambda\)CDM parameters (\(H_0 \approx 67.4\) km/s/Mpc, \(\Omega_\Lambda \approx 0.685\)):

\begin{itemize}
\tightlist
\item
  \(\Lambda \approx 1.09 \times 10^{-52}\) m\(^{-2}\)
\item
  \(r_{dS} \approx 1.66 \times 10^{26}\) m
\item
  Therefore:
\end{itemize}

\[
\boxed{
a_0^{\mathrm{(OPH)}} \approx 1.03 \times 10^{-10} \text{ m/s}^2
}
\]

For comparison, observational fits to galaxy regularities (RAR/MDAR/MOND phenomenology) quote \(a_0 \sim 1.2 \times 10^{-10}\) m/s\(^2\). This numerical proximity does not by itself derive galaxy dynamics.

\textbf{One illustrative response ansatz.} If one further assumes that \(T_{00}^{\mathrm{anom}}\) is the dominant deep-IR source and that the response organizes into a MOND/RAR-like law, the Newtonian limit could be written as:

\[
\nabla^2 \Phi = 4\pi G (\rho_b + \rho_{\mathrm{anom}}),
\]

i.e., baryons plus an effective extra density. Under the same extra ansatz the radial acceleration relation (RAR) could be written as:

\[
g_{\mathrm{obs}} \approx g_b + \sqrt{a_0 \cdot g_b},
\quad
g_{\mathrm{DM}} := g_{\mathrm{obs}} - g_b \approx \sqrt{a_0 \cdot g_b}.
\]

With \(a_0 = a_0^{\mathrm{(OPH)}}\) fixed, that same ansatz would imply:

\textbf{(i) Baryonic Tully-Fisher relation.}

\[
V^4 \approx G \cdot M_b \cdot a_0^{\mathrm{(OPH)}}
\]

where \(V\) is the asymptotic rotation velocity and \(M_b\) is baryonic mass.

\textbf{(ii) Flat rotation curves.} For a point mass \(M_b\):

\[
g_{\mathrm{DM}}(r) = \frac{\sqrt{G M_b \, a_0^{\mathrm{(OPH)}}}}{r}
\quad \Rightarrow \quad
M_{\mathrm{DM}}(r) \propto r
\]

i.e., inferred dark mass grows linearly with radius, producing flat rotation curves.

\textbf{(iii) Characteristic surface density.}

\[
\Sigma_0^{\mathrm{(OPH)}} = \frac{a_0^{\mathrm{(OPH)}}}{2\pi G}
\approx 0.25 \text{ kg/m}^2 \approx 120 \, M_\odot/\text{pc}^2.
\]

This lies in the range of observed central halo surface densities, but it is not a derived halo theorem.

\textbf{Status.} What is grounded in the framework developed here:

\begin{itemize}
\tightlist
\item
  The modular anomaly term exists with fixed coefficient \(\frac{15}{8\pi^2}\)
\item
  \(\Lambda\) and \(r_{dS}\) are determined by screen capacity
\end{itemize}

What remains missing for theorem-level closure:

\begin{itemize}
\tightlist
\item
  A controlled nonrelativistic limit from the anomaly term to galaxy observables
\item
  A derived response law selecting the MOND/RAR functional rather than alternative IR behavior
\item
  Lensing, cluster, and Bullet-Cluster phenomenology
\item
  Cosmological abundance and structure-formation checks
\item
  Environment-dependence and stability control
\end{itemize}

This branch contains structural ingredients and an IR benchmark, but it does not yet derive the precise galaxy-scale response from the recovered core.

\textbf{Falsifiability.} If later work cannot supply the missing derivation steps, or if any resulting continuation is incompatible with galaxy, lensing, cluster, Bullet-Cluster, or cosmological data, then the modular-anomaly dark-sector continuation retracts. The recovered core theorem package does not.

\subsubsection{De Sitter holography: static patch vs boundary-at-infinity}\label{516-de-sitter-holography-static-patch-vs-boundary-at-infinity}

A natural question arises: how does this framework relate to the ``unsolved problem'' of de Sitter holography?

\textbf{What the usual dS holography problem is.} When people say ``dS holography is unsolved,'' they typically mean that we do not have anything as sharp as AdS/CFT, where the bulk has a timelike asymptotic boundary supporting a well-defined dual CFT with a precise dictionary. For de Sitter, there is no asymptotic timelike boundary in the static patch where one can simply place the dual theory. The classic dS/CFT proposal at future infinity has familiar difficulties, including non-unitarity worries and complex conformal weights.

\textbf{What this model does differently.} The framework begins with an observer's static patch and its horizon screen \(S^2\), building a net of subregion algebras on that screen. At finite cutoff those algebras are type-I regulators. The Lorentz branch, when invoked, is a scaling-limit statement about the refinement-limit observer net, and that limit may leave the regulator class. On the extracted-prime-geometric branch of Theorem~\ref{42-theorem-mathrmbw_s2-from-markov-locality-symmetry-regularity}, the realized scaling-limit cap modular action is geometric and, in the non-type-I case of interest, generally outer. That identification is \emph{not} derived from MaxEnt alone; it is exactly the explicit \textbf{T2} plus extraction-and-rigidity boundary.

This is therefore a fundamental fork away from AdS/CFT-style holography:

{\def\LTcaptype{none}
\begin{longtable}[]{@{}>{\RaggedRight\arraybackslash}p{0.480\textwidth}>{\RaggedRight\arraybackslash}p{0.480\textwidth}@{}}
\toprule\noalign{}
AdS/CFT & This framework \\
\midrule\noalign{}
\endhead
\bottomrule\noalign{}
\endlastfoot
Codimension-1 boundary at infinity & Codimension-2 horizon screen (\(S^2\)) \\
Single global boundary theory & Observer-dependent patches that overlap \\
Dual CFT required & Only algebras + consistency conditions \\
Negative \(\Lambda\) & Positive \(\Lambda\) natural \\
\end{longtable}
}

This aligns with the static-patch/complementarity intuition in the dS literature, where the fundamental description is patch-based and different static patches are related by consistency rules, not by a single global boundary theory.

\textbf{The mechanism: \(\Lambda\) as global capacity, not local physics.} A key structural result is that null modular data reconstruct the stress tensor only up to an additive metric term. This is the statement that vacuum-energy or cosmological-constant shifts are invisible to the local null-data route. The Einstein equation derived from entanglement equilibrium is therefore fixed only up to \(\Lambda g_{ab}\).

For that reason \(\Lambda\) cannot be determined by local consistency alone. It must be fixed by a \emph{global} input: the total number of degrees of freedom on the screen. The de Sitter link is
\[
\Lambda=\frac{3\pi}{G N_{\mathrm{scr}}}.
\]
If the screen Hilbert space has finite dimension, the natural semiclassical interpretation is a cosmological horizon with finite entropy. The local/global stack is therefore: local null-modular reconstruction leaves the metric ambiguity, global screen capacity fixes the remaining constant. This logic is compatible with the BW scaling branch but does not depend on the scaling-limit algebra remaining type I.
In the dependency ledger this is the D5\(\to\)D6 handoff: D5 supplies the local Einstein recovery only modulo \(\Lambda g_{ab}\), and D6 closes that same branch only after the global capacity input \(N_{\mathrm{scr}}\) is declared.

\textbf{What this solves vs.\ what it assumes.} The model does \textbf{not} solve the classic ``give me a unitary CFT at future infinity'' problem. It does not aim there. It also does \textbf{not} prove that every refinement-stable MaxEnt branch lands in the BW/canonical static-patch phase. What it does provide is a coherent route to \textbf{patch holography} in which de Sitter static patches are natural:
\begin{enumerate}
\def\labelenumi{\arabic{enumi}.}
\item
the fundamental object is a horizon screen in a static-patch description;
\item
\(\Lambda\) is a capacity parameter tied to finite Hilbert-space dimension, not a locally reconstructible vacuum-energy term;
\item
Einstein-like dynamics emerge up to \(\Lambda g_{ab}\);
\item
on the explicit \textbf{T2} plus extraction-and-rigidity branch, the realized scaling-limit cap modular action is geometric and may be outer on a non-type-I observer algebra.
\end{enumerate}

\textbf{BW-side status.} The Lorentz side stays at the explicit conditionality boundary recorded elsewhere in the manuscript. The internal extension route is a two-object scaffold on the geometric subnet: first a realized scaling-limit cap-pair extraction from transported marginals, then ordered cut-pair rigidity on that realized limit. The pressure point is the first object. Until that realized scaling-limit geometric cap pair exists, the BW/canonical language remains theorem-external shorthand for the \textbf{T2} plus extraction-and-rigidity branch and does not become a fully internalized static-patch theorem.

\textbf{Many observers, one \(\Lambda\).} The philosophical stance of the framework maps naturally onto static-patch intuition: each timelike observer has a horizon and a patch, and there is no operational access to a single global description. The de Sitter parameter \(\Lambda\) is the one thing that all overlap-consistent descriptions share as a global capacity constraint.

\textbf{Summary.} The model gets de Sitter by moving the holographic screen from ``infinity'' to an observer's horizon and by elevating de Sitter entropy, i.e.\ finite screen capacity, to a fundamental input. The usual dS-holography obstacles are precisely the ones avoided by refusing the boundary-at-infinity viewpoint. This is not a claim of a solved dS/CFT dual. It is a static-patch holography program with an explicit local/global split and an explicit extracted-geometric-subnet conditionality boundary.

\begin{center}\rule{0.5\linewidth}{0.5pt}\end{center}
\subsection{Gauge Reconstruction and Standard Model Structure}\label{6-gauge-reconstruction-and-standard-model-structure}

\subsubsection{Edge sector category and gauge group reconstruction}\label{61-edge-sector-category-and-gauge-group-reconstruction}

At any fixed UV cutoff, edge-center completion provides finite tensor categories of transportable edge charges. The local MaxEnt / collar-mixing package established earlier controls only fixed-cutoff recoverability, modular-support localization, and carried error terms on the realized branch. It is logically separate from the question whether transportable edge sectors survive refinement, and it supplies no theorem that the refinement-limit sector category is trivial or nontrivial. A later nontrivial transportable colimit is therefore treated here as additional gauge data rather than as something decided by the mixing estimate itself. Whenever compact gauge reconstruction is invoked, that transportable refinement-directed sector system is exactly the content of \textbf{T6}.

Assume therefore that in the EFT regime there is a directed system \((\mathsf{Sect}_r)_r\) of transportable edge-sector categories with refinement functors, and write

\[
\mathsf{Sect}_\infty := \varinjlim_r \mathsf{Sect}_r,
\]

for the directed colimit retaining the sectors and intertwiners that persist in that transportable system, with objectwise finite-dimensional fibers. The theorem below is neutral about whether \(\mathsf{Sect}_\infty\) is trivial or nontrivial: if \textbf{T6} supplied only the tensor unit, the reconstructed compact group would be the trivial group.

This gives a tensor category \(\mathsf{Sect}_\infty\) of edge charges:

\begin{itemize}
\tightlist
\item
  objects: transportable sector labels that persist in the assumed directed system,
\item
  morphisms: intertwiners between sectors,
\item
  tensor product: fusion by collar concatenation, \(\alpha \otimes \beta = \bigoplus_\gamma N_{\alpha\beta}^{\ \ \gamma}\,\gamma\),
\item
  duals: orientation reversal \(\alpha \leftrightarrow \bar\alpha\),
\item
  symmetric braiding in the EFT regime (no anyonic statistics in 3+1D).
\end{itemize}

Let \(\mathcal F: \mathsf{Sect}_\infty \to \mathsf{Hilb}_{\mathrm{fd}}\) be the bosonic fiber functor that sends each sector to its edge multiplicity space. The fibers are finite-dimensional objectwise even though \(\mathsf{Sect}_\infty\) can have infinitely many simple objects. The theorem begins only once that transportable colimit and its fiber functor are given; those data are not consequences of Axiom 3 alone.

\textbf{Theorem 6.1 (Tannaka/DR reconstruction).} If \(\mathsf{Sect}_\infty\) is a rigid symmetric \(C^*\) tensor category with a faithful bosonic fiber functor \(\mathcal F\), then

\[
G := \mathrm{Aut}_\otimes(\mathcal F)
\]

is a compact group and \(\mathsf{Sect}_\infty \simeq \mathrm{Rep}(G)\). In particular \(G\) is unique up to isomorphism, and this group is a compact subgroup of a product of unitary groups.

\textbf{Proof.} The MaxEnt/local-Gibbs/collar-mixing package is not used here to prove existence of the transportable colimit or to decide whether it is trivial or nontrivial; those data are supplied by the directed-colimit premise \textbf{T6}. Fix a small skeleton of \(\mathsf{Sect}_\infty\). For every object \(X\), a monoidal natural automorphism \(\eta \in \mathrm{Aut}_\otimes(\mathcal F)\) has a unitary component \(\eta_X \in U(\mathcal F(X))\), so

\[
\mathrm{Aut}_\otimes(\mathcal F) \hookrightarrow \prod_X U(\mathcal F(X)), \qquad \eta \mapsto (\eta_X)_X.
\]

For each intertwiner \(f:X\to Y\), naturality gives the closed relation \(\mathcal F(f)\eta_X=\eta_Y\mathcal F(f)\). For each pair \(X,Y\), the monoidal structure isomorphism \(J_{X,Y}\) of \(\mathcal F\) gives the closed relation

\[
\eta_{X\otimes Y}=J_{X,Y}\,(\eta_X\otimes\eta_Y)\,J_{X,Y}^{-1},
\qquad
\eta_{\mathbf 1}=\mathrm{id}_{\mathbb C}.
\]

Hence \(G=\mathrm{Aut}_\otimes(\mathcal F)\) is a closed subgroup of the compact product \(\prod_X U(\mathcal F(X))\), so \(G\) is compact. The remaining DR/Tannaka hypotheses are exactly the stated ones: \(\mathsf{Sect}_\infty\) is rigid, symmetric, and \(C^*\), and \(\mathcal F\) is a faithful bosonic fiber functor with finite-dimensional values. Therefore the Doplicher--Roberts/Tannaka reconstruction theorem applies and gives a symmetric \(C^*\)-tensor equivalence \(\mathsf{Sect}_\infty \simeq \mathrm{Rep}(G)\) \cite{dr89,dr90}. If another compact group \(G'\) gave the same category through a faithful bosonic fiber functor, the induced symmetric tensor equivalence \(\mathrm{Rep}(G)\simeq\mathrm{Rep}(G')\) compatible with the forgetful functors would identify both groups as the automorphism group of the same fiber functor. Thus \(G\cong G'\), so the reconstruction is unique up to isomorphism. QED.

\textbf{Corollary 6.1 (field algebra reconstruction, conditional).} If in the small-region limit the edge sectors are localized and transportable in the DHR sense (i.e., charges can be moved between patches without changing their fusion), then there exists a field algebra \(\mathcal F\) and a compact group \(G\) such that \(\mathcal A = \mathcal F^G\). This is the Doplicher--Roberts reconstruction of local gauge symmetry from sectors. QED.

\textbf{Proposition 6.1a (Transportability from gluing obstruction on the central branch).} On the central-defect branch, the gluing framework identifies the explicit loop-coherence criterion that accompanies DHR transportability in the later ordinary compact-gauge route. In the gluing framework (Section 3), transportability is the statement that charges can be moved between patches without changing fusion rules. The obstruction to path-independent transport is the \v{C}ech-type central cocycle \(z_{ijk}\) from the central-defect subbranch.

Explicitly: the gluing framework gives a gauge-invariant \v{C}ech 2-cocycle class \([z]\) built from triple overlaps (Section 6.6). Transportability holds iff this class vanishes:

\[
\text{DHR transportable} \iff [z] = 0 \iff \text{loop-coherent gluing}.
\]

\textbf{Proof.} Transportability means charges can be moved along any path without affecting the result. In gluing language, this is path-independent parallel transport of edge labels. Lemma 6.12 shows that loop-coherent global gluing exists iff \([z] = 0\). But loop-coherent gluing is exactly path-independent transport, so the equivalence holds. QED.

\textbf{Corollary.} On the central branch, the DHR-transportability condition in Corollary 6.1 is tracked explicitly by the equivalent loop-coherence condition \([z]=0\). This does not make transportability theorem-internal to OPH; it makes the central gluing-side criterion explicit. On the genuinely noncentral branch, Corollary 3.4a gives the parallel fixed-cutoff criterion in terms of the higher-gauge class \(q_\Sigma\), but that branch is handled by crossed-module data rather than by the ordinary DR field-algebra corollary.

\subsubsection{Selecting the SM factors (derived from MAR)}\label{62-selecting-the-sm-factors-derived-from-mar}

Theorem 6.1 yields \emph{some} compact \(G\). Axiom 5 (MAR, Minimal Admissible Realization) uniquely determines which \(G\) is realized.

\begin{center}\rule{0.5\linewidth}{0.5pt}\end{center}

\textbf{Axiom 5 (MAR): Minimal Admissible Realization.} Among all OPH-realizable sector packages \(\mathfrak S\) consisting of the connected Lie gauge-sector image relevant in the low-energy EFT, its admissible light chiral matter content, and one Higgs doublet, and which are (i) loop-coherent / transportable (vanishing relevant obstruction: \([z]=0\) on the central branch or \(q_\Sigma=0\) on the higher-gauge branch), (ii) anomaly-free, (iii) refinement-stable with light chiral matter, (iv) single-Higgs Yukawa-completable with one connected abelian charge factor acting nontrivially on the coupled carrier, (v) intrinsically CP-capable, (vi) weak-sector UV-completable, the realized low-energy package is the lexicographically minimal one under

\[C(\mathfrak S) = (\chi_{\mathrm{cpl}},\; N_{\mathrm{nonab}},\; N_c,\; N_g).\]

The complete formal statement, admissibility definitions, and proof route are given in the rest of Section 6.2.

Here \(\chi_{\mathrm{cpl}}\) is the \textbf{coupled edge capacity}: the dimension of the minimal unitary carrier containing a common irreducible block on which the admissible pseudoreal and complex nonabelian charge types both act nontrivially. This is intentionally stronger than the abstract minimal faithful representation dimension. For \(S(U(3)\times U(2))\), the block-diagonal action on \(\mathbb{C}^3 \oplus \mathbb{C}^2\) is faithful of dimension \(5\), but it is not coupled and therefore does not enter MAR. The object minimized by MAR is the sector package \(\mathfrak S\), not the bare tensor category by itself. MAR is thus an explicit structural-economy selector on the admissible class, not a theorem derived from the earlier axioms.

\textbf{Note on transport obstruction.} The loop-coherent gluing condition is an explicit premise of the gauge derivation. On the central branch this is \([z]=0\) (Proposition 6.1a); on the genuinely noncentral branch it is \(q_\Sigma=0\) (Corollary 3.4a). In either form it ensures that the reconstructed compact group acts as a genuine global gauge symmetry.

\textbf{What MAR derives.} Product structure, minimal sector content, and coupled edge-capacity minimality are consequences of MAR applied to the admissible class:

\begin{itemize}
\tightlist
\item
  \textbf{Product structure}: follows from the minimal coupled carrier \(\mathbb{C}^3 \otimes \mathbb{C}^2\), which enforces commuting color and weak actions.
\item
  \textbf{Minimal sector content}: the pseudoreal doublet and complex triplet are the minimal nonabelian representations satisfying the admissibility conditions, while the connected abelian charge factor is an explicit admissibility input.
\item
  \textbf{Coupled edge-capacity minimality} (formerly S): is the first component of MAR\textquotesingle s complexity vector.
\end{itemize}

\begin{center}\rule{0.5\linewidth}{0.5pt}\end{center}

With MAR stated, the SM derivation proceeds via standard lemmas:

\textbf{Lemma 6.2 (Product factorization implies product group).} If \(\mathsf{Sect} \simeq \mathrm{Rep}(G)\) and \(\mathsf{Sect} \simeq \mathsf{Sect}_1 \boxtimes \mathsf{Sect}_2\), then

\[
G \cong G_1 \times G_2,
\qquad
\mathsf{Sect}_i \simeq \mathrm{Rep}(G_i).
\]

QED.

\textbf{Lemma 6.3 (SU(2) from a pseudoreal doublet).} Let \(H\) be a positive-dimensional compact connected Lie group with a faithful 2D pseudoreal unitary representation \(V\). Then the semisimple part of \(H\) contains an \(SU(2)\) factor acting as the fundamental doublet. Finite or disconnected counterexamples are not part of the theorem package, which only applies the lemma to the identity component on the relevant nonabelian image. QED.

\textbf{Lemma 6.4 (SU(3) from an irreducible triplet).} Let \(H\) be a positive-dimensional compact connected Lie group with a faithful irreducible complex 3D unitary representation \(W\). Then the semisimple image contains an \(SU(3)\) factor acting as the fundamental triplet. Finite or disconnected counterexamples are not part of the theorem package, which only applies the lemma to the identity component on the relevant nonabelian image. QED.

\textbf{Lemma 6.5 (Connected abelian factor criterion).} If the admissible sector package contains a connected abelian charge factor acting nontrivially on the coupled carrier, then the identity component of the abelianized reconstructed group contains a one-torus. Under the single connected abelian-factor admissibility premise, this factor is \(U(1)\). QED.

\textbf{Proposition 6.6 (physical group quotient).} If the realized matter spectrum has hypercharges quantized in sixths, then the kernel acting trivially on all realized sectors is Z\textsubscript{6}, so

\[
G_{\mathrm{phys}} = \frac{\mathrm{SU}(3)\times \mathrm{SU}(2)\times \mathrm{U}(1)}{\mathbb Z_6}.
\]

QED.

\textbf{Proposition 6.6a (SM from MAR).} Under Axiom 5 (MAR):

\begin{itemize}
\tightlist
\item
  Admissibility conditions (iii)--(iv) require both a pseudoreal nonabelian charge type and a complex nonabelian charge type, and include one connected abelian charge factor acting nontrivially on the coupled carrier (Lemma 6.7, Corollary 6.8).
\item
  The minimal faithful pseudoreal representation is the doublet (\ensuremath{\chi} = 2), giving \(SU(2)\). No intrinsically complex 2D irrep qualifies in the connected compact Lie class, because the connected derived image of any irreducible 2D unitary representation lies in \(SU(2)\), so the nonabelian action is pseudoreal up to abelian twist. The first intrinsically complex case is therefore the triplet (\ensuremath{\chi} = 3), giving \(SU(3)\).
\item
  The minimal coupled carrier for both is \(\mathbb{C}^3 \otimes \mathbb{C}^2\), giving coupled edge capacity \(\chi_{\mathrm{cpl}} = 6\).
\item
  The block-diagonal faithful representation \(\mathbb{C}^3 \oplus \mathbb{C}^2\) of \(S(U(3)\times U(2))\) has dimension \(5\), but it is not coupled and therefore is not the MAR minimizer.
\item
  The maximal compact subgroup of \(U(6)\) acting on \(\mathbb{C}^3 \otimes \mathbb{C}^2\) with commuting actions is \((SU(3) \times SU(2) \times U(1))/(\text{finite center})\).
\item
  The commutant of \(SU(3) \times SU(2)\) inside \(U(6)\) is exactly \(U(1)\), so any connected abelian factor acting nontrivially on the coupled carrier is necessarily a single \(U(1)\), and no additional continuous factors appear without increasing \(\chi_{\mathrm{cpl}}\).
\item
  Product structure is not separately assumed: it follows from the tensor product structure of the minimal coupled carrier.
\end{itemize}

Combined with Proposition 6.6 (hypercharges quantized in sixths from the realized spectrum), this yields:

\[
G_{\mathrm{phys}} = \frac{SU(3) \times SU(2) \times U(1)}{Z_6}.
\]

The full proof route is given in the remainder of Section 6.2.

\subsubsection{Refinement stability and unprotected relevant operators}\label{63-refinement-stability-and-unprotected-relevant-operators}

The refinement-stability used later in the gauge branch is contained in the third OPH axiom. Because the same finite local constraint family is preserved across cutoffs, the realized UV states lie in one common finite-dimensional MaxEnt family rather than in a new coupling space at every refinement step. The selected stable branch of this family is therefore not an extra axiom beyond the MaxEnt branch itself.

Concretely, if
\[
\omega_\ell(\lambda)
=
Z_\ell(\lambda)^{-1}
\exp\!\left(
-\sum_x \sum_a \lambda_a O_a(x)-\sum_b \mu_b Q_b
\right),
\]
then any refinement channel \(\Phi_{\ell\to L}\) compatible with Axiom 3 acts on the realized branch by an induced finite-dimensional map
\[
\Phi_{\ell\to L}\bigl(\omega_\ell(\lambda)\bigr)
=
\omega_L\!\bigl(R_{\ell\to L}(\lambda)\bigr).
\]
The resulting refinement-stable branch is simply a trajectory or invariant subset of this finite-dimensional multiplier map. Accordingly, when later sections speak of a refinement-stable directed colimit of sectors, the refinement-stable qualifier means persistence along this MaxEnt branch itself; the remaining extra ingredients are transportability, symmetry, and the bosonic fiber-functor clauses, and the MaxEnt/mixing package does not decide whether the resulting colimit is trivial or nontrivial.

Relevant operators that are neither symmetry-forbidden nor retained in the selected constraint family are precisely the directions that try to push the flow off that stable branch.

\textbf{Lemma 6.7 (refinement-stable MaxEnt branch forbids unprotected relevant operators).} Assume the local finite-constraint MaxEnt/refinement branch of the third OPH axiom. Let \(\mathcal{O}\) be a gauge-invariant Lorentz-scalar relevant deformation in the emergent EFT sense (\(\Delta<4\) in \(3+1\)D), allowed by symmetry and absent from the retained constraint family. Then the selected branch can keep the coupling of \(\mathcal{O}\) at zero only if symmetry forbids that direction or the constraint family explicitly retains it. Otherwise generic refinement induces a nonzero component along that direction and the flow leaves the selected branch. This is a branch-persistence statement, not a universal entropy-ordering theorem for arbitrary off-branch phases.

\textbf{Proof sketch.} Linearize the induced refinement map \(R_{\ell\to L}\) on the finite-dimensional multiplier space around the selected branch. A relevant operator gives an unstable direction with scaling exponent \(y>0\). If \(\mathcal{O}\) is not fixed by symmetry and is not part of the retained constraint family, generic UV mismatch produces a nonzero component along that direction. Because \(y>0\), repeated coarse-graining amplifies the component and pushes the flow off the selected branch unless one fine-tunes it away at every scale or protects it by symmetry or by the declared constraint family. QED.

\textbf{Corollary 6.8 (chirality selector).} A gauge-invariant Dirac mass term is a relevant scalar. If both chiralities exist in conjugate representations, the mass term is allowed and will be generated under refinement unless symmetry-forbidden or explicitly retained as a protected constraint. Therefore keeping light fermions on the selected refinement-stable branch requires chiral matter content or an explicit protecting mechanism. QED.

\subsubsection{Generation number from CP violation and refinement stability (derived from MAR)}\label{64-generation-number-from-cp-violation-and-refinement-stability-derived-from-mar}

Anomaly cancellation is generation-by-generation, so it does not fix the number of generations. The admissibility conditions constrain the window; MAR selects the minimum.

\textbf{Proposition 6.9 (The number of generations is N\_g = 3).} Under (i) intrinsic CP violation in the quark sector, (ii) UV-completability of SU(2)\_L (asymptotic freedom at one loop), (iii) MaxEnt/refinement stability selecting minimal viable spectrum, and (iv) the derived N\_c = 3 from Theorem 6.14, the generation number is

\[
N_g = 3.
\]

\textbf{Inputs.}

\begin{enumerate}
\def\labelenumi{\arabic{enumi}.}
\tightlist
\item
  \textbf{Intrinsic CP violation exists} in the quark sector (empirical fact; also, the framework treats "intrinsic CP violation" as a selector input).
\item
  \textbf{UV-completability proxy}: SU(2)\_L is asymptotically free at one loop in the emergent EFT.
\item
  \textbf{MaxEnt + refinement stability} penalizes unnecessary unfixed flavor structure, selecting the minimal viable spectrum.
\item
  Use the derived N\_c = 3 from Theorem 6.14.
\end{enumerate}

\textbf{Step 1: CP violation lower bound.} The number of physical CP-violating phases in an N\_g \ensuremath{\times} N\_g CKM matrix is:

\[
\#\text{(CP phases)} = \frac{(N_g - 1)(N_g - 2)}{2}.
\]

\begin{itemize}
\tightlist
\item
  For N\_g = 1, 2: this is 0 \ensuremath{\to} \textbf{no intrinsic CP violation possible}.
\item
  For N\_g = 3: this is 1 \ensuremath{\to} \textbf{intrinsic CP violation possible}.
\end{itemize}

So intrinsic CP violation requires:

\[
N_g \ge 3.
\]

\textbf{Step 2: SU(2) asymptotic freedom upper bound.} The one-loop coefficient is:

\[
b_{\mathrm{SU}(2)} = \frac{22}{3} - \frac{1}{3}N_g(N_c + 1) - \frac{1}{6},
\]

where the final \(-1/6\) is the contribution of one complex Higgs doublet. Asymptotic freedom means \(b_{\mathrm{SU}(2)} > 0\), i.e.,

\[
N_g(N_c + 1) < \frac{43}{2}.
\]

With N\_c = 3, we have N\_c + 1 = 4, so:

\[
4 N_g < \frac{43}{2} \quad \Rightarrow \quad N_g \le 5.
\]

Combining: 3 \ensuremath{\leq} N\_g \ensuremath{\leq} 5.

\textbf{Step 3: MAR selection.} Given the allowed window \{3, 4, 5\}, MAR (fourth component of the complexity vector \(C(\mathfrak S)\)) selects the smallest viable choice:

\[
N_g = 3.
\]

QED.

\textbf{Why this is convincing.}

\begin{itemize}
\tightlist
\item
  It predicts a \textbf{single integer}.
\item
  It uses \textbf{two admissibility conditions} (CP violation exists; weak sector is UV-completable) plus MAR\textquotesingle s lexicographic minimality.
\item
  It is not a fit to a continuous number.
\item
  Under the stated gauge-selection hypotheses, this is a derived result, not conditional.
\end{itemize}

\subsubsection{Hilbert-space formulation of gluing data}\label{65-hilbert-space-formulation-of-gluing-data}

Let \{P\_i\} be a good cover of the screen. For each patch, fix a representation

\[
\pi_i: \mathcal{A}_i \to \mathcal B(\mathcal H_i).
\]

For each overlap, choose a unitary intertwiner

\[
U_{ij}: \mathcal H_j \to \mathcal H_i
\]

such that for all \(O \in \mathcal{A}_{ij}\),

\[
\pi_i(O) = U_{ij} \pi_j(O) U_{ij}^\dagger.
\]

Normalize U\_ii = 1 and U\_ji = U\_ij\ensuremath{\dagger}.

\textbf{Lemma 6.10 (centrality on triple overlaps).} On a triple overlap define

\[
\Omega_{ijk} := U_{ij} U_{jk} U_{ki}.
\]

For all \(O \in \mathcal{A}_{ijk}\),

\[
\Omega_{ijk} \pi_i(O) = \pi_i(O) \Omega_{ijk}.
\]

\textbf{Proof.} Conjugation by U\_ki sends \ensuremath{\pi}\_i(O) to \ensuremath{\pi}\_k(O), by U\_jk to \ensuremath{\pi}\_j(O), by U\_ij back to \ensuremath{\pi}\_i(O). Thus conjugation by \ensuremath{\Omega}\_ijk fixes \ensuremath{\pi}\_i(O), so \ensuremath{\Omega}\_ijk commutes with \ensuremath{\pi}\_i(O). QED.

\textbf{Lemma 6.11 (gauge behavior).} If \~{U}\_ij = V\_i U\_ij V\_j\ensuremath{\dagger} with V\_i acting trivially on overlap observables, then

\[
\tilde{\Omega}_{ijk} = V_i \Omega_{ijk} V_i^\dagger.
\]

In particular, if \ensuremath{\Omega}\_ijk is central, its class is gauge invariant. QED.

\subsubsection{Loop obstruction class (central defect)}\label{66-loop-obstruction-class-central-defect}

Assume the defect is central and write

\[
\varphi_{ij} := \mathrm{Ad}(U_{ij}) \quad \text{on } \mathcal{A}_{ij}.
\]

This is the abelian truncation of the full 2-group obstruction in Section 3.4.

Then there exist central unitaries \(z_{ijk}\) such that

\[
\varphi_{ij} \varphi_{jk} \varphi_{ki} = \mathrm{Ad}(z_{ijk}).
\]

\textbf{Theorem 6.12 (loop-coherent gluing iff vanishing obstruction).} The family \(\{z_{ijk}\}\) is a \v{C}ech 2-cocycle, and its class \([z]\) is gauge invariant. On any quadruple overlap \(P_{ijkl}\),

\[
z_{jkl} z_{ikl}^{-1} z_{ijl} z_{ijk}^{-1} = 1.
\]

A loop-coherent global gluing exists iff {[}z{]} = 0.

\textbf{Proof.} Compare two parenthesizations of \ensuremath{\phi}\_ij \ensuremath{\phi}\_jk \ensuremath{\phi}\_kl \ensuremath{\phi}\_li on a quadruple overlap to obtain the cocycle condition above. Gauge changes shift z by a coboundary. If {[}z{]}=0, rephase by a 1-cochain to eliminate defects and obtain path-independent transport. Conversely, loop-coherent gluing implies z\_ijk = 1. QED.

\subsubsection{EFT reduction to anomaly cancellation}\label{67-eft-reduction-to-anomaly-cancellation}

Assume ExtEFT: a low-energy 3+1D chiral gauge theory exists with group \(G\). Then the obstruction class \([z]\) is structurally analogous to the EFT \textquotesingle t Hooft anomaly class and is expected to map to it after a separate anomaly-descent construction. That map is not constructed in this manuscript, so \([z]=0\) is used here as an internal gluing/transportability condition rather than a proved equivalence to EFT anomaly cancellation.

\subsubsection{Hypercharge from anomaly freedom and Yukawas}\label{68-hypercharge-from-anomaly-freedom-and-yukawas}

\textbf{Theorem 6.13 (Hypercharge from anomaly freedom and Yukawas).} Assume gauge group SU(N\_c) \ensuremath{\times} SU(2) \ensuremath{\times} U(1)\_Y and one generation of left-handed Weyl fermions (Q, u\textsuperscript{c}, d\textsuperscript{c}, L, e\textsuperscript{c}), with a Higgs doublet H and Yukawa terms

\[
Q H u^c,\qquad Q H^\dagger d^c,\qquad L H^\dagger e^c.
\]

Then anomaly freedom and Yukawa invariance fix the hypercharges up to an overall normalization, yielding the Standard Model pattern for N\_c = 3.

\textbf{Proof.} Yukawa invariance gives

\[
Y_u = -(Y_Q + Y_H),\quad
Y_d = -Y_Q + Y_H,\quad
Y_e = -Y_L + Y_H.
\]

Anomaly cancellation yields

\[
\begin{aligned}
&SU(2)^2 U(1): && N_c Y_Q + Y_L = 0,\\
&\mathrm{grav}^2 U(1): && 2 N_c Y_Q + N_c Y_u + N_c Y_d + 2 Y_L + Y_e = 0.
\end{aligned}
\]

Solving gives

\[
Y_L = -N_c Y_Q,\quad
Y_H = N_c Y_Q,\quad
Y_u = -(N_c + 1) Y_Q,\quad
Y_d = (N_c - 1) Y_Q,\quad
Y_e = 2 N_c Y_Q.
\]

With these relations, SU(N\_c)\textsuperscript{2}U(1) and U(1)\textsuperscript{3} anomalies vanish automatically. Fixing the normalization by Q = T\textsubscript{3} + Y and Q(\ensuremath{\nu}\_L)=0 gives

\[
Y_Q = \frac{1}{2 N_c}.
\]

For N\_c = 3,

\[
Y_Q = \frac{1}{6},\quad Y_L = -\frac{1}{2},\quad Y_e = 1,\quad
Y_u = -\frac{2}{3},\quad Y_d = \frac{1}{3},\quad Y_H = \frac{1}{2}.
\]

Without Yukawas, the cubic anomaly leaves two discrete branches (Y\_u, Y\_d exchange). Yukawa invariance selects the branch with a single Higgs doublet. QED.

\textbf{Corollary 6.13a (Exact rational hypercharges).} With the derived N\_c = 3, the hypercharge assignments are uniquely fixed to exact rational values:

\[
Y_Q = \tfrac{1}{6}, \quad
Y_L = -\tfrac{1}{2}, \quad
Y_u = -\tfrac{2}{3}, \quad
Y_d = \tfrac{1}{3}, \quad
Y_e = 1, \quad
Y_H = \tfrac{1}{2}.
\]

\textbf{Why this is convincing.}

\begin{itemize}
\tightlist
\item
  These are \textbf{exact rationals}, not approximate numbers.
\item
  Their ratios are fixed by anomaly freedom + Yukawa invariance, and the absolute lattice is fixed by the standard normalization, with no continuous parameters to adjust.
\item
  This high-precision set of numbers strongly constrains the realized matter package and matches observation exactly.
\end{itemize}

\subsubsection{Witten anomaly and the number of colors}\label{69-witten-anomaly-and-the-number-of-colors}

\textbf{Theorem 6.14 (The number of colors is N\_c = 3).} Under the gauge structure SU(N\_c) \ensuremath{\times} SU(2)\_L \ensuremath{\times} U(1)\_Y with one left-handed quark doublet Q per color and one left-handed lepton doublet L per generation, the global SU(2) anomaly (Witten, 1982) requires N\_c to be odd. Under MAR (lexicographic minimization of \(C(\mathfrak S)\), where \(N_c\) is the third component), this yields:

\[
N_c = 3.
\]

\textbf{Inputs.}

\begin{enumerate}
\def\labelenumi{\arabic{enumi}.}
\tightlist
\item
  Low-energy gauge group contains an SU(2)\_L factor and an SU(N\_c) color factor.
\item
  The matter content per generation includes:

  \begin{itemize}
  \tightlist
  \item
    one left-handed quark doublet Q which is an SU(2) doublet and carries color,
  \item
    one left-handed lepton doublet L which is an SU(2) doublet and color singlet.
  \end{itemize}
\item
  \textbf{Witten\textquotesingle s global SU(2) anomaly constraint} (Witten, 1982): the number of left-handed SU(2) doublets must be even.
\item
  \textbf{MAR} (Selection Axiom): among allowed values, the third component of the complexity vector \(C(\mathfrak S)\) is minimized.
\end{enumerate}

\textbf{Proof.} Count SU(2) doublets per generation:

\begin{itemize}
\tightlist
\item
  Quark doublets: N\_c copies (one per color),
\item
  Lepton doublets: 1 copy.
\end{itemize}

Total doublets per generation:

\[
N_c + 1.
\]

Witten anomaly cancellation requires this to be even:

\[
N_c + 1 \equiv 0 \pmod{2} \quad \Rightarrow \quad N_c \text{ is odd}.
\]

The Witten constraint alone allows N\_c \ensuremath{\in} \{1, 3, 5, 7, ...\}. N\_c = 1 fails admissibility (SU(1) is trivial, no complex nonabelian charge type). MAR (input 4) then selects N\_c = 3 as the smallest nontrivial value. QED.

\textbf{Status.} The Witten anomaly derives N\_c odd. The specific value N\_c = 3 follows from MAR applied to the admissible class. Under the stated gauge-selection hypotheses, this is derived, not assumed.

\textbf{Why this is convincing.}

\begin{itemize}
\tightlist
\item
  It predicts a \textbf{single integer} given the minimality selector.
\item
  The odd constraint is independent of continuous parameters, RG running, masses, or Yukawa values.
\item
  It cannot be adjusted without changing the basic notion of electroweak doublets and color replication.
\end{itemize}

\subsubsection{Bond-dimension gatekeeping}\label{610-bond-dimension-gatekeeping}

In tensor-network or code realizations, gauge actions act on edge factors of size \ensuremath{\chi}, so emergent compact gauge groups embed in U(\ensuremath{\chi}). This shows a capacity constraint: accommodating SU(3) color and SU(2) weak factors shows \ensuremath{\chi} \ensuremath{\geq} 6 in the minimal case, consistent with the MAR-derived gauge group.

\subsubsection{Inevitability of photon and graviton}\label{611-inevitability-of-photon-and-graviton}

The model requires photons and gravitons.

\textbf{Photon inevitability chain:}

\begin{enumerate}
\def\labelenumi{\arabic{enumi}.}
\tightlist
\item
  Gauge-as-gluing states that overlap identifications have redundancy forming a local groupoid.
\item
  Theorem 2.3 (edge-center completion) decomposes collar Hilbert spaces into sectors labeled by boundary gauge representations.
\item
  Theorem 6.1 (Tannaka/DR reconstruction) recovers a compact gauge group G from the transportable bosonic edge-sector category defined by these fusion rules.
\item
  Corollary 6.1 (conditional on DHR transportability) reconstructs a field algebra with G as a local gauge symmetry.
\item
  For the Standard Model, G includes U(1)\_em after electroweak symmetry breaking.
\item
  A gauge boson is the quantum of the gauge field. Once U(1)\_em emerges from overlap redundancy, its gauge field exists, and its quantum (the photon) must exist.
\end{enumerate}

The photon is not postulated. It is forced by the axioms through the chain above. The photon mediates the correlations between charged excitations in different patches; it is how the U(1) redundancy structure propagates through the algebra net.

\textbf{Graviton inevitability chain:}

\begin{enumerate}
\def\labelenumi{\arabic{enumi}.}
\tightlist
\item
  Theorem 4.2 (BW\_\{S\textsuperscript{2}\}) shows that on the stated OPH Bisognano--Wichmann refinement branch, and under collar Markov locality plus the controlled refinement-limit premises, modular flow on caps becomes geometric conformal dilation.
\item
  Theorem 4.3 identifies the induced kinematic group as Conf\textsuperscript{+}(S\textsuperscript{2}) \ensuremath{\cong} PSL(2,\ensuremath{\mathbb{C}}) \ensuremath{\cong} SO\textsuperscript{+}(3,1), the Lorentz group.
\item
  Theorem 5.1 shows that, under the stated scaling-limit null-stress and reference-state premises, the condition \ensuremath{\delta}S\_gen = 0 implies the rest-frame first-variation Einstein relation, with overlap consistency upgrading it to the semiclassical Einstein equations in the EFT regime.
\item
  The metric tensor emerges as the compression of modular flow data, and its dynamics are fixed by entanglement equilibrium.
\item
  A dynamical metric in a quantum theory requires a spin-2 quantum field. Its quantum (the graviton) must exist.
\end{enumerate}

The graviton is not postulated. It is forced by the axioms through the chain above. Diffeomorphism invariance emerges because the bulk spacetime description is a compression of screen data; different coordinate descriptions are redundancies in how that compression is presented.

\subsubsection{Quotient-protected charge quantization}\label{612-mass-predictions-from-symmetry}

\textbf{Theorem 6.19 (Charge quantization and no fractional color singlets).} If the global gauge group is

\[
G_{\mathrm{phys}} = \frac{\mathrm{SU}(3) \times \mathrm{SU}(2) \times \mathrm{U}(1)}{Z_6},
\]

as derived in Proposition 6.6, then

\[
\text{All color-singlet states have integer electric charge.}
\]

Equivalently: no stable isolated particles with charges like \ensuremath{\pm}1/3 can exist as color singlets.

\textbf{Proof.} The Z\textsubscript{6} quotient identifies the center elements (e\^{}\{2\ensuremath{\pi}i/3\}, -1, e\^{}\{i\ensuremath{\pi}/3\}) \ensuremath{\in} SU(3) \ensuremath{\times} SU(2) \ensuremath{\times} U(1) with the identity. For a color-singlet state (\ensuremath{\tau} = 0), the SU(3) factor acts trivially. The remaining identification requires the SU(2) \ensuremath{\times} U(1) quantum numbers to satisfy

\[
(-1)^{2j} \cdot e^{i\pi n/3} = 1,
\]

where j is the SU(2) spin and n = 6Y is the integer hypercharge label. This gives n \ensuremath{\equiv} -6j (mod 6), i.e., n \ensuremath{\equiv} 0 (mod 6) for integer j and n \ensuremath{\equiv} 3 (mod 6) for half-integer j. Equivalently: Y is integer when j is integer, and Y is half-integer when j is half-integer.

After electroweak breaking, Q = T\textsubscript{3} + Y. For integer j, T\textsubscript{3} \ensuremath{\in} \ensuremath{\mathbb{Z}} and Y \ensuremath{\in} \ensuremath{\mathbb{Z}}, so Q \ensuremath{\in} \ensuremath{\mathbb{Z}}. For half-integer j, T\textsubscript{3} \ensuremath{\in} \ensuremath{\mathbb{Z}} + 1/2 and Y \ensuremath{\in} \ensuremath{\mathbb{Z}} + 1/2, so Q = (half-integer) + (half-integer) \ensuremath{\in} \ensuremath{\mathbb{Z}}. In both cases, Q \ensuremath{\in} \ensuremath{\mathbb{Z}}. QED.

\textbf{Experimental status}: No fractionally charged color-singlet particles have been observed. Three independent high-precision bounds confirm this:

\begin{enumerate}
\def\labelenumi{\arabic{enumi}.}
\item
  \textbf{Neutrality of matter} (PDG 2024): The proton-electron charge sum satisfies \[|q_p + q_e|/e < 1 \times 10^{-21},\] confirming charge quantization to 21 decimal places.
\item
  \textbf{Fractional charge searches in bulk matter}: Silicone oil drop experiments limit fractionally charged particle abundance to \[(\text{fractionally charged particles})/\text{nucleon} \lesssim 10^{-22}.\]
\item
  \textbf{Collider searches} (CMS, PRL 134, 2025): Exclusions for stable particles with \(q \in [e/3, 0.9e]\) up to masses \(\sim 640\) GeV (95\% CL).
\end{enumerate}

Existing searches are fully consistent with these structural consequences.

These are exact symmetry- or quotient-protected outputs: the empirical content is that violations remain excluded up to available experimental sensitivity.

\subsubsection{Coupling extraction from edge-sector probabilities}\label{613-coupling-extraction-from-edge-sector-probabilities}

The edge-center completion (Theorem 2.3) yields sector probabilities p\_\ensuremath{\alpha} on collar boundaries. These probabilities encode the renormalized gauge coupling through a heat-kernel/Laplacian weighting law.

\textbf{Abelian case (Z\_n).} For a Z\_n gauge theory, the edge sectors are labeled by charge q \ensuremath{\in} \{0, 1, ..., n-1\}. The correct "Casimir" eigenvalue is the Laplacian eigenvalue of the boundary random walk:

\[
\lambda_q = 4 \sin^2\left(\frac{\pi q}{n}\right).
\]

Note: only in the limit n \ensuremath{\to} \ensuremath{\infty} and q \ensuremath{\ll} n does \ensuremath{\lambda}\_q \ensuremath{\approx} (2\ensuremath{\pi}q/n)\textsuperscript{2} \ensuremath{\propto} q\textsuperscript{2}. For finite n, the exact form is essential.

The sector probabilities follow a heat-kernel law:

\[
p_q \propto e^{-t(\mu) \lambda_q},
\]

where t(\ensuremath{\mu}) is the "modular time" parameter encoding the scale. The extraction formula is:

\[
t(\mu) = -\frac{\log(p_q/p_0)}{\lambda_q}, \qquad
g_{\mathrm{ent}}^2(\mu) = \frac{t(\mu)}{2\pi}.
\]

Consistency requires that t extracted from different charges q agrees; this has been verified numerically (see Section 6.14).

\textbf{Electric-center measurement.} The edge sectors are measured using the \emph{electric-center} prescription. For a region A and boundary vertex v \ensuremath{\in} \ensuremath{\partial}A, define the restricted star operator:

\[
Q_v^{(A)} = \prod_{\ell \in \mathrm{star}(v) \cap A} X_\ell^{\pm 1},
\]

where X\_\ensuremath{\ell} is the shift operator on link \ensuremath{\ell}. The sector projectors are:

\[
P_{v,q} = \frac{1}{n} \sum_{m=0}^{n-1} \omega^{-mq} \left(Q_v^{(A)}\right)^m,
\qquad \omega = e^{2\pi i/n},
\]

and the probabilities are p\_\{v,q\} = \ensuremath{\langle}P\_\{v,q\}\ensuremath{\rangle}. This electric-center operator, built from X\textquotesingle s rather than Z\textquotesingle s, correctly captures the boundary gauge charge/flux that labels entanglement edge sectors.

\textbf{Non-abelian generalization.} For SU(N) gauge theories, the edge sectors are labeled by irreducible representations with probabilities:

\[
p_j \propto d_j \, e^{-t(\mu) C_2(j)},
\]

where d\_j is the dimension and C\textsubscript{2}(j) the quadratic Casimir. Extraction:

\[
t(\mu) = -\frac{\log(p_j/p_0)}{C_2(j)}, \qquad
g_{\mathrm{ent}}^2(\mu) = \frac{t(\mu)}{2\pi}.
\]

\textbf{Theoretical derivation.} The fixed-cutoff same-overlap edge law is carried by the
regulated microphysics package once the declared overlap-sector projectors, reversible thermal
branch, and quasi-local local-Gibbs generator of Theorem 2.6 are in hand. On this synthesis
surface, the point is to summarize that fixed-cutoff source package and to state the
refinement/Peter--Weyl lift boundary explicitly rather than to relocate the leaf-level source
surface.

\textbf{Theorem 6.20 (Fixed-cutoff edge-sector law and refinement lift boundary).} Under the OPH
axioms, the fixed-cutoff overlap-gauge realization of Sections 2.3 and 3.2, and the quasi-local
local-Gibbs form supplied by Theorem 2.6, the regulated same-overlap edge-sector probability
distribution satisfies:

\[
p_R = \frac{d_R \, e^{-t \lambda_R}}{\sum_{R'} d_{R'} \, e^{-t \lambda_{R'}}}
\]

where \ensuremath{\lambda}\_R is the Laplacian eigenvalue on the R-isotypic component and t is determined by the collar Gibbs parameter.

\textbf{Proof.}

\emph{Step 1 (Edge Hilbert space).} From the fixed-cutoff overlap-gauge realization, the edge
degrees of freedom at a boundary circle \ensuremath{\Sigma} = \ensuremath{\partial}C live in a
Hilbert space transforming under the gauge group G. Microscopically this is a finite-dimensional
quantum-link edge space on the declared overlap interface; the effective
representation-theoretic description used only for the refinement lift is modeled by \(L^2(G)\).
By the Peter--Weyl theorem~\cite{peterWeyl27}:

\[
L^2(G) \cong \bigoplus_R V_R \otimes V_R^*
\]

where V\_R is the carrier space of irrep R.

\emph{Step 2 (Gauge invariance).} The Gauss law constrains physical states. For an entanglement cut at \ensuremath{\Sigma}, the physical edge Hilbert space decomposes as \ensuremath{\mathcal{H}}\_\{edge\}\^{}\{phys\} = \ensuremath{\oplus}\_R W\_R where W\_R contains states with flux in representation R.

\emph{Step 3 (Natural Hamiltonian).} From the local-Gibbs form, the MaxEnt generator restricted to edge modes takes the form H\_\{edge\} = \ensuremath{\Sigma}\_R h\_R P\_R where P\_R is the projector onto the R-sector. The key claim is that h\_R = \ensuremath{\lambda}\_R.

\emph{Justification:} For a compact simple factor, the group Laplacian \ensuremath{\Delta}\_G = -\ensuremath{\Sigma}\_a (T\^{}a)\textsuperscript{2} is the unique bi-invariant second-order differential operator up to scale. For a product group \(G=\prod_i G_i\), the most general bi-invariant second-order operator is a positive linear combination \(\sum_i c_i \Delta_{G_i}\), with one coefficient per factor. Any other gauge-invariant local choice would require higher derivatives, violating locality. For finite groups, the Cayley graph Laplacian plays the same role: \ensuremath{\lambda}\_R = \textbar S\textbar{} - (1/d\_R) \ensuremath{\Sigma}\_\{s\ensuremath{\in}S\} \ensuremath{\chi}\_R(s).

\emph{Step 4 (MaxEnt selection).} MaxEnt selects the Gibbs state:

\[
\rho_{\mathrm{edge}} = \frac{1}{Z} e^{-t H_{\mathrm{edge}}} = \frac{1}{Z}
\sum_R e^{-t \lambda_R} P_R.
\]

\emph{Step 5 (Sector probabilities).} The probability of sector R is p\_R = Tr(\ensuremath{\rho}\_\{edge\} P\_R). The effective dimension for entanglement is d\_R (not d\_R\textsuperscript{2}) because we trace over one side of the cut. This gives:

\[
p_R = \frac{d_R \, e^{-t \lambda_R}}{Z}.
\]

QED.

\textbf{Why the entropy rank is \(d_R\) (instead of \(d_R^2\)).} The full edge space has dimension \(d_R^2\) in sector \(R\) (from \(V_R \otimes V_R^\ast\)). Entanglement entropy, however, measures correlations \emph{across} the cut. After tracing over one side, the reduced density matrix has effective rank \(d_R\). Mathematically: in the Markov normal form, the edge factor on one side contributes \(\log d_R\) to the entropy.

\textbf{Status.} The fixed-cutoff same-overlap derivation is the theorem-bearing part of the
edge-law package. The quasi-local local-Gibbs generator used here is derived from Theorem 2.6:
if MaxEnt constraints are expectations of finitely many quasi-local operators, the entropy
maximizer is automatically a Gibbs state with a quasi-local generator. The additional
representation-theoretic step summarized here is the compact-group / Peter--Weyl refinement
lift, together with its factorwise Laplacian identification. What remains open above the
fixed-cutoff source package is refinement-stable universality/uniqueness of that lift, not the
existence of the fixed-cutoff same-overlap Casimir law itself.

\textbf{Normalization anchor: 2D Yang-Mills.} The parameter t can be exactly matched to a conventional coupling in 2D Yang-Mills, where the physical Hamiltonian is literally the group Laplacian:

\[
H = \frac{g^2}{2} \Delta_G, \qquad
\Delta_G \chi_R = -C_2(R) \chi_R
\quad \Rightarrow \quad
E_R = \frac{g^2}{2} C_2(R).
\]

Euclidean evolution for "time" A (the area of a cylinder in 2D YM) gives weight(R) \ensuremath{\propto} exp(-A E\_R) = exp(-g\textsuperscript{2} A C\textsubscript{2}(R)/2). Comparing with the heat-kernel expansion K\_t(U) = \ensuremath{\Sigma}\_R d\_R \ensuremath{\chi}\_R(U) e\^{}\{-t C\textsubscript{2}(R)\} yields the exact identification:

\[
t_\mathrm{phys} = \frac{g^2 A}{2} \quad \text{(in 2D YM, no ambiguity).}
\]

This shows that the Laplacian + MaxEnt \ensuremath{\to} heat-kernel structure is exact in a solvable continuum Yang-Mills case. The coefficient in front of C\textsubscript{2} is fixed. In any regime where the edge theory reduces to an effective 2D YM with known "Euclidean thickness" A\_eff:

\[
g^2(\mu) = \frac{2}{A_\mathrm{eff}(\mu)} \cdot \frac{\Delta_R(\mu)}{C_2(R)},
\]

and the RHS must be R-independent. This R-independence is an internal precision consistency test; the formula itself is the normalization map that connects t to the conventional gauge coupling.

\subsubsection{Numerical validation of the heat-kernel law}\label{614-numerical-validation-of-the-heat-kernel-law}

The heat-kernel/Laplacian weighting of edge sectors has been validated in explicit 2D Z\_n gauge models on closed geometries.

\textbf{Model.} A 2\ensuremath{\times}2 periodic lattice gauge theory (8 links) with Z\_n link Hilbert spaces and Hamiltonian:

\[
H = -K \sum_p \mathrm{Re}(B_p) - h \sum_\ell \mathrm{Re}(X_\ell) - \Gamma \sum_v \mathrm{Re}(A_v),
\]

where X\_\ensuremath{\ell} is the Z\_n shift on link \ensuremath{\ell}, B\_p is the oriented plaquette operator (product of Z\textquotesingle s around plaquette p), and A\_v is the oriented star/Gauss operator (outgoing X, incoming X\ensuremath{\dagger}). With K = 1 and \ensuremath{\Gamma} = 5, the ground state satisfies \ensuremath{\langle}A\_v\ensuremath{\rangle} = 1 at all vertices to numerical precision.

\textbf{Region and edge operator.} Region A consists of links whose tail has x = 0 ("half-lattice" cut). At each boundary vertex v, the electric-center edge charge is the restricted star Q\_v\^{}\{(A)\} = \ensuremath{\prod}\_\{\ensuremath{\ell} \ensuremath{\in} star(v) \ensuremath{\cap} A\} X\_\ensuremath{\ell}\^{}\{\ensuremath{\pm}1\}.

\textbf{Results for Z\textsubscript{2}.} With \ensuremath{\lambda}\textsubscript{1} = 4sin\textsuperscript{2}(\ensuremath{\pi}/2) = 4:

{\def\LTcaptype{none} % do not increment counter
\begin{longtable}[]{@{}>{\RaggedLeft\arraybackslash}p{0.192\textwidth}>{\RaggedLeft\arraybackslash}p{0.192\textwidth}>{\RaggedLeft\arraybackslash}p{0.192\textwidth}>{\RaggedLeft\arraybackslash}p{0.192\textwidth}>{\RaggedLeft\arraybackslash}p{0.192\textwidth}@{}}
\toprule\noalign{}
h & p\textsubscript{0} & p\textsubscript{1} & t & g\_ent \\
\midrule\noalign{}
\endhead
\bottomrule\noalign{}
\endlastfoot
0.5 & 0.8266 & 0.1734 & 0.391 & 0.249 \\
1.0 & 0.9612 & 0.0388 & 0.803 & 0.357 \\
2.0 & 0.9917 & 0.0083 & 1.194 & 0.436 \\
\end{longtable}
}

\textbf{Results for Z\textsubscript{3} (overconstrained test).} With \ensuremath{\lambda}\textsubscript{1} = \ensuremath{\lambda}\textsubscript{2} = 4sin\textsuperscript{2}(\ensuremath{\pi}/3) = 3:

{\def\LTcaptype{none} % do not increment counter
\begin{longtable}[]{@{}>{\RaggedLeft\arraybackslash}p{0.112\textwidth}>{\RaggedLeft\arraybackslash}p{0.112\textwidth}>{\RaggedLeft\arraybackslash}p{0.112\textwidth}>{\RaggedLeft\arraybackslash}p{0.112\textwidth}>{\RaggedLeft\arraybackslash}p{0.112\textwidth}>{\RaggedLeft\arraybackslash}p{0.112\textwidth}>{\RaggedLeft\arraybackslash}p{0.112\textwidth}>{\RaggedLeft\arraybackslash}p{0.112\textwidth}@{}}
\toprule\noalign{}
h & p\textsubscript{0} & p\textsubscript{1} & p\textsubscript{2} & t(q=1) & t(q=2) & g\_ent & m\_plaq \\
\midrule\noalign{}
\endhead
\bottomrule\noalign{}
\endlastfoot
0.2 & 0.4395 & 0.2803 & 0.2803 & 0.1500 & 0.1500 & 0.154 & 2.22 \\
0.5 & 0.7509 & 0.1245 & 0.1245 & 0.5989 & 0.5989 & 0.309 & 1.75 \\
1.0 & 0.9606 & 0.0197 & 0.0197 & 1.2956 & 1.2956 & 0.454 & 4.07 \\
1.5 & 0.9851 & 0.0074 & 0.0074 & 1.6288 & 1.6288 & 0.509 & 7.06 \\
2.0 & 0.9921 & 0.0039 & 0.0039 & 1.8440 & 1.8440 & 0.542 & 10.10 \\
\end{longtable}
}

The equality p\textsubscript{1} = p\textsubscript{2} is exact (charge conjugation symmetry in Z\textsubscript{3}). The equality t\_\{q=1\} = t\_\{q=2\} is the crucial \textbf{overconstrained} check: at h = 1.0, extracting t from q = 1 and q = 2 independently gives t\_\{q=1\} \ensuremath{\approx} 1.2956389318579 and t\_\{q=2\} \ensuremath{\approx} 1.2956389318521. The agreement to \textasciitilde10\textsuperscript{-}\textsuperscript{1}\textsuperscript{4} (machine precision) confirms that the edge distribution genuinely follows the heat-kernel/Laplacian form.

\textbf{Region-choice robustness.} At h = 1, the extracted g\_ent is nearly independent of region size:

\begin{itemize}
\tightlist
\item
  2 links (one vertex\textquotesingle s outgoing links): g\_ent \ensuremath{\approx} 0.453
\item
  4 links (half-lattice): g\_ent \ensuremath{\approx} 0.454
\item
  6 links (three vertices): g\_ent \ensuremath{\approx} 0.453
\end{itemize}

This locality confirms that the coupling is dominated by physics near the cut, not global bookkeeping, exactly what is expected if this behaves like a local QFT observable.

\textbf{Results for Z\textsubscript{5} (golden ratio test).} The Z\textsubscript{5} case provides a stringent test because the Laplacian eigenvalues have a distinctive ratio involving the golden ratio \ensuremath{\phi} = (1+\ensuremath{\sqrt{}}5)/2:

\[
\lambda_q = 4\sin^2\left(\frac{\pi q}{5}\right), \qquad
\frac{\lambda_2}{\lambda_1} = \frac{\sin^2(72^\circ)}{\sin^2(36^\circ)} = \phi^2 \approx 2.618.
\]

This ratio distinguishes the Laplacian law from naive alternatives: a linear model (\ensuremath{\lambda}\_q \ensuremath{\propto} q) would predict ratio 2, while a quadratic model (\ensuremath{\lambda}\_q \ensuremath{\propto} q\textsuperscript{2}) would predict ratio 4.

Simulations on a 2\ensuremath{\times}2 torus in the dual/flux basis (125 states in the zero-winding sector) give:

{\def\LTcaptype{none} % do not increment counter
\begin{longtable}[]{@{}>{\RaggedLeft\arraybackslash}p{0.320\textwidth}>{\RaggedLeft\arraybackslash}p{0.320\textwidth}>{\RaggedLeft\arraybackslash}p{0.320\textwidth}@{}}
\toprule\noalign{}
h & Measured ratio ln(p\textsubscript{2}/p\textsubscript{0})/ln(p\textsubscript{1}/p\textsubscript{0}) & Deviation from \ensuremath{\phi}\textsuperscript{2} \\
\midrule\noalign{}
\endhead
\bottomrule\noalign{}
\endlastfoot
0.5 & 2.25 & 14\% \\
1.0 & 2.51 & 4\% \\
2.0 & 2.619 & \textless{} 0.1\% \\
\end{longtable}
}

In the weak-field limit (h \ensuremath{\to} 0, strong magnetic coupling), the simulation converges to the golden ratio squared. This confirms that the vacuum entanglement spectrum encodes the precise geometric structure of the gauge group Laplacian.

\textbf{Significance.} This validates the mathematical law (sector probabilities weighted by Laplacian eigenvalues) in explicit 2D gauge-invariant models with non-flat sector distributions. The Z\textsubscript{3} and Z\textsubscript{5} tests are structurally identical to SU(2)/SU(3): multiple irreps overconstrain the slope, and agreement confirms the mechanism works before jumping to nonabelian groups.

\textbf{Results for S\textsubscript{3} (first nonabelian test).} The abelian tests above use charge-sector projectors that reduce to Fourier modes. For nonabelian groups, the edge-sector projector must be generalized to character projectors:

\[
P_{v,R} = \frac{d_R}{|G|} \sum_{h \in G} \chi_R(h^{-1}) Q_v^{(A)}(h),
\]

where d\_R is the dimension of irrep R, \ensuremath{\chi}\_R is its character, and Q\_v\^{}\{(A)\}(h) is the restricted gauge action at boundary vertex v acting only on links in region A.

For S\textsubscript{3} (the smallest nonabelian group, order 6), there are three irreps: trivial (d=1), sign (d=1), and standard (d=2). The Cayley-graph Laplacian eigenvalues for the transposition generating set are:

\[
\lambda_{\mathrm{triv}} = 0, \qquad \lambda_{\mathrm{sign}} = 6, \qquad
\lambda_{\mathrm{std}} = 3.
\]

\textbf{Exact reduction on one plaquette.} For the single-plaquette model (4 links), imposing Gauss\textquotesingle s law at all vertices means the physical wavefunction depends only on the plaquette holonomy\textquotesingle s conjugacy class. Since S\textsubscript{3} has exactly 3 conjugacy classes, the gauge-invariant Hilbert space is 3-dimensional, spanned by the character states \{\textbar \ensuremath{\chi}\_R\ensuremath{\rangle}\}. In this basis, the edge-sector probabilities are exactly p\_R = \textbar c\_R\textbar\textsuperscript{2} where \textbar \ensuremath{\psi}\textsubscript{0}\ensuremath{\rangle} = \ensuremath{\Sigma}\_R c\_R \textbar \ensuremath{\chi}\_R\ensuremath{\rangle}. This is not an approximation; it is an exact identity for the one-plaquette gauge-invariant sector.

The heat-kernel ansatz predicts p\_R \ensuremath{\propto} d\_R exp(-t \ensuremath{\lambda}\_R). Extracting t independently from the sign and standard irreps provides an overconstrained test: the ratio \ensuremath{\lambda}\_sign/\ensuremath{\lambda}\_std = 6/3 = 2 is a parameter-free prediction. Results from a single-plaquette S\textsubscript{3} lattice gauge model (K=1, \ensuremath{\Gamma}=5):

{\def\LTcaptype{none} % do not increment counter
\begin{longtable}[]{@{}>{\RaggedLeft\arraybackslash}p{0.112\textwidth}>{\RaggedLeft\arraybackslash}p{0.112\textwidth}>{\RaggedLeft\arraybackslash}p{0.112\textwidth}>{\RaggedLeft\arraybackslash}p{0.112\textwidth}>{\RaggedLeft\arraybackslash}p{0.112\textwidth}>{\RaggedLeft\arraybackslash}p{0.112\textwidth}>{\RaggedLeft\arraybackslash}p{0.112\textwidth}>{\RaggedLeft\arraybackslash}p{0.112\textwidth}@{}}
\toprule\noalign{}
h & p\_triv & p\_sign & p\_std & t (sign) & t (std) & \ensuremath{\Delta}t/t & log-ratio \\
\midrule\noalign{}
\endhead
\bottomrule\noalign{}
\endlastfoot
0.5 & 0.909 & 0.0013 & 0.089 & 1.09 & 1.01 & 8.4\% & 2.17 \\
1.0 & 0.980 & 7.5\ensuremath{\times}10\textsuperscript{-}\textsuperscript{5} & 0.020 & 1.58 & 1.54 & 2.8\% & 2.06 \\
2.0 & 0.996 & 4.3\ensuremath{\times}10\textsuperscript{-}\textsuperscript{6} & 0.004 & 2.06 & 2.04 & 1.0\% & 2.02 \\
5.0 & 0.9993 & 1.0\ensuremath{\times}10\textsuperscript{-}\textsuperscript{7} & 0.00066 & 2.68 & 2.67 & 0.3\% & 2.006 \\
12 & 0.9999 & 3.0\ensuremath{\times}10\textsuperscript{-}\textsuperscript{9} & 0.00011 & 3.27 & 3.27 & 0.1\% & 2.002 \\
100 & 1.0000 & 6.1\ensuremath{\times}10\textsuperscript{-}\textsuperscript{1}\textsuperscript{3} & 2.0\ensuremath{\times}10\textsuperscript{-}\textsuperscript{6} & 4.69 & 4.69 & 0.009\% & 2.0002 \\
\end{longtable}
}

The "\ensuremath{\Delta}t/t" column shows the fractional difference (t\_sign - t\_std) / \={t}. The "log-ratio" column shows log(p\_sign/p\textsubscript{0}) / log(p\_std/(2 p\textsubscript{0})), which should equal \ensuremath{\lambda}\_sign/\ensuremath{\lambda}\_std = 2 if the heat-kernel form holds exactly.

As h increases, both diagnostics converge: \ensuremath{\Delta}t/t drops below 10\textsuperscript{-}\textsuperscript{4} and the log-ratio approaches 2.000. This is exactly the expected behavior: finite-size corrections are largest at strong coupling; the heat-kernel form becomes exact as the perturbative regime is approached.

This is the first nonabelian validation of the edge-sector extraction mechanism. The structure (character projectors, Laplacian eigenvalues from the group\textquotesingle s Cayley graph, overconstrained t extraction) is identical to what will be used for SU(2) and SU(3).

\textbf{Parameter-free predictions for SU(2) and SU(3).} The heat-kernel law yields exact, parameter-free ratio predictions that require no scheme matching. Define the "Casimir log-gap":

\[
\Delta_R \equiv \ln\left(\frac{p_0}{d_0}\right) - \ln\left(\frac{p_R}{d_R}\right)
= t \, C_2(R).
\]

Ratios of \ensuremath{\Delta}\_R cancel all unknowns (t, partition function):

\[
\frac{\Delta_{R_1}}{\Delta_{R_2}} = \frac{C_2(R_1)}{C_2(R_2)}
\quad \text{(exact, parameter-free).}
\]

\emph{SU(2) predictions.} Irreps labeled by spin j have d\_j = 2j+1 and C\textsubscript{2}(j) = j(j+1). The framework predicts:

\begin{itemize}
\tightlist
\item
  \ensuremath{\Delta}\textsubscript{1}/\ensuremath{\Delta}\textsubscript{1}/\textsubscript{2} = 2/(3/4) = \textbf{8/3 \ensuremath{\approx} 2.667}
\item
  \ensuremath{\Delta}\textsubscript{3}/\textsubscript{2}/\ensuremath{\Delta}\textsubscript{1}/\textsubscript{2} = (15/4)/(3/4) = \textbf{5}
\item
  \ensuremath{\Delta}\textsubscript{3}/\textsubscript{2}/\ensuremath{\Delta}\textsubscript{1} = (15/4)/2 = \textbf{15/8 = 1.875}
\end{itemize}

\emph{SU(3) predictions.} Irreps labeled by Dynkin indices (p,q) have C\textsubscript{2}(p,q) = (p\textsuperscript{2} + q\textsuperscript{2} + pq + 3p + 3q)/3. Using the fundamental \textbf{3} = (1,0) with C\textsubscript{2} = 4/3 as the reference:

\begin{itemize}
\tightlist
\item
  \ensuremath{\Delta}\textsubscript{8}/\ensuremath{\Delta}\textsubscript{3} = 3/(4/3) = \textbf{9/4 = 2.25}
\item
  \ensuremath{\Delta}\textsubscript{6}/\ensuremath{\Delta}\textsubscript{3} = (10/3)/(4/3) = \textbf{5/2 = 2.5}
\item
  \ensuremath{\Delta}\textsubscript{1}\textsubscript{0}/\ensuremath{\Delta}\textsubscript{3} = 6/(4/3) = \textbf{9/2 = 4.5}
\item
  \ensuremath{\Delta}\textsubscript{1}\textsubscript{5}/\ensuremath{\Delta}\textsubscript{3} = (16/3)/(4/3) = \textbf{4}
\item
  \ensuremath{\Delta}\textsubscript{2}\textsubscript{7}/\ensuremath{\Delta}\textsubscript{3} = 8/(4/3) = \textbf{6}
\end{itemize}

These are the SU(2)/SU(3) analogs of the Z\textsubscript{5} golden-ratio test: exact rational numbers fixed entirely by group theory, with no adjustable parameters.

\textbf{Preliminary SU(3) results.} A one-plaquette SU(3) "quantum link" model (finite truncated irrep basis, n\_max = 12, \ensuremath{\kappa} = 2) has been used to extract t from 14 different irreps simultaneously. The results show internal consistency at the 1-3\% level:

{\def\LTcaptype{none} % do not increment counter
\begin{longtable}[]{@{}>{\RaggedLeft\arraybackslash}p{0.240\textwidth}>{\RaggedLeft\arraybackslash}p{0.240\textwidth}>{\RaggedLeft\arraybackslash}p{0.240\textwidth}>{\RaggedLeft\arraybackslash}p{0.240\textwidth}@{}}
\toprule\noalign{}
bare g\textsuperscript{2} & extracted t (mean\ensuremath{\pm}std) & g\_ent & gap \\
\midrule\noalign{}
\endhead
\bottomrule\noalign{}
\endlastfoot
0.3 & 0.314\ensuremath{\pm}0.0005 & 0.224 & 1.92 \\
0.5 & 0.539\ensuremath{\pm}0.0025 & 0.293 & 1.83 \\
0.8 & 0.896\ensuremath{\pm}0.012 & 0.378 & 1.72 \\
1.0 & 1.144\ensuremath{\pm}0.025 & 0.427 & 1.64 \\
\end{longtable}
}

The standard deviation across irreps provides a built-in error estimate. This is a QCD proton-physics surrogate rather than a full proton calculation; it lacks dynamical quarks and operates on a single plaquette. It nevertheless demonstrates that the nonabelian extraction machinery produces self-consistent outputs without tuning.

\textbf{Extracting the normalization factor A\_eff.} The 2D YM anchor (Section 6.13) gives t = g\textsuperscript{2} A / 2, so the "effective Euclidean thickness" is

\[
A_\mathrm{eff} = \frac{2t}{g^2}.
\]

Computing this from the SU(3) table:

{\def\LTcaptype{none} % do not increment counter
\begin{longtable}[]{@{}>{\RaggedLeft\arraybackslash}p{0.320\textwidth}>{\RaggedLeft\arraybackslash}p{0.320\textwidth}>{\RaggedLeft\arraybackslash}p{0.320\textwidth}@{}}
\toprule\noalign{}
bare g\textsuperscript{2} & extracted t & A\_eff \\
\midrule\noalign{}
\endhead
\bottomrule\noalign{}
\endlastfoot
0.3 & 0.314 & 2.093 \\
0.5 & 0.539 & 2.156 \\
0.8 & 0.896 & 2.240 \\
1.0 & 1.144 & 2.288 \\
\end{longtable}
}

\textbf{Mean}: A\_eff \ensuremath{\approx} 2.19 with point-to-point scatter \textasciitilde4\%.

\textbf{Extrapolation to weak coupling.} The systematic drift in A\_eff shows fitting A\_eff(g\textsuperscript{2}) = A\textsubscript{0} + a \ensuremath{\cdot} g\textsuperscript{2}. A weighted linear fit gives:

\[
A_0 = 2.004 \pm 0.012
\]

with \ensuremath{\chi}\textsuperscript{2}/dof \ensuremath{\approx} 0.09, indicating excellent consistency. This strongly shows that, in this toy UV completion, the "missing normalization" converges to A\_eff \ensuremath{\to} 2 as g\textsuperscript{2} \ensuremath{\to} 0.

The normalization factor behaves like a quasi-constant and extrapolates to a simple value (\ensuremath{\approx} 2) in the weak-coupling limit. This provides a concrete path to absolute coupling predictions: once A\_eff is determined from microphysics, the conversion g\textsuperscript{2} = 2t/A\_eff fixes the gauge coupling without additional free parameters.

\textbf{Internal validation summary.} The heat-kernel law has been validated with increasing precision across multiple gauge groups:

\begin{itemize}
\tightlist
\item
  \textbf{Z\textsubscript{3}}: Overconstrained t extraction (q=1 vs q=2), precision \textasciitilde10\textsuperscript{-}\textsuperscript{1}\textsuperscript{4}
\item
  \textbf{Z\textsubscript{5}}: Golden ratio squared (\ensuremath{\lambda}\textsubscript{2}/\ensuremath{\lambda}\textsubscript{1} = \ensuremath{\phi}\textsuperscript{2}), precision 0.04\%
\item
  \textbf{S\textsubscript{3}}: Casimir log-ratio (\ensuremath{\lambda}\_sign/\ensuremath{\lambda}\_std = 2), precision 0.01\%
\item
  \textbf{SU(3)}: 14-irrep simultaneous extraction, precision 1-3\%
\end{itemize}

The Z\textsubscript{3} test achieves machine precision because it is exactly overconstrained. The Z\textsubscript{5} and S\textsubscript{3} tests converge to their predicted ratios as coupling decreases. This provides strong internal validation of the mechanism "MaxEnt + Laplacian =\textgreater{} heat-kernel sector weights" before applying it to physical gauge groups.

\subsubsection{IBM Quantum Cloud hardware benchmarks}\label{614a-ibm-quantum-cloud-hardware-benchmarks}

The lattice calculations above are internal numerical validations. The IBM Quantum Cloud benchmarks add a hardware check of whether OPH-motivated reduced-sector structures survive preparation and readout on real superconducting-qubit devices. These runs are small-sector consistency benchmarks rather than a standalone confirmation of the full framework. Their value is that they probe the same recoverability and heat-kernel observables on physical hardware and in simulation.

Detailed circuits, representative raw outputs, and rerun instructions are public in the project write-up \href{https://github.com/FloatingPragma/observer-patch-holography/blob/main/extra/IBM_QUANTUM_CLOUD.md}{\texttt{extra/IBM\_QUANTUM\_CLOUD.md}} and the associated code/data bundle \href{https://github.com/FloatingPragma/observer-patch-holography/tree/main/code/ibm_quantum_cloud}{\texttt{code/ibm\_quantum\_cloud/}}.

\begin{itemize}
\tightlist
\item
  \textbf{Stage 1 (Markov/recoverability benchmark).} On \texttt{ibm\_marrakesh} and \texttt{ibm\_fez}, the structured state reconstructs below both controls in conditional mutual information and above both controls in Petz fidelity. On \texttt{ibm\_marrakesh}, the observed ordering is \(0.2309 < 0.3890 < 0.9474\) for CMI and \(0.9297 > 0.8649 > 0.6066\) for Petz fidelity; on \texttt{ibm\_fez}, it is \(0.1498 < 0.4992 < 0.9166\) and \(0.9479 > 0.8117 > 0.6449\). This is the OPH-favored direction on two independent real backends.
\item
  \textbf{Z\textsubscript{3} hardware sanity check.} The overconstrained \(Z_3\) extraction passes cleanly on hardware: across prepared \(t = 0.30, 0.60, 0.90\), the mean extracted \(t\) remains within about \(0.02\) of the target, the two independent extractions agree to within \(0.0142\), \(0.0034\), and \(0.0043\), and leakage stays below \(0.1\%\). This shows that the reduced-sector preparation and readout path is internally coherent on-device before sharper ratio claims are interpreted.
\item
  \textbf{Z\textsubscript{5} exact-ratio test.} The parameter-free OPH target is \(\Delta_2/\Delta_1 = \varphi^2 \approx 2.618\). Across repeated sweeps on \texttt{ibm\_marrakesh} and one replication on \texttt{ibm\_fez}, the best \texttt{ibm\_marrakesh} points land within \(0.8\%\), \(1.2\%\), and \(1.5\%\) of the target, while a focused high-shot rerun at \(t = 0.90\) lands within \(2.6\%\). The \texttt{ibm\_fez} replication is noisier but remains in the same neighborhood rather than producing a clean contradictory ratio.
\item
  \textbf{S\textsubscript{3} nonabelian ratio test.} The first nonabelian hardware run revealed a real layout-dependent bias. After reversing the qubit layout, the mitigated ratio moves from \(1.8724\) to \(2.0299\) against the OPH target \(2.0000\), and a direct confirmation run returns \(2.0657\). The nonabelian target reappears once the hardware mapping is corrected, indicating an identifiable device/layout effect instead of a structural failure of the OPH prediction.
\end{itemize}

These IBM runs are best read as hardware benchmarks for the local OPH edge-sector picture. They show that the OPH-favored recoverability ordering survives on two real backends, the abelian exact-ratio tests land near the predicted parameter-free values, and the first nonabelian test returns to the expected ratio once a diagnosed layout bias is corrected. They add real-device evidence to the simulation-based checks, but they do not validate the full framework by themselves.

\begin{center}\rule{0.5\linewidth}{0.5pt}\end{center}
\subsection{Status, Tests, and Open Problems}\label{8-status-tests-and-open-problems}

\subsubsection{Classification of results and remaining dependence}\label{81-classification-of-results}

This section states which node of the program has been proved, which standard mathematics is imported, and which extra physical premises or external inputs remain load-bearing. The documentary node labels D1--D12 match the ledger in Section~\ref{16-summary-complete-axiom-set}. The three-phase boundary is
\[
\text{Phase I}=(D1\text{--}D5)\cup(D7\text{--}D9),\qquad
\text{Phase II}=D6\cup D10,\qquad
\text{Phase III}=D12.
\]

\begingroup
\scriptsize
{\def\LTcaptype{none}
\setlength{\tabcolsep}{2pt}
\begin{longtable}[]{@{}>{\RaggedRight\arraybackslash}p{0.07\textwidth}>{\RaggedRight\arraybackslash}p{0.22\textwidth}>{\RaggedRight\arraybackslash}p{0.15\textwidth}>{\RaggedRight\arraybackslash}p{0.20\textwidth}>{\RaggedRight\arraybackslash}p{0.17\textwidth}@{}}
\toprule\noalign{}
Node & Representative outputs in this manuscript & Standard mathematics used & Extra physical premises / external inputs required & Claim tier \\
\midrule\noalign{}
\endhead
\bottomrule\noalign{}
\endlastfoot
\textbf{D1} & Theorem 3.1 schedule-independent normal form on the physical quotient together with observable-level confluence on the declared fixed-cutoff physical algebras for the recovery-derived repair relation & local-to-global Lyapunov descent on the inconsistency potential; local-diamond completion from overlap-associative center-sector / higher-gauge union-collar gluing; Newman's lemma & repair completeness; stated Petz support/CPTP control where that branch is used & Phase I structural theorem \\
\textbf{D2} & EC decomposition, exact/approximate Markov collar, and generalized-entropy split & HJPW; Petz; Fawzi--Renner & explicit recoverability control when exact identities are used; coarse-grained \(A/(4G)\) dictionary & Phase I structural lemma/proposition layer \\
\textbf{D3} & Theorems 4.2--4.3 Lorentz branch & Bisognano--Wichmann modular geometry; conformal-group classification & rotationally invariant constraints together with T2 and the theorem-local geometric cap-pair extraction and ordered cut-pair rigidity hypotheses of Theorem 4.2 & Phase I conditional scaling-limit theorem \\
\textbf{D4} & Section 5.2 null modular bridge & Borchers--Wiesbrock; distribution theory for half-line differentiation & the theorem-local inherited-strip decomposition and exact-or-controlled Markov hypotheses of Section 5.2, plus the same-half-line local modular-Hamiltonian input; the endpoint/propagation controls are internal to the local finite-constraint MaxEnt branch, and bounded-interval formulas still require the separate interval-preserving projective branch & Phase I conditional bridge theorem \\
\textbf{D5} & Theorem 5.1 rest-frame Einstein relation plus the conditional tensor upgrade of Corollary 5.1 & Jacobson small-ball mathematics; null-to-tensor reconstruction modulo a metric term & fixed-cap stationarity, the locally Lorentzian \(d=4\) scaling regime, the separate interval-preserving projective branch used in the small-ball step, and the all-directions/all-reference-states premise used in the tensor upgrade & Phase I conditional scaling-limit theorem/corollary \\
\textbf{D6} & capacity/\(\Lambda\) statements, plus a legacy capacity-level neutrino side estimate & de Sitter entropy relation and dimensional analysis & external input \(N_{\mathrm{scr}}\) and the screen-capacity identification & Phase II input-dependent corollary / estimate \\
\textbf{D7} & Theorem 6.1 compact gauge reconstruction & Doplicher--Roberts / Tannaka reconstruction & bosonic symmetric transportable colimit and fiber-functor premises on the ordinary or central-defect branch; the genuinely noncentral fixed-cutoff branch is handled separately by crossed-module data, and the refinement-stable qualifier is internal to the local finite-constraint MaxEnt/refinement branch & Phase I conditional structural theorem \\
\textbf{D8} & product gauge structure up to finite quotient under MAR & compact Lie representation classification; Schur's lemma & the same ordinary or central-defect bosonic branch as D7, together with MAR, a connected admissible class, one connected abelian factor, and faithful action on the minimal coupled carrier & Phase I realized-branch theorem \\
\textbf{D9} & realized \(\SU(3)\times\SU(2)\times\U(1)/\mathbb Z_6\), exact hypercharges, the counting chain \(N_g=3\) then \(N_c=3\), and no gauge-mediated proton decay as a product-group corollary & anomaly algebra; Witten anomaly; CKM CP counting & realized one-generation/one-Higgs package and MAR admissibility premises & Phase I realized-branch theorem/corollary chain \\
\textbf{D10} & forward gauge-coupling closure and target-free source-only electroweak repair theorem of the D10 calibration branch & RG running and matching conventions & external input \(P\), pixel closure, the source-only forward transmutation solve \(\mathcal F(\alpha_U;P)=0\), the fixed-cutoff edge heat-kernel / Casimir theorem on the microphysics surface together with the compact-group / Peter--Weyl lift, and the printed threshold conventions & Phase II calibration sector \\
\textbf{D12} & charged-lepton centered-readback / centered-to-affine frontier, textures, dark-sector, baryogenesis, spectroscopy, the string/worldsheet lane beyond the fixed-cutoff heat-kernel / Yang--Mills form on the declared compact-gauge branch, and other downstream continuations & branch-specific EFT and phenomenological manipulations & explicit additional ans\"atze beyond D6, D9, or D10, including any controlled large-\(N_{\mathrm{edge}}\) regime and any extra worldsheet/CFT premises needed for the conjectural critical-superstring lift & Phase III phenomenological continuation / program branch \\
\end{longtable}
}
\endgroup

The ledger implies the following reading rules.

\begin{enumerate}
\def\labelenumi{\arabic{enumi}.}
\item
Phase I, the recovered core carried by this manuscript, is D1--D5 together with D7--D9. These are the structural or realized-branch claims defended at theorem/corollary level once all stated premises are included.
\item
Phase II contains the separate input-dependent capacity corollary D6 and the D10 calibration sector. Failure there does not by itself erase Phase I.
\item
Phase III contains D12 continuations only. Nothing in D12 is theorem-level unless its extra ans\"atze are replaced by proved premises.
\item \textbf{Existence / strange-loop narratives}: self-simulation or strange-loop closure stories are interpretive epilogues only. The internally available ingredient is a theorem-usable OPH state-and-law habitat built from overlap-consistent patch states together with finite Axiom-3 law coordinates. The recovered core does \emph{not} derive an existence theorem, an OPH-internal closure map on an invariant admissible observer-supporting sector, or any uniqueness/stability result for such a map.
\end{enumerate}

The corresponding low-energy effective-action form is
\[
\mathcal L_{\mathrm{eff}}^{\mathrm{OPH}}
\approx
\sqrt{-g}\left[
\frac{1}{16\pi G}(R-2\Lambda)
+
\mathcal L_{\mathrm{SM}}^{\mathrm{realized\ branch}}
\right]
+
\sum_i \frac{c_i}{M_*^{\Delta_i-4}}\mathcal O_i.
\]
Here \(\mathcal L_{\mathrm{SM}}^{\mathrm{realized\ branch}}\) denotes the realized \(\SU(3)\times\SU(2)\times\U(1)/\mathbb Z_6\) sector with the exact hypercharge lattice and the realized counting chain \(N_g=3\) then \(N_c=3\). Phase I secures the conditional Einstein branch and the realized Standard Model structural chain, D6 supplies \(\Lambda\) once the screen-capacity identification is adopted, and the higher-dimension operators absorb UV-sensitive or downstream corrections not fixed by the recovered core itself.

Symmetry-protected exact zeros such as the masslessness of the photon, gluons, and graviton, together with the absence of gauge-mediated proton decay, are corollaries inside Phase I. They do not define separate axiom-only nodes.
\subsubsection{Structural assessment}\label{82-structural-assessment}

\begin{itemize}
\tightlist
\item
  \textbf{Dynamics}: The GR chain requires modular covariance plus the fixed-cutoff null-surface modular bridge. Section 5.2 carries that bridge through the derived positive null-translation stage and the exact half-line generator/charge identification on the declared geometric scaling branch. The remaining inputs are modular covariance, the bounded-interval projective transport branch, and the tensor upgrade modulo the null-invisible metric term.
\item
  \textbf{Gauge structure}: The gauge group is reconstructed from the transportable bosonic sector package, and the anomaly/gluing link is precise. Axiom 5 (MAR) fixes the SM factors; the CP/UV admissibility window and MAR give \(N_g = 3\), and Witten's parity constraint with that realized generation count gives \(N_c = 3\). On the central branch, Proposition 6.1a identifies the explicit loop-coherence criterion \([z]=0\) that accompanies DHR transportability in the ordinary compact-group route. On the genuinely noncentral branch, Corollary 3.4a gives the parallel fixed-cutoff criterion \(q_\Sigma=0\), but that branch is handled separately by crossed-module data rather than by the ordinary DR field-algebra route. The chirality selector still has to be justified in explicit models.
\item
  \textbf{Microscopic theory}: Quantum link models (Section 2.6) provide explicit fixed-cutoff realized presentations and give EC + Markov collars automatically; they are not proofs of a unique microscopic completion. Open items are (i) a microscopic derivation of Recoverable Generalized Entropy, and (ii) the continuum lift on the extracted geometric subnet, namely \textbf{T2} together with geometric cap-pair extraction and ordered cut-pair rigidity.
\item
  \textbf{Loop gluing beyond central defect}: At fixed cutoff the genuinely noncentral branch is closed by the crossed-module higher-gauge EC and transport package; what remains open is the refinement-limit transportable-sector lift together with quantitative matching to EFT anomalies.
\end{itemize}

\subsubsection{Testable Items and Conditional Phenomenological Branches}\label{83-novel-testable-predictions}

The framework makes several directly testable statements, together with weaker continuation-level templates that can also be confronted with data:

\textbf{GW horizon spectroscopy template (Section 5.11; continuation-level).} If the extra discrete-horizon premises of Section 5.11 hold, the log-integer area spectrum gives discrete resonant frequencies for Kerr black hole horizons:

f\_\{k,m\}(M,\ensuremath{\chi}) = (m \ensuremath{\Omega}\_H)/(2\ensuremath{\pi}) + (c\textsuperscript{3} g(\ensuremath{\chi}))/(16\ensuremath{\pi}\textsuperscript{2} GM) \ensuremath{\cdot} ln(k) for k = 2, 3, 4, ...

After rescaling by remnant parameters, all events should stack at universal coordinates x\_k = ln(k)/(8\ensuremath{\pi}). This is checkable with public LIGO/Virgo data as a continuation-level test of that extra branch. Absence of coherent stacking would challenge the discrete-horizon continuation, not the recovered-core package by itself.

\textbf{Discrete Hawking template (Section 5.11; continuation-level).} If the same discrete-horizon branch is realized, primordial black holes in the final evaporation stage should show comb structure at E\_k/E\textsubscript{2} = ln(k)/ln(2). Current PBH burst searches (Fermi, H.E.S.S.) can constrain this with dedicated template analysis.

\textbf{Casimir ratio precision (Section 8.1).} Future lattice measurements of SU(3) edge-sector probabilities should confirm \ensuremath{\Delta}\textsubscript{8}/\ensuremath{\Delta}\textsubscript{3} = 9/4 exactly, not 2.67 (dimension-only) or 5.06 (Casimir-squared). The full set of parameter-free SU(3) ratio predictions is:

\begin{itemize}
\tightlist
\item
  \ensuremath{\Delta}\textsubscript{8}/\ensuremath{\Delta}\textsubscript{3} = 9/4 = 2.25
\item
  \ensuremath{\Delta}\textsubscript{6}/\ensuremath{\Delta}\textsubscript{3} = 5/2 = 2.5
\item
  \ensuremath{\Delta}\textsubscript{1}\textsubscript{0}/\ensuremath{\Delta}\textsubscript{3} = 9/2 = 4.5
\item
  \ensuremath{\Delta}\textsubscript{1}\textsubscript{5}/\ensuremath{\Delta}\textsubscript{3} = 4
\item
  \ensuremath{\Delta}\textsubscript{2}\textsubscript{7}/\ensuremath{\Delta}\textsubscript{3} = 6
\end{itemize}

These exact rationals are fixed entirely by group theory (Casimir eigenvalue ratios), with no adjustable parameters. Any deviation identifies a measurement contradiction with the heat-kernel edge-sector mechanism.

\textbf{Conditional Z$_6$ center-label entropy scale (Section 6.18).} Under the uniform sixfold center-label ensemble used in the flavor continuation, the associated entropy scale is \(\log_2 6 \approx 2.585\) bits. This is a phenomenological ansatz tied to the quotient label set, not a theorem-level direct observable of the quotient alone.

\textbf{Black-hole spectroscopy bookkeeping within the continuation template (Section 5.11).} If one adopts the extra discrete-horizon selection rule, the semiclassical linewidth estimate, and the greybody matching used in Section 5.11, then the template carries secondary structure:

\begin{enumerate}
\def\labelenumi{\arabic{enumi}.}
\item
  \emph{Universal energy ratios}: E\_k/E\_2 = ln(k)/ln(2) exactly within the assumed log-integer rule. For example, E\_3/E\_2 = ln(3)/ln(2) \ensuremath{\approx} 1.585. This arithmetic pattern of ratios distinguishes the template from other "quantized area" proposals that have different functional forms or free spacing parameters.
\item
  \emph{Mass-independent fractional linewidth estimate}: The continuation-level linewidth estimate \(\Gamma/\Delta E_k \approx 3\text{--}5\%\) is approximately independent of black hole mass.
\item
  \emph{Fixed weight hierarchy}: Within the same template, line weights follow a (k-1)/k pattern from detailed balance in the log-integer transition rule, on top of the GR greybody envelope.
\end{enumerate}

\textbf{Inequality bounds on GR deviations (Section 5.8).} The modular additivity defect satisfies the exact identity \ensuremath{\langle}\ensuremath{\Delta}K\ensuremath{\rangle} = -I(A:D\textbar B), where I(A:D\textbar B) is the conditional mutual information. Under the Markov/mixing assumptions, this defect is exponentially small in collar thickness:

\textbar\ensuremath{\langle}\ensuremath{\Delta}K\ensuremath{\rangle}\textbar{} \ensuremath{\leq} 2\textbar A\textbar{} \ensuremath{\cdot} \ensuremath{\eta}\^{}\{w/\ensuremath{\xi}\}

This propagates into an explicit upper bound on how far the Einstein equation can deviate from GR in regimes where the emergence proof applies. Unlike typical beyond-GR frameworks that postulate corrections, OPH provides a quantitative ceiling: given the information-theoretic primitives, corrections decay exponentially with collar width. This "UV ignorance \ensuremath{\to} rigorous inequality" structure is distinctive.

\paragraph{Shared excitation dictionary benchmark (Section 6.21).} The same-family continuations are organized by one common excitation dictionary instead of a primitive one-number suppression law. The proof-facing data are the same-label gap/overlap scalars
\[
(g_e,\omega_e),\qquad
d_e=1-\omega_e,\qquad
q_e=\sqrt{g_e d_e},\qquad
\eta_e=\log q_e-\frac13\sum_f \log q_f,
\]
equivalently the centered eta-class \([\eta_e]\) or the normalized weights
\[
\mu_e:=\frac{e^{\eta_e}}{\frac13\sum_f e^{\eta_f}}.
\]
Let \(s_e\) denote the realized sector-charge step carried by the arrow \(e\).
The internal integer data are the excitation multiplicities
\[
m_{\gamma,e}\in\mathbb Z_{\ge0}
\]
with which a given readout path \(\gamma\) traverses the realized same-label arrows, equivalently the integers satisfying
\[
s_\gamma=\sum_e m_{\gamma,e}s_e,
\]
and the corresponding excitation action
\[
A_\gamma=-\sum_e m_{\gamma,e}\log q_e.
\]
Old ``base-6 exponents'' are only the compression
\[
n_\gamma^{(6)}:=A_\gamma/\log 6,
\]
which becomes an actual defect count only after the extra sixfold-uniform collapse compatible with the realized \(\mathbb Z_6\) quotient. The intrinsic neutrino lane factors directly through the same-label scalar certificate, the quark mass-side continuation compresses to \((\Delta_{ud}^{\mathrm{overlap}},\eta_Q^{\mathrm{centered}})\) and the emitted ray \(\mathcal R_{D12}^{ud}\), and the charged lane fixes centered logs once a charged source pair exists while the common-shift quotient is already closed negatively. The remaining blocker structure therefore sits above the dictionary: \(\widehat C_e^{\mathrm{cand}}\to\widehat C_e\) and then the post-promotion lift whose descended scalar is \(\mu_{\mathrm{phys}}(Y_e)\), with the physical identity-mode equalizer beneath that scalar, for charged leptons; the maximal theorem-emitted quark package on the present ledger is the D12 mass ray, the negative selector closure \(\sigma_{\mathrm{ref}}\), and the restricted-scope affine mean package with \(g_{\mathrm{ch}}=0.9231656602589082\) on \texttt{shared\_budget\_only} and \((g_u,g_d)=(0.7797392875757557,0.12172551081512113)\) on \texttt{current\_family\_only}; the stronger physical closure objects \(\Theta_{ud}^{\mathrm{mass}},\Theta_{ud}^{\mathrm{phys}},\Theta_{ud}^{\mathrm{abs}}\) are absent on that ledger, and the exact minimal extension above that theorem-emitted package is the triple \((H_{\mathrm{mass}},H_{\mathrm{phys}},H_{\mathrm{abs}})\) with \(H_{\mathrm{mass}}:\ell_{ud}=\log(c_d/c_u)\), \(H_{\mathrm{phys}}:s_{ud}^{\mathrm{phys}}:\mathfrak M_{ud}^{\mathrm{CR,phys}}\to\Sigma_{ud}^{\mathrm{phys}}\), and \(H_{\mathrm{abs}}:A_q^{\mathrm{phys}}:\Sigma_{ud}^{\mathrm{phys}}\to\mathbb R\); while the weighted-cycle neutrino lane fixes the bridge-rigidity invariant \(C_\nu=G^2Q S^{-1/2}=0.9994295999075177\) above the emitted proxy \(P_\nu=6.699825740519345\), with \(B_\nu=P_\nu C_\nu=6.696004159297337\) retained only as the paper-facing amplitude parameterization. The scalar benchmark \(\varepsilon=1/6\) survives only as the extra uniform sixfold collapse of this richer dictionary, so exponents \(-\log(y_f)/\log 6\) are compare-only summaries of excitation action, not recovered-core outputs or independent flavor inputs.

\paragraph{No gauge-mediated proton decay (Section 6.11).} If the gauge group is genuinely a product (from sector factorization), coupling unification is geometric (shared edge diffusion parameter) rather than simple-group embedding. This predicts unification-like coupling relations \emph{without} GUT leptoquark bosons, hence no gauge-mediated proton decay. The combination "coupling unification + no proton decay" is a crisp discriminator against classic GUT predictions.

\subsubsection{What is not predicted (gaps)}\label{84-what-is-not-predicted-gaps}

Section~\ref{81-classification-of-results} fixes the three-phase boundary. The list below collects the items that remain outside Phase I even when they are discussed elsewhere in the manuscript.

\textbf{Not predicted by the recovered core.}
\begin{itemize}
\item \textbf{Flavor and Yukawa data}: Yukawa matrices, flavor hierarchies, charged-lepton continuation ans\"atze beyond exact centered readback, quark-texture exponents, CKM/PMNS closure, and related downstream flavor fits are not Phase-I outputs. The common proof-facing family base is the same-label excitation dictionary
\[
(g_e,\omega_e)\longmapsto d_e=1-\omega_e,\qquad q_e=\sqrt{g_e d_e},\qquad [\eta_e].
\]
Its internal integer data are the realized traversal multiplicities
\[
m_{\gamma,e}\in\mathbb Z_{\ge0}
\]
with
\[
s_\gamma=\sum_e m_{\gamma,e}s_e.
\]
The associated excitation action is
\[
A_\gamma=-\sum_e m_{\gamma,e}\log q_e,
\]
from which any compare-only base-6 compression is formed.
The lane-specific readout boundary above that dictionary is explicit: on the emitted local quark surface the same-label selector is fixed to \(\sigma_{\mathrm{ref}}\), which leaves the physical CKM-shell no-go intact on the selected D12 sheet, and an explicit same-family underdetermination theorem says that the mass side emits only the one-parameter ray \(D12_{ud}^{\mathrm{mass}}\). The maximal theorem-emitted quark package on the present ledger is therefore the D12 mass ray, the negative selector closure \(\sigma_{\mathrm{ref}}\), and the restricted-scope affine mean package with \(g_{\mathrm{ch}}=0.9231656602589082\) on \texttt{shared\_budget\_only} and \((g_u,g_d)=(0.7797392875757557,0.12172551081512113)\) on \texttt{current\_family\_only}. The stronger physical closure objects \(\Theta_{ud}^{\mathrm{mass}}, \Theta_{ud}^{\mathrm{phys}}, \Theta_{ud}^{\mathrm{abs}}\) are absent on that ledger. The exact minimal extension above that theorem-emitted package is \(H_{\mathrm{mass}}:\ell_{ud}=\log(c_d/c_u)\), \(H_{\mathrm{phys}}:s_{ud}^{\mathrm{phys}}:\mathfrak M_{ud}^{\mathrm{CR,phys}}\to\Sigma_{ud}^{\mathrm{phys}}\), and \(H_{\mathrm{abs}}:A_q^{\mathrm{phys}}:\Sigma_{ud}^{\mathrm{phys}}\to\mathbb R\). With that triple adjoined, \(H_{\mathrm{mass}}\) emits \(\Theta_{ud}^{\mathrm{mass}}=\text{\texttt{quark\_d12\_t1\_value\_law}}\) with \(t_1=\tfrac56\ell_{ud}\), \(\Delta_{ud}^{\mathrm{overlap}}=\tfrac16\ell_{ud}=t_1/5\), \(\eta_Q^{\mathrm{centered}}=-\tfrac{1-x_2^2}{27}t_1\), \(\kappa_Q=-t_1/54\), and the odd source package functorially, with \texttt{intrinsic\_scale\_law\_D12} as the derived wrapper; \(H_{\mathrm{phys}}\) emits the sector-attached same-label left-handed lift to the physical CKM shell; and \(H_{\mathrm{abs}}\) emits the target-free absolute sector readout on that physical sheet, promoting the affine mean split \(g_u=g_{\mathrm{ch}}\exp(-(A_{ud}\sigma_{\mathrm{seed}}^{ud}-B_{ud}\eta_{ud}))\), \(g_d=g_{\mathrm{ch}}\exp(-(A_{ud}\sigma_{\mathrm{seed}}^{ud}+B_{ud}\eta_{ud}))\). The exact same-family selected-sheet witness hits the running quark reference sextet on \texttt{current\_family\_only}, and a separate continuation-only D12 internal backread sidecar fixes the mass-side scalar package numerically on its own surface, but neither surface replaces the public theorem frontier or removes the wrong-sheet CKM boundary. The charged-lepton lane carries exact centered readback and a closed common-shift no-go, then waits on promotion of \(\widehat C_e^{\mathrm{cand}}\) and the post-promotion lift whose descended scalar is \(\mu_{\mathrm{phys}}(Y_e)\), with the physical identity-mode equalizer beneath that scalar, from which the uncentered lift, determinant-line section, and affine anchor \(A_{\mathrm{ch}}\) follow canonically. The weighted-cycle neutrino theorem boundary is carried by the bridge-rigidity invariant \(C_\nu=G^2Q S^{-1/2}=0.9994295999075177\) with \(G=\mathrm{sum\_gap}\), \(Q=\mathrm{prod\_qbar}\), and \(S=\mathrm{solar\_response\_over\_mstar}\), above the emitted proxy \(P_\nu=6.699825740519345\); \(B_\nu=P_\nu C_\nu=6.696004159297337\) is retained only as the paper-facing amplitude parameterization, and absolute attachment fixes \(\lambda_\nu=(m_{\star,\mathrm{eV}}/q_{\mathrm{mean}}^{p_\nu})P_\nu C_\nu=1.7237014208357415\), which emits the absolute family and splittings on that branch. The exact positive-segment adapter, the bridge corridor, and the reduced-correction audit remain diagnostic-only beneath this bridge-rigidity surface. The \(\varepsilon = 1/6\) texture language is the extra uniform-center-label collapse of this richer dictionary and is not a recovered-core theorem.
\item \textbf{Hadron masses and resonances}: the particle branch supplies the seeded-family, correlator, and stable-channel readout bridge, but promotable hadron rows remain execution-contract-frozen. The current particle roadmap does not include hadron-mass closure, because the remaining step is one production backend export bundle followed by real nonperturbative production computation and production systematics rather than an additional symbolic theorem alone, and that execution burden is not a realistic target within the current project compute scope.
\item \textbf{Neutrino data}: beyond any legacy, input-dependent capacity-level side estimate tied to $N_{\mathrm{scr}}$, the recovered core does not derive neutrino masses or mixings. A continuation branch emits intrinsic masses, splittings, and PMNS data; its weighted-cycle branch reaches the physical PMNS/hierarchy regime, the normalized same-label overlap-defect weight section closes beneath that branch, and the weighted-cycle bridge-rigidity theorem emits \(C_\nu=G^2Q S^{-1/2}=0.9994295999075177\) above the emitted proxy \(P_\nu=6.699825740519345\). Absolute attachment then fixes the paper-facing amplitude parameterization \(B_\nu=P_\nu C_\nu=6.696004159297337\), \(\lambda_\nu=(m_{\star,\mathrm{eV}}/q_{\mathrm{mean}}^{p_\nu})P_\nu C_\nu=1.7237014208357415\), and therefore the absolute family \(m_i=\lambda_\nu\hat m_i\), \(\Delta m_{ij}^2=\lambda_\nu^2\widehat{\Delta m_{ij}^2}\). Because the positive selector segment is explicit, a compare-only two-parameter continuation adapter still hits the representative central splittings exactly on the same family, but it is strictly diagnostic-only beneath the bridge-rigidity invariant \(C_\nu\) and its induced paper-facing amplitude \(B_\nu\). Together with the exact \(W/Z\) sidecar pair, the exact Higgs/top sidecar, and the exact charged/quark same-family witnesses, the exact non-hadron output lane is complete even though the theorem-grade charged and quark lanes remain open. The attached weighted-cycle stack factors through \(q_e = q_{\mathrm{mean}} qbar_e\), and the best closed constructive local object beneath that bridge-rigidity surface is the defect-weighted same-label edge family \(q_e=\sqrt{g_e d_e}\) together with its induced \(\mu_e\) family; that family is constructive beneath the bridge-rigid theorem surface and does not replace it.
\item \textbf{Dark-sector phenomenology}: the sign, closure, and galaxy-scale response law of any modular-anomaly stress tensor are not derived by the recovered gravity chain.
\item \textbf{Baryogenesis}: suppression-counting estimates are not a substitute for a derived out-of-equilibrium mechanism.
\item \textbf{Black-hole spectroscopy and evaporation closure}: beyond the structural statements retained in Section 5.11 (edge-center obstruction to naive factorization, recoverability-style interior encoding at small CMI, Hawking/KMS normalization, discrete area spectrum, and the Schwarzschild transition identity \(\Delta E = k_B T_H \ln(d'/d)\) for actual sector changes), the manuscript does not derive integer-transition combs, \(k=3\)/QNM selection, Page-type linewidth estimates, PBH burst templates, Kerr/LIGO horizon spectroscopy combs, or any Page-curve/island closure.
\item \textbf{String/worldsheet reorganization}: the Gross-Taylor-type interpretation~\cite{grossTaylor93} requires a controlled large-\(N_{\mathrm{edge}}\) regime distinct from the proved compact-gauge chain, and any further lift to critical superstring structure remains a conjectural extension requiring extra worldsheet ingredients.
\item \textbf{Other downstream continuations}: strong-CP proposals, proton-spin fractions, and similar late-stage continuations appearing on other OPH surfaces are not derived by the recovered core.
\end{itemize}

\textbf{Core/Implementation Boundary.}
\begin{itemize}
\item \textbf{Scaling-limit control}: the Lorentz/null-modular/Einstein chain depends on the explicit scaling-limit premises declared earlier in the manuscript; proving or constructing that regime is an open task.
\item \textbf{Fermionic gauge branch}: the manuscript proves compact gauge reconstruction only in the bosonic internal-gauge branch; a full fermionic/super-Tannakian extension remains open.
\item \textbf{External inputs}: $P$ and $N_{\mathrm{scr}}$ are not derived in this manuscript. The numerical value of $\Lambda$ is not a local recovered prediction because local null data determine the Einstein equation only up to a metric term $\phi g_{ab}$; fixing $\Lambda$ requires the global capacity input $N_{\mathrm{scr}}$.
\item \textbf{Calibration sector}: the D10 numerical implementation depends on printed running, matching, and threshold conventions together with the integrated heat-kernel calibration appendices. Its internal transmutation data are fixed forward from the OPH input \(P\) through the source-only pixel-closure solve, not read back from measured low-energy couplings. The displayed carrier closes its own exact \(W/Z\) chart, and the promoted target-free source-only repair theorem emits the public electroweak quintet on the Phase II calibration branch. The freeze-once coherent repair surface is retained only as compare-only validation; on that frozen authoritative surface the canonical \(W/Z\) pair is hit exactly, but it remains below the target-free theorem.
\item \textbf{Hadron execution boundary}: the hadron lane is not blocked by one small symbolic gap. The next object is one production backend export bundle with publication-complete manifest provenance and real stable-channel correlator arrays; the normalized production dump and then production statistical/systematic control only follow after that bundle lands. Because that remaining nonperturbative execution burden is not a realistic target within the current project compute scope, hadron rows remain outside the current particle roadmap and outside the derived output surface.
\end{itemize}

The compact Higgs/top forward seed downstream of the D10 gauge core is closed on its emitted one-scalar branch and supplies the public D11 rows. It remains a secondary quantitative descendant rather than a recovered-core theorem or a completed particle-spectrum closure. A separate compare-only inverse slice on the same D11 Jacobian can be solved to hit the canonical Higgs/top reference pair exactly, but that sidecar is not the live forward branch and does not change the theorem status of the public D11 rows.

Failures localize by tier. A failure in a Phase-III flavor, dark-sector, baryogenesis, spectroscopy, string, strong-CP, proton-spin, or similar continuation retracts that continuation only. A failure in D10 challenges the corresponding calibration sector. Only a failure in Phase I would overturn the recovered-core claim set.

\subsubsection{Comparison with other unification approaches}\label{85-comparison-with-other-unification-approaches}

Unified models attempting to tie together QFT, gravity, and SM structure tend to encounter a repeatable set of conceptual difficulties. This subsection examines how the observer-patch holography framework addresses these common pitfalls.

\textbf{1. Subsystem factorization in gauge theory and gravity.}

In gauge theories and gravity, the Hilbert space does not cleanly split as "inside \ensuremath{\otimes} outside" across a cut. This infects entanglement entropy definitions, area terms, edge modes, and observable identification. Many unification attempts handwave this or patch it with conventions.

\emph{Resolution:} The framework builds from a net of von Neumann algebras on patches plus overlap consistency. It does not start from na\"{i}ve tensor factorization. The gauge-as-gluing + regulator package yields edge-center completion: a canonical block decomposition on collars where the center captures superselection data at the cut, and the state becomes (exactly or approximately) Markov across the collar. The entropy split S(\ensuremath{\rho}\_C) = S\_bulk + \ensuremath{\langle}L\_C\ensuremath{\rangle} follows from having a center with sector labels. This replaces the ad hoc "add an area term" move.

\textbf{2. Modular Hamiltonian nonlocality.}

Many entanglement-based gravity derivations depend on modular Hamiltonians that look like local stress-tensor charges (true only in special states/regions). In generic QFT states, modular Hamiltonians are nonlocal, making "first law of entanglement \ensuremath{\Rightarrow} Einstein equation" arguments fragile.

\emph{Resolution:} The Markov collar condition does heavy lifting: approximate Markov implies approximate modular additivity, with the defect controlled by conditional mutual information. This makes "modular locality" a controlled approximation. It is not treated as an assumption. On the extracted prime geometric subnet of Theorem 4.2, the controlled tangent-half-space comparison then locks modular flow to geometric dilations with rigid \(2\pi\) normalization.

\textbf{3. Lorentz invariance as a derived output.}

Discrete microscopic models generally break Lorentz symmetry, and many unified proposals simply postulate Lorentz invariance in the IR.

\emph{Resolution:} Lorentz kinematics are tied to geometric modular flow on caps. On the extracted prime geometric subnet, once modular flow acts as conformal transformations on S\textsuperscript{2}, we get Conf\textsuperscript{+}(S\textsuperscript{2}) \ensuremath{\cong} PSL(2,\ensuremath{\mathbb{C}}) \ensuremath{\cong} SO\textsuperscript{+}(3,1), so the Lorentz group appears as a theorem-level output of modular structure on that branch. No external spacetime symmetry axiom is added.

\textbf{4. Dynamics beyond kinematics.}

Many approaches produce emergent geometry/kinematics but stall at dynamics: why Einstein\textquotesingle s equations (with the right coefficient) rather than some other geometric PDE?

\emph{Resolution:} The framework combines MaxEnt entanglement equilibrium, the derived K\_C = 2\ensuremath{\pi}B\_C structure with rigid normalization, and an EFT bridge identifying modular energy with stress-tensor charges. The null modular additivity route reduces much of that bridge to Markov/edge-center mechanisms on null strips plus explicit scaling-limit regularity premises. It does not rely only on "assume a UV CFT."

\textbf{5. Gauge symmetry origin and compactness.}

Most unification stories pick a gauge group and work out consequences. Emergent-gauge approaches sometimes produce noncompact groups or uncontrolled redundancies.

\emph{Resolution:} Gauge symmetry is recast as redundancy in overlap identifications (gauge-as-gluing). From edge sectors and fusion, a tensor category is reconstructed; Tannaka-Krein / Doplicher-Roberts reconstruction then yields a compact group G given the categorical hypotheses. "Gauge symmetry" names the gluing redundancy at the conceptual level. "Compact group" is the mathematical form compatible with finite-dimensional sector/fiber-functor structure.

\textbf{6. Massless photon and graviton usually hand-imposed.}

Getting massless gauge bosons is easy if exact gauge invariance is assumed, but that restates the problem. Massless graviton is more delicate (mass terms, vDVZ discontinuity, strong coupling scales).

\emph{Resolution:} Once gauge and diffeomorphism invariance are emergent redundancies of description (from gluing consistency / emergent geometry), hard mass terms are forbidden: "a coordinate system\textquotesingle s Jacobian can\textquotesingle t show up as a physical mass." These symmetry-protected zeros emerge from the same consistency machinery that gives the symmetries.

\textbf{7. Global consistency, anomalies, and loop patching.}

Building physics from local patches hits loop/holonomy problems: consistent gluing on a tree but obstructions around loops. These obstructions are often anomalies or global topological constraints.

\emph{Resolution:} This is elevated to a first-class organizing principle: gluing data on overlaps define cocycles; central defects define a \v{C}ech obstruction class {[}z{]} (and more generally a 2-group/crossed-module cocycle for noncentral defects). "Global consistency exists iff the obstruction class vanishes" becomes the universal statement. Anomalies become "failure to glue" rather than a mysterious quantum pathology.

\textbf{8. Charge quantization without a GUT.}

Without embedding into a simple GUT group, explaining charge quantization (why all isolated color singlets are integer charged) is awkward. Standard lore requires grand unification or monopoles.

\emph{Resolution:} The framework leans on global group structure (the Z\textsubscript{6} quotient) and derives congruence/selection rules for allowed representations/hypercharges. This gives a structural explanation for integer-charged color singlets without introducing the proton-decay channel of simple-GUT models.

\textbf{9. Coupling unification usually forces proton decay.}

Traditional simple-group unification introduces leptoquark gauge bosons (X, Y) mediating proton decay. Experiment keeps pushing limits up, pressuring minimal GUTs.

\emph{Resolution:} "Unification" here is geometric/entropic (shared edge diffusion parameter, heat-kernel weights). It is not an embedding into a simple Lie group. If the reconstructed gauge group genuinely factorizes as a product (sector factorization selector), there are no mixed generators playing the X/Y role. "Unify couplings" and "unify groups" therefore come apart.

\textbf{10. Cosmological constant locality.}

The cosmological constant problem is a graveyard of unified theories: local QFT estimates are enormous, and tiny observed \ensuremath{\Lambda} seems to demand absurd fine tuning.

\emph{Resolution:} From null modular data, T\_ab is reconstructed only up to \ensuremath{\phi}g\_ab. Local consistency conditions and null focusing are blind to vacuum-energy shifts, so the Einstein equation is fixed only up to \ensuremath{\Lambda}g\_ab. \ensuremath{\Lambda} becomes a global "capacity" parameter of the static patch (tied to log dim H\_tot). It is not a locally computable quantity. This resolves the conceptual tension: local microphysics \emph{cannot} fix \ensuremath{\Lambda} by structural information-theoretic reasons.

\textbf{11. UV infinities and nonrenormalizability.}

Unified programs struggle to give sharp, finite microscopic definitions. Formal continuum structures, infinite entropies, and regularization dependence abound.

\emph{Resolution:} The regulator premises explicitly require local patch algebras to be type-I and finite-dimensional, with a MaxEnt branch whose generator is quasi-local and obeys a Lieb-Robinson bound. So the fixed-cutoff UV branch is interacting in the ordinary finite-range sense, and the fundamental degrees of freedom are finite and live on the screen. The framework does \emph{not} claim a unique microscopic UV completion: physical uniqueness is only modulo gauge or implementation hiding together with inert ancillary stabilization. The genuinely noncentral topological branch is also closed at fixed cutoff by the higher-gauge crossed-module collar theorem; the remaining burdens are the continuum BW lift and the refinement-limit transportable gauge branch.

\textbf{12. Predictivity vs. parameter explosion.}

Unified models often explode in parameters, sectors, or vacua, becoming hard to test directly because everything depends on choices.

\emph{Resolution:} The framework compresses freedom into a "pixel area" (resolution) parameter and a total Hilbert space capacity (size) parameter, then derives structure from consistency (Lorentz, Einstein form, compact gauge group reconstruction, exact zeros, quantization patterns). The explicit MAR selector then picks the SM factors and sector factorization on the admissible class discussed in Section 6.2.

\textbf{Structural pattern.} The framework treats locality, Lorentz invariance, gauge symmetry, and
gravity as consequences of consistency conditions among overlapping descriptions together with
information-theoretic properties of states. Modular rigidity then supplies the familiar symmetry
and dynamical structures.

\textbf{Engineering deliverables.} Certain problems can be stated as explicit closure tasks:

\begin{itemize}
\tightlist
\item
  \(\Lambda\) is structurally explained as a global capacity parameter and is not numerically predicted
\item
  A full microphysical derivation of geometric modular action is still required
\end{itemize}

These are shared challenges across unification approaches. The framework provides an explicit map of where they live and what would resolve them.

\begin{center}\rule{0.5\linewidth}{0.5pt}\end{center}

\ifdefined\OPHSkipEmbeddedReferences
\else
\subsection{References}\label{references-paper}

\subsubsection{Foundational results used in this work}\label{foundational-results-used-in-this-work}

\textbf{Modular theory and spacetime:}

\begin{itemize}
\tightlist
\item
  Bisognano, J. J. and Wichmann, E. H. (1975). "On the duality condition for a Hermitian scalar field." \emph{J. Math. Phys.} 16, 985-1007.
\item
  Bisognano, J. J. and Wichmann, E. H. (1976). "On the duality condition for quantum fields." \emph{J. Math. Phys.} 17, 303-321.
\item
  Unruh, W. G. (1976). "Notes on black-hole evaporation." \emph{Phys. Rev. D} 14, 870-892.
\item
  Brunetti, R., Guido, D., and Longo, R. (1993). "Modular Structure and Duality in Conformal Quantum Field Theory." \emph{Commun. Math. Phys.} 156, 201-219. arXiv:funct-an/9302008.
\item
  Wiesbrock, H.-W. (1993). "Half-Sided modular inclusions of von-Neumann-Algebras." \emph{Commun. Math. Phys.} 157, 83-92.
\end{itemize}

\textbf{Gravity from thermodynamics/entanglement:}

\begin{itemize}
\tightlist
\item
  Jacobson, T. (1995). "Thermodynamics of spacetime: The Einstein equation of state." \emph{Phys. Rev. Lett.} 75, 1260-1263. arXiv:gr-qc/9504004.
\item
  Jacobson, T. (2016). "Entanglement equilibrium and the Einstein equation." \emph{Phys. Rev. Lett.} 116, 201101. arXiv:1505.04753.
\item
  Bousso, R., Fisher, Z., Koeller, J., Leichenauer, S., and Wall, A. C. (2016). "Proof of the quantum null energy condition." \emph{Phys. Rev. D} 93, 024017. arXiv:1509.02542.
\end{itemize}

\textbf{Strong subadditivity:}

\begin{itemize}
\tightlist
\item
  Lieb, E. H. and Ruskai, M. B. (1973). "Proof of the strong subadditivity of quantum-mechanical entropy." \emph{J. Math. Phys.} 14, 1938-1941.
\item
  Lieb, E. H. and Robinson, D. W. (1972). "The finite group velocity of quantum spin systems." \emph{Commun. Math. Phys.} 28, 251-257.
\end{itemize}

\textbf{Regulator and edge examples:}

\begin{itemize}
\tightlist
\item
  Chandrasekharan, S. and Wiese, U.-J. (1997). "Quantum link models: A discrete approach to gauge theories." \emph{Nucl. Phys. B} 492, 455-471. arXiv:hep-lat/9609042.
\item
  Donnelly, W. and Wall, A. C. (2015). "Entanglement entropy of electromagnetic edge modes." \emph{Phys. Rev. Lett.} 114, 111603. arXiv:1412.1895.
\item
  Pastawski, F., Yoshida, B., Harlow, D., and Preskill, J. (2015). "Holographic quantum error-correcting codes: toy models for the bulk/boundary correspondence." \emph{JHEP} 06 (2015) 149. arXiv:1503.06237.
\item
  Levin, M. A. and Wen, X.-G. (2005). "String-net condensation: A physical mechanism for topological phases." \emph{Phys. Rev. B} 71, 045110. arXiv:cond-mat/0404617.
\end{itemize}

\textbf{Quantum recovery and Markov chains:}

\begin{itemize}
\tightlist
\item
  Petz, D. (1986). "Sufficient subalgebras and the relative entropy of states of a von Neumann algebra." \emph{Commun. Math. Phys.} 105, 123-131.
\item
  Petz, D. (1988). "Sufficiency of channels over von Neumann algebras." \emph{Quart. J. Math.} 39, 97-108.
\item
  Fawzi, O. and Renner, R. (2015). "Quantum conditional mutual information and approximate Markov chains." \emph{Commun. Math. Phys.} 340, 575-611. arXiv:1410.0664.
\item
  Hayden, P., Jozsa, R., Petz, D., and Winter, A. (2004). "Structure of states which satisfy strong subadditivity of quantum entropy with equality." \emph{Commun. Math. Phys.} 246, 359-374.
\end{itemize}

\textbf{Superselection sectors and gauge reconstruction:}

\begin{itemize}
\tightlist
\item
  Doplicher, S. and Roberts, J. E. (1989). "A new duality theory for compact groups." \emph{Invent. Math.} 98, 157-218.
\item
  Doplicher, S. and Roberts, J. E. (1990). "Why there is a field algebra with a compact gauge group describing the superselection structure in particle physics." \emph{Commun. Math. Phys.} 131, 51-107.
\end{itemize}

\textbf{Tannaka-Krein duality:}

\begin{itemize}
\tightlist
\item
  Tannaka, T. (1938). "\"{U}ber den Dualit\"{a}tssatz der nichtkommutativen topologischen Gruppen." \emph{Tohoku Math. J.} 45, 1-12. (Some sources cite 1939.)
\item
  Krein, M. G. (1949). "A principle of duality for a bicompact group and a square block algebra." \emph{Dokl. Akad. Nauk SSSR} 69, 725-728.
\end{itemize}

\subsubsection{Standard Model and unification (borrowed results)}\label{standard-model-and-unification-borrowed-results}

\textbf{Grand Unified Theories:}

\begin{itemize}
\tightlist
\item
  Georgi, H. and Glashow, S. L. (1974). "Unity of all elementary-particle forces." \emph{Phys. Rev. Lett.} 32, 438-441.
\end{itemize}

\textbf{GIM mechanism:}

\begin{itemize}
\tightlist
\item
  Glashow, S. L., Iliopoulos, J., and Maiani, L. (1970). "Weak interactions with lepton-hadron symmetry." \emph{Phys. Rev. D} 2, 1285-1292.
\end{itemize}

\textbf{Witten anomaly:}

\begin{itemize}
\tightlist
\item
  Witten, E. (1982). "An SU(2) anomaly." \emph{Phys. Lett. B} 117, 324-328.
\end{itemize}

\textbf{MSSM gauge unification:}

\begin{itemize}
\tightlist
\item
  Dimopoulos, S., Raby, S., and Wilczek, F. (1981). "Supersymmetry and the scale of unification." \emph{Phys. Rev. D} 24, 1681-1683.
\item
  Amaldi, U., de Boer, W., and F\"{u}rstenau, H. (1991). "Comparison of grand unified theories with electroweak and strong coupling constants measured at LEP." \emph{Phys. Lett. B} 260, 447-455.
\end{itemize}

\textbf{String/worldsheet continuation:}

\begin{itemize}
\tightlist
\item
  Gross, D. J. and Taylor, W. (1993). "Two-dimensional QCD is a string theory." \emph{Nucl. Phys. B} 400, 181-208. arXiv:hep-th/9301068.
\end{itemize}

\textbf{Experimental inputs:}

\begin{itemize}
\tightlist
\item
  Laghi, D., Carullo, G., Veitch, J., and Del Pozzo, W. (2021). "Quantum black hole spectroscopy: probing the quantum nature of the black hole area using LIGO-Virgo ringdown detections." \emph{Class. Quantum Grav.} 38, 095005. arXiv:2011.03816.
\item
  Particle Data Group (2024). "Review of Particle Physics." \emph{Phys. Rev. D} 110, 030001. \url{https://pdg.lbl.gov/2024/download/db2024.pdf}
\item
  Particle Data Group (2024). "Electroweak Model and Constraints on New Physics." \url{https://pdg.lbl.gov/2024/reviews/rpp2024-rev-standard-model.pdf}
\end{itemize}
\fi

\input{tex_fragments/OBSERVERS_SYNTHESIS_SECTIONS.tex}

\appendix
\section{Candidate Microphysics Reference Architecture}

The synthesis paper keeps the observer-facing theorem package in the main text, but the concrete
finite screen architecture belongs in an appendix. The current theorem-bearing fixed-cutoff
source surface for the edge heat-kernel / Casimir leaf is the microphysics side paper; the
appendix records only the synthesis-level restatement imported from that surface. The reference
architecture used here is a finite cellulation of the horizon screen with boundary-fixed patch
algebras, boundary-fixed overlap algebras on collars, declared shared readout packets, local
repair instruments, and record registers exposed on the observer-facing interface.

\paragraph{Theorem (Regulated patch-net embedding).}
On a fixed finite cellulation and finite patch cover, the regulated OPH microphysics realizes an
explicit finite patch-net layer: patch vertices are genuine screen patches, overlap edges are
genuine collars, packet fields are actual readouts of declared overlap observables, and allowed
repair branches are concrete local channels on the same finite regulator.

\paragraph{Proof sketch.}
All patch and overlap algebras are finite-dimensional at fixed cutoff, and the declared packet
fields are readouts of explicit overlap observables in those algebras. The candidate repair menu
is implemented by local finite-support channels on the same neighborhoods. Therefore the abstract
finite patch-net syntax is embedded into an explicit regulator object rather than introduced as a
separate formal layer.

\paragraph{Theorem (Fixed-cutoff edge heat-kernel branch).}
On the thermalized overlap branch of the same regulated architecture, the stationary sector law is
\[
\pi_\beta(\alpha)\propto d_\alpha e^{-\beta C_2(\alpha)},
\]
and along refinement ladders converging to a compact-group Peter--Weyl decomposition~\cite{peterWeyl27} the same
weights converge cylinderwise to the heat-kernel/Casimir law used in the OPH edge-sector branch.

\paragraph{Proof sketch.}
The regulated proposal kernel is chosen so that the weighted proposal symmetry
\(d_\alpha q_{\alpha\beta}=d_\beta q_{\beta\alpha}\) holds. Detailed balance with the Casimir
acceptance rule then fixes the stationary distribution in finite dimension, and the compact-group
heat-kernel law is the corresponding refinement lift.

\paragraph{Extended derivation.}
The full fixed-cutoff architecture, validation package, and theorem stacks are developed in
\emph{Screen Microphysics, Patches, and Observer Synchronization in OPH}~\cite{ophmicrophysics}.

\section{Interpretive Epilogue: Habitat and Strange-Loop Closure}

This appendix is interpretive. It uses the OPH state-and-law habitat theorem and leaves the
closure map on an observer-supporting sector open.

The OPH inputs used here are:
\begin{enumerate}
\item the patch net \(P\mapsto A(P)\) of von Neumann algebras together with isotony and overlap
restriction maps;
\item the global state on the inductive-limit algebra, whose restrictions witness a nonempty
overlap-consistent local sector;
\item the finite Axiom-3 MaxEnt/refinement branch, which supplies a common finite family of
gauge-invariant local constraint observables across cutoffs;
\item on the fixed-cutoff branch, the derived finite type-I presentations and compact boundary
gluing groups used elsewhere in the OPH package.
\end{enumerate}

The law slot is encoded by the finite expectation-value coordinates of the retained Axiom-3
constraint family. That choice furnishes a compact convex state-and-law habitat within OPH.

\paragraph{Theorem (OPH internal state-and-law habitat theorem).}
Fix finitely many screen patches \(P_1,\dots,P_N\) in the standard OPH setup. For each \(i\), let
\[
M_i:=A(P_i),
\qquad
M_{ij}:=A(P_i\cap P_j),
\]
where \(M_i\) is the patch von Neumann algebra and \(M_{ij}\subset M_i,M_j\) is the overlap
algebra. Let
\[
S_i:=S(M_i)\subset M_i^*
\]
be the full state space of \(M_i\), endowed with the weak* topology \(\sigma(M_i^*,M_i)\).

Choose a finite family of self-adjoint gauge-invariant local constraint observables
\[
\mathcal O=\{O_1,\dots,O_m\}
\]
from the Axiom-3 MaxEnt branch, with each \(O_a\) supported in one of the chosen patches; write
\(i(a)\) for an index such that \(O_a\in M_{i(a)}\).

Define the overlap-consistent state sector
\[
X_{\mathrm{ov}}
:=
\left\{
(\omega_1,\dots,\omega_N)\in \prod_{i=1}^N S_i:
\omega_i|_{M_{ij}}=\omega_j|_{M_{ij}}\ \text{for all }i,j
\right\}.
\]
Define the OPH law-coordinate map
\[
c:X_{\mathrm{ov}}\to \mathbb R^m,
\qquad
c(\omega_1,\dots,\omega_N)
:=
\bigl(\omega_{i(1)}(O_1),\dots,\omega_{i(m)}(O_m)\bigr).
\]
Let
\[
X_{\mathcal O}
:=
\left\{
(\omega_1,\dots,\omega_N,\ell)\in
\left(\prod_{i=1}^N S_i\right)\times \mathbb R^m:
(\omega_1,\dots,\omega_N)\in X_{\mathrm{ov}},
\ \ell=c(\omega_1,\dots,\omega_N)
\right\}.
\]
Then:
\begin{enumerate}
\item the ambient product
\[
E:=
\left(\prod_{i=1}^N M_i^*\right)\times \mathbb R^m
\]
is a Banach space for any product norm;
\item \(X_{\mathcal O}\) is nonempty and convex;
\item \(X_{\mathcal O}\) is norm-closed in \(E\);
\item \(X_{\mathcal O}\) is compact in the product topology
\[
\tau
:=
\left(\prod_{i=1}^N \sigma(M_i^*,M_i)\right)\times
\text{Euclidean topology on }\mathbb R^m;
\]
\item therefore every \(\tau\)-continuous self-map
\[
\Phi:X_{\mathcal O}\to X_{\mathcal O}
\]
has a fixed point.
\end{enumerate}

\begin{proof}
Each \(M_i\) is a von Neumann algebra in the OPH patch net, hence a Banach space, so each dual
\(M_i^*\) is Banach. A finite product of Banach spaces is Banach, hence so is \(E\).

For each \(i\), the state space
\[
S_i=\{\omega\in M_i^*:\omega\ge 0,\ \omega(1)=1\}
\]
is convex and weak* compact by Banach--Alaoglu, because it is a weak* closed subset of the dual
unit ball. For every pair \((i,j)\), restriction along \(M_{ij}\subset M_i\) and
\(M_{ij}\subset M_j\) defines affine weak* continuous maps
\[
r_{ij}:S_i\to S(M_{ij}),
\qquad
r_{ji}:S_j\to S(M_{ij}).
\]
Hence \(X_{\mathrm{ov}}\) is an intersection of affine equalizer sets, so it is convex and
\(\tau\)-closed inside \(\prod_i S_i\). Since \(\prod_i S_i\) is compact and convex in the
product weak* topology, \(X_{\mathrm{ov}}\) is compact and convex as well.

Nonemptiness comes from the OPH global state on the inductive-limit algebra: its restrictions
\[
\bigl(\omega|_{M_1},\dots,\omega|_{M_N}\bigr)
\]
belong to \(X_{\mathrm{ov}}\) because the restrictions agree on every overlap by construction.

Each coordinate \(\omega_{i(a)}\mapsto \omega_{i(a)}(O_a)\) is affine and weak* continuous, since
evaluation at a fixed algebra element is weak* continuous. Thus \(c\) is affine and
\(\tau\)-continuous. The graph map
\[
G:X_{\mathrm{ov}}\to E,
\qquad
G(x):=(x,c(x)),
\]
is affine and continuous, so its image \(X_{\mathcal O}=G(X_{\mathrm{ov}})\) is convex and
\(\tau\)-compact.

To see norm-closedness in \(E\), let \((x_n,c(x_n))\to (x,\ell)\) in norm. Norm convergence
implies weak* convergence on each coordinate, so \(x\in X_{\mathrm{ov}}\) because
\(X_{\mathrm{ov}}\) is weak* closed. Since \(c\) is weak* continuous, \(c(x_n)\to c(x)\), while
norm convergence in \(\mathbb R^m\) also gives \(c(x_n)\to \ell\). Hence \(\ell=c(x)\), so
\((x,\ell)\in X_{\mathcal O}\).

Finally, \(X_{\mathcal O}\) is a compact convex subset of the locally convex topological vector
space \(E\) equipped with \(\tau\). Therefore every \(\tau\)-continuous self-map
\(\Phi:X_{\mathcal O}\to X_{\mathcal O}\) has a fixed point by the Schauder--Tychonoff
fixed-point theorem.
\end{proof}

\paragraph{Corollary (Fixed-cutoff OPH Brouwer corollary).}
On the fixed-cutoff realized presentation of OPH, each chosen patch algebra \(M_i=A(P_i)\) is
finite-dimensional. Then \(E\) is finite-dimensional, the weak* and norm topologies coincide on
\(E\), and \(X_{\mathcal O}\) is a nonempty compact convex subset of a finite-dimensional real
vector space. Hence every continuous self-map
\[
\Phi:X_{\mathcal O}\to X_{\mathcal O}
\]
has a fixed point by Brouwer. If, moreover, \(\Phi\) is a contraction for some metric on
\(X_{\mathcal O}\), then the fixed point is unique and the iterates converge to it from every
starting point by Banach's contraction theorem.

\begin{proof}
On that fixed-cutoff realized presentation, each chosen patch algebra is identified
with a finite-dimensional type-I presentation modulo compact boundary redundancy, so each \(M_i\)
and \(M_i^*\) is finite-dimensional. Hence \(E\) is finite-dimensional and all relevant locally
convex topologies coincide. The habitat theorem above shows that \(X_{\mathcal O}\) is nonempty,
compact, and convex, so Brouwer gives a fixed point for every continuous self-map. The
contraction statement is the usual Banach fixed-point theorem.
\end{proof}

The additional inputs needed for a strange-loop closure theorem are:
\begin{enumerate}
\item an actual OPH closure map \(\Phi\) built from internal operations such as overlap repair,
Axiom-3 MaxEnt reprojection, collar recoverability, and record-sector coarse-graining;
\item a proof that the stronger ``observer-supporting'' or ``selected-world'' criteria carve out a
nonempty compact convex \(\Phi\)-invariant subset of \(X_{\mathcal O}\);
\item any uniqueness, contraction, or Lyapunov-type stability estimate for that closure map.
\end{enumerate}

The habitat theorem supplies the Banach/compact-convex setting. A strange-loop closure theorem
would additionally require the closure map and the stronger invariant-sector and stability
inputs.

\subsection{Timeless Causal Chain}

The hypothesis is represented by the chain
\[
\begin{aligned}
\text{information constraints}
&\to \text{effective physics}
\to \text{complex chemistry and biology}\\
&\to \text{observers with records}
\to \text{formal OPH reconstruction}\\
&\to \text{OPH-compatible simulation substrate}
\to \text{information constraints}.
\end{aligned}
\]
In this appendix the chain is interpreted as a consistency loop. The discussion does not use it as
a linear first-cause narrative. The Escher drawing-hands image is used only as a geometric
analogy for this closure structure.

\subsection{Why Reality Exists and Why It Has This Form}

Under this hypothesis, existence is identified with nontrivial fixed points of the OPH
consistency map. The observed form of reality is selected by stability requirements:
\begin{enumerate}
\item overlap-consistent sharable records,
\item recoverability under local information loss,
\item long-lived structured sectors that support observer formation,
\item internal reconstructibility of the governing rule set.
\end{enumerate}
The appearance of intelligent observers enters the closure mechanism itself. Observer-capable
worlds are the sectors that complete the loop.

\subsection{Additional Problem Closures}

Using the documentary node labels defined in Section~\ref{16-summary-complete-axiom-set}, Phase I
is D1--D5 together with D7--D9; D6 is the separate input-dependent corollary; and D12
collects phenomenological continuations. Ref.~\cite{ophmicrophysics} proves fixed-cutoff theorem
packages for the measurement/Born-rule interface and for checkpoint/restoration/backup.
Table~\ref{tab:additional-closures} records the status of the broader problem-closure claims
discussed in the OPH synthesis surfaces. It also records the appendix-tier closure stories so they
cannot float outside the official status ledger.

\begin{longtable}{p{0.24\textwidth} p{0.19\textwidth} p{0.50\textwidth}}
\caption{Problem-closure statements sorted by claim tier.}\label{tab:additional-closures}\\
\toprule
Problem domain & Status & OPH mechanism / scope note \\
\midrule
\endfirsthead
\toprule
Problem domain & Status & OPH mechanism / scope note \\
\midrule
\endhead
\bottomrule
\endfoot
Quantum gravity consistency & Proved corollary of earlier theorems & Conditional Lorentz and Jacobson-type Einstein branches already follow once the stated scaling-limit, null-bridge, and fixed-cap premises are imported. \\
UV completion / microscopic uniqueness & Explicit future program & Fixed-cutoff theorem packages exist, but a unique microscopic representative, the continuum BW lift on the extracted geometric subnet, and the refinement-limit transportable compact-gauge lift remain open. \\
Measurement problem & Explicit future program & Ref.~\cite{ophmicrophysics} closes the fixed-cutoff central-record / Born-L\"uders interface, but a full philosophical or continuum-level measurement closure is not derived here. \\
Cosmological principle and horizon homogeneity & Explicit future program & MaxEnt symmetry inheritance is suggestive only; an independent derivation of cosmological isotropy and homogeneity remains outside the recovered-core theorems. \\
Problem of time & Proved corollary of earlier theorems & On the BW\(_{S^2}\) geometric branch, observer-relative time is read from cap modular flow. This is a corollary-level interpretation of D3 rather than a separate theorem package. \\
Black-hole information paradox & Explicit future program & Edge-center sectorization and recoverability motivate an internal resolution template, but the full black-hole-information branch remains outside Phase I. \\
Cosmological constant problem & Proved corollary of earlier theorems & Once the screen-capacity identification is declared, \(\Lambda\) is fixed as the separate D6 input-dependent corollary rather than by the local recovered core. \\
Magnetic monopole expectation from simple GUT groups & Proved corollary of earlier theorems & The realized product-group branch \(\SU(3)\times\SU(2)\times\U(1)/\mathbb Z_6\) removes the standard simple-GUT symmetry-breaking monopole channel. \\
Proton stability against gauge-mediated decay & Proved corollary of earlier theorems & The realized product-group branch removes the standard simple-GUT gauge-mediated proton-decay channel. \\
Proton spin fraction & Explicit future program & Closing the proton-spin claim still requires nonperturbative QCD completion beyond Phase I. \\
Dark matter phenomenology & Explicit future program & Modular-anomaly response laws remain conjectural beyond the recovered gravity chain. \\
Baryon asymmetry scale & Explicit future program & Suppression-counting estimates are not a substitute for a derived out-of-equilibrium mechanism. \\
Three generations & Proved corollary of earlier theorems & \(N_g=3\) is recovered on the realized MAR-admissible branch. \\
Downstream flavor structure, Koide, and charged-lepton fits & Explicit future program & Later flavor constructions require extra ans\"atze and remain outside the recovered core. \\
String/worldsheet reorganization & Explicit future program & Requires a controlled large-\(N_{\mathrm{edge}}\) regime distinct from the physical color rank \(N_c=3\); any further lift to critical superstring structure remains conjectural. \\
Why anything exists / strange-loop closure & Explicit future program & Appendix B provides only the habitat theorem and fixed-point setting. The OPH closure map, invariant observer-supporting sector, and uniqueness/stability proofs remain open. \\
Observer continuation / backup & Explicit future program & Ref.~\cite{ophmicrophysics} proves fixed-cutoff checkpoint/restoration for observer-accessible event algebras, but stronger continuation and substrate-transfer claims remain speculative. \\
\end{longtable}

\section{Observer Continuation and Backup}

This appendix uses the fixed-cutoff checkpoint/restoration/error/identity theorem package proved
in Ref.~\cite{ophmicrophysics} and summarizes its algebraic interface in observer language. It
does not define a strange-loop closure map on an invariant observer-supporting sector or establish
uniqueness and stability for such a map.

\subsection{Observer as Algebraic Pattern}

Here, an observer is represented by
\[
O=(P,\mathcal{A}(P),\rho,R),
\]
where \(P\) is a screen patch, \(\mathcal{A}(P)\) is its local algebra, \(\rho\) is the local
state, and \(R\) is the record algebra. Records are carried exactly by central record projectors
and, on practical readout surfaces, by approximately commuting projectors in overlap centers, so
they are shareable without violating no-cloning constraints for generic quantum states.

Ref.~\cite{ophmicrophysics} proves a fixed-cutoff observer-facing measurement interface: a
finite central record algebra, Born probabilities for its event projectors, and L\"uders
conditioning on that same commuting record algebra. This appendix uses that fixed-cutoff
measurement package in observer language. If a practical readout uses projectors that are
\(\delta_{\mathrm{rec}}\)-close in operator norm to that central reference algebra and the
accessible state is perturbed by at most \(\varepsilon\) in trace norm, each declared elementary
record-event probability shifts by at most \(\varepsilon+\delta_{\mathrm{rec}}\).

\subsection{Markov Collar Factorization}

For a collar tripartition \(A\)-\(B\)-\(D\) with conditional mutual information
\(I(A:D|B)\le\varepsilon\), the exact Markov limit gives
\[
\rho_{ABD}
=\bigoplus_{\alpha} p_{\alpha}\,
\rho^{(\alpha)}_{A b_L^{\alpha}}
\otimes
\rho^{(\alpha)}_{b_R^{\alpha} D}.
\]
The sector label \(\alpha\) is classical center data. This decomposition is the mathematical
basis for extracting an interior observer state with a controlled boundary interface.

\subsection{Checkpoint and Restoration Map}

A checkpoint is the tuple
\[
\mathcal{C}=\bigl(R,\alpha,\rho_{\mathrm{int}}^{(\alpha)}\bigr),
\]
with \(\rho_{\mathrm{int}}^{(\alpha)}\) the interior conditional state. Given a compatible target
environment state \(\sigma_{\mathrm{env}}^{(\alpha)}\), a restored state is
\[
\rho_{\mathrm{new}}^{(\alpha)}
=
\rho_{\mathrm{int}}^{(\alpha)}
\otimes
\sigma_{\mathrm{env}}^{(\alpha)}.
\]
For approximate Markov collars, recovery is controlled by the standard bound
\[
\left\|
\rho_{ABD}
-
(\mathrm{id}_A\otimes\mathcal{R}_{B\to BD})(\rho_{AB})
\right\|_1
\le
2\sqrt{\ln 2\,\varepsilon}.
\]
This gives quantitative trace-distance control on the recovered state under finite conditional
mutual information (finite collar error). Interpreting that control as observer continuation
requires additional modeling assumptions beyond the bound itself.

The fixed-cutoff microphysics claim is stronger than the bare recoverability estimate written
above. In Ref.~\cite{ophmicrophysics}, exact checkpoint restoration preserves the full future
law of observer-accessible events, while an \(\varepsilon\)-accurate restoration changes that
future law by at most \(\varepsilon\) in total variation. The stronger strange-loop closure story
remains open.

\subsection{Physical Meaning}

At fixed cutoff, Ref.~\cite{ophmicrophysics} proves a checkpoint/restoration theorem and backup
corollary for observer-accessible event algebras. This appendix states the observer-facing meaning
of that theorem and its algebraic prerequisites. Open items are the stronger substrate-selection,
strange-loop closure, uniqueness, and stability package.

\input{tex_fragments/OBSERVERS_TECHNICAL_SUPPLEMENT_PORT.tex}

\input{tex_fragments/OBSERVERS_REFERENCES.tex}
\end{document}
